% DER_XVII-b: Grothendieck's text with the commentary set in, then the overview, inside a group that applies the chunk's own definitions
\dgrpart{1}{XVII-b}{Part 2 of 2, manuscript pages 58--118}{XVII-b\enspace\textperiodcentered\enspace Chapter XVII, part 2 of 2\enspace\textperiodcentered\enspace manuscript pages 58--118}{Chapter XVII \textperiodcentered\ \textup{XVII-b}}
\dgrchunkbegin
% DER_XVII-b: the definitions of the chunk's preamble, applied inside its group
\let\printedpage\relax\newcommand{\printedpage}[1]{\marginpar{\footnotesize p.~#1}}
\let\msp\relax\newcommand{\msp}[1]{\marginpar{\footnotesize [#1]}}
\let\Esp\relax\newcommand{\Esp}{\operatorname{Esp}}
\let\Fl\relax\newcommand{\Fl}{\operatorname{Fl}}
\let\Ob\relax\newcommand{\Ob}{\operatorname{Ob}}
\let\Hom\relax\newcommand{\Hom}{\operatorname{Hom}}
\let\Ho\relax\newcommand{\Ho}{\operatorname{Ho}}
\let\id\relax\newcommand{\id}{\mathrm{id}}
\let\pr\relax\newcommand{\pr}{\mathrm{pr}}
\let\M\relax\newcommand{\M}{\mathcal{M}}
\let\circled\relax\newcommand{\circled}[1]{\tikz[baseline=(c.base)]\node[draw,circle,inner sep=1pt](c){#1};}
\let\defeq\relax\newcommand{\defeq}{\overset{\text{déf}}{=}}
\let\postilde\relax\newcommand{\postilde}{^{\sim}}
\let\Cof\relax\newcommand{\Cof}{\mathrm{Cof}}
\let\fnm\relax\newcommand{\fnm}[1]{\textsuperscript{(}\footnotemark[#1]\textsuperscript{)}}
\let\cl\relax\newcommand{\cl}{^{\sim}}% tilde printed as exponent after a symbol or a parenthesis
\let\encircled\relax\newcommand{\encircled}[1]{\tikz[baseline=(X.base)]\node[draw,circle,inner sep=1pt](X){#1};}
\let\encoval\relax\newcommand{\encoval}[1]{\tikz[baseline=(X.base)]\node[draw,rounded corners=4pt,inner sep=2pt](X){#1};}
\let\Dl\relax\newcommand{\Dl}{\boldsymbol{\Delta}}
\let\tl\relax\newcommand{\tl}{^{\sim}}
\let\cercle\relax\newcommand{\cercle}[1]{\tikz[baseline=(c.base)]\node[draw,circle,inner sep=1pt](c){#1};}

% ---- from der-ch17_p031-035.tex
\printedpage{29}
% head: A. Grothendieck : Dérivateurs, Ch. XVII, Édition des Universités de Montpellier II et Paris VII
% foot: Version 9 février 2019
lesquels sont visiblement adjoint l'un de l'autre (puisque induits par un couple de foncteurs % [sic] adjoint -> adjoints, accord avec « lesquels »
adjoints $u_!$, $u^*$), et de plus compatibles aux localiseurs. Donc on peut appliquer le lemme des foncteurs adjoints pour les localisés. Pour résumer :

\textbf{Corollaire 3.5.} \textit{Supposons que $I$, $I'$ soient artiniens, et qu'on soit dans l'un des deux cas suivants :}
\begin{enumerate}
\item[a)] \textit{Les espaces} $\Esp(I)$, $\Esp(I')$ \textit{sont noethériens (p.ex. $I$, $I'$ finis).}
\msp{page 58}
\item[b)] \textit{$C$ et $TC = C\cap W$ sont stables par composition transfinie.}
\end{enumerate}
\textit{Supposons de plus que $u^*$ applique $\M(I'^\circ)_{\mathrm{c}}$ dans $\M(I^\circ)_{\mathrm{c}}$, ce qui est le cas notamment si $u$ est localement cofiltrant, i.e. les ensembles ordonnés localisés $i'\backslash I$ sont localement cofiltrants. Alors le foncteur $u_!^W$ construit dans (3.3.1) est adjoint à gauche de $u_W^*$.}

\textbf{Remarque 3.6.} Quand on a un couple de foncteurs adjoints
\[
\begin{tikzcd}
\mathcal{C} \arrow[r, shift left, "u_!"] & \mathcal{C}' \arrow[l, shift left, "u^*"]
\end{tikzcd}
\ ,
\]
avec les localiseurs $W$, $W'$ de $\mathcal{C}$, $\mathcal{C}'$, $u^*$ compatible aux localiseurs, et si de plus on a une sous-catégorie pleine $\mathcal{C}_0$ de $\mathcal{C}$ telle que
\begin{enumerate}
\item[a)] $u_!|\mathcal{C}_0$ transforme flèches de $W$ en flèches de $W'$, et
\item[b)] l'inclusion $\mathcal{C}_0 \hookrightarrow \mathcal{C}$ induit une \textit{équivalence} $\mathcal{C}_0W_0^{-1} \xrightarrow{\sim} \mathcal{C}W^{-1}$ (où $W_0 = W\cap\Fl(\mathcal{C}_0)$),
\end{enumerate}
on trouve deux foncteurs canoniques
\[
\begin{tikzcd}
\mathcal{C}W^{-1} \arrow[r, shift left, "u_!^W"] & \mathcal{C}'W'^{-1} \arrow[l, shift left, "u_W^*"]
\end{tikzcd}
\ ,
\]
où $u_!^W$ est défini seulement à isomorphisme canonique près.

\dgrblock{dgrb931}{The adjunction must survive the choice of resolutions}
\begin{dcomment}{exposition}{p58:est défini seulement à isomorphisme canonique près.}{The adjunction must survive the choice of resolutions}
Grothendieck has constructed the candidate for a left homotopy Kan extension and now asks for the property that makes this candidate useful: maps out of it must correspond to maps into the restricted diagram. The distinction matters because an equivalence between two localized categories remembers neither a preferred direction of replacement nor the comparison maps with the original functor. Corollary~3.5 has enough compatibility to descend an actual adjunction. Remark~3.6 removes that compatibility and exposes what is lost.

Write $F\dashv G$, let $j:\mathcal C_0\hookrightarrow\mathcal C$ be the chosen full subcategory, and write $\gamma$ and $\delta$ for localization. If $\operatorname{Ho}(j)$ is an equivalence and $Fj$ preserves weak equivalences, the composite
\[
 \overline F=\operatorname{Ho}(Fj)\operatorname{Ho}(j)^{-1}
\]
is certainly a functor. Its construction has not supplied a natural transformation $\overline F\gamma\to\delta F$. Such a comparison is part of the data of a left derived functor, with a universal property; it cannot be reconstructed from the objectwise choice alone.

The smallest example makes the distinction quite concrete. Let $\mathcal C$ be the walking arrow $0\to1$, use all arrows as $W$, only isomorphisms as $W'$, and take $F=G=1_{\mathcal C}$. Inverting the single nonidentity arrow makes $\mathcal C[W^{-1}]$ equivalent to the terminal category. Both subcategories $\{0\}$ and $\{1\}$ therefore give localization equivalences. Choosing $\mathcal C_0=\{1\}$ makes $\overline F$ select~$1$, while $\overline G:\mathcal C\to1$ is the unique functor. An adjunction would require
\[
 \operatorname{Hom}_{\mathcal C}(1,0)
 \simeq \operatorname{Hom}_{1}(*,*) .
\]
The first set is empty and the second is a singleton. Choosing $\{0\}$ instead gives an adjunction, since $0$ is initial. Even then $\overline G$ is not fully faithful: the empty hom-set from~$1$ to~$0$ again maps to a singleton. This also proves the second warning on manuscript p.~68. The two failures concern the formal inferences; they do not disprove a correctly constructed homotopy Kan extension or a correctly justified homotopy-cofinality theorem.

A useful sufficient repair is a functor $Q:\mathcal C\to\mathcal C_0$ and a natural weak equivalence $q:jQ\to1_{\mathcal C}$, with $F$ homotopical on $\mathcal C_0$ and $G$ homotopical on all of its domain. The comparison is then $\delta Fq$, rather than an unspecified isomorphism after localization. This is the left-deformation construction of \citet[Theorems~2.2.8 and~2.2.11]{Riehl2014}. A more intrinsic sufficient condition is that the total derived functors are absolute Kan extensions; their adjunction is the unnumbered theorem of \citet{Maltsiniotis2007Adjunction}.

The direction of $q$ explains the example. The constant choice $Q(X)=0$ admits $q_X:0\to X$; the constant choice $Q(X)=1$ admits no component $1\to0$. Thus the obstruction is visible before any hom-set calculation in the localization. The passage to the categorical mapping cylinder in the next proof is a way of recovering compatible adjunctions without first choosing a functorial resolution of every diagram.
\end{dcomment}


\msp{page 59}
Contrairement à ce que prétend \textsc{Anderson}, il n'est par vrai en général que ces deux % [sic] par -> pas, faute de frappe
foncteurs soient adjoints, sans hypothèses supplémentaires. Voici un contrexemple : on prend $\mathcal{C} = \mathcal{C}'$, $u_! = u^* = \id_{\mathcal{C}}$, $W' = \operatorname{Is}(\mathcal{C})$, $W = \Fl(\mathcal{C})$, donc $u_W^*$ est le foncteur canonique
\[
\underbrace{\Pi\,\mathcal{C}}_{\text{groupoïde fondamental}} \qquad \longleftarrow \mathcal{C}\ .
\]
Supposons que $\mathcal{C}$ soit 1-connexe, i.e. $\Pi\,\mathcal{C}$ soit équivalente à $e$. Soit $e\in\Ob\mathcal{C}$ tel que $\operatorname{End}_{\mathcal{C}}(e)$ soit le groupe unité, et prenons $\mathcal{C}_0 = \{e\}$. Alors l'hypothèse sur $\mathcal{C}_0$ est satisfaite. Mais si $u_!^W$ était adjoint à gauche de $u_W^*$, $e$ serait objet initial de $\mathcal{C}$. Par exemple si $\mathcal{C}$ a un objet final $e$ (ce qui implique que $\mathcal{C}$ est 1-connexe et $e$ rigide), cet objet final devrait être initial si \textsc{Anderson} avait raison avec “it is clear that \ldots”.

\msp{page 60}
Néanmoins :

\dgrblock{dgrb932}{Why the mapping cylinder repairs the argument}
\begin{dcomment}{exposition}{p59}{Why the mapping cylinder repairs the argument}
Grothendieck replaces the given map of indexing posets by two maps for which the adjunction can be controlled separately. The construction called a mapping-cone here is a categorical mapping cylinder. It is also the collage of the bifunctor $\operatorname{Hom}_{I'}(-,u-)$, the gluing construction already encountered in the commentary on Chapter~XII, manuscript pp.~59--61. Here its particular purpose is to make one factor preserve cofibrancy and the other preserve weak equivalences without deriving it.

Keep the convention of Theorem~3.7 explicit: $\mathcal M(I^\circ)=\operatorname{Fun}(I,\mathcal M)$ consists of covariant diagrams. The objects of $C(u)$ are the disjoint union of those of $I'$ and $I$. Within either piece the original arrows remain; a mixed arrow from $j\in I'$ to $i\in I$ is an arrow $j\to u(i)$ in $I'$; there are no mixed arrows in the reverse direction. Composition is induced by composition in $I'$ and by $u$. Thus the inclusion $w:I\hookrightarrow C(u)$ is a cosieve, and the inclusion $s:I'\hookrightarrow C(u)$ is a sieve. The retraction $v$ is the identity on $I'$ and equals $u$ on $I$.

The essential feature of $v$ can be read directly from its hom-sets:
\[
 \operatorname{Hom}_{C(u)}(s j,x)
 \simeq\operatorname{Hom}_{I'}(j,vx).
\]
Thus $s\dashv v$, with $vs=1$. In particular $(j,1_j)$ is initial in $(j\downarrow v)$. For an object $(x,\beta:j\to vx)$ of this comma category, its unique arrow from $(j,1_j)$ is the arrow represented by $\beta$ itself. This proves Lemma~3.8, including uniqueness, for arbitrary categories. An initial object gives a common predecessor for every finite diagram and equalizes each parallel pair; hence the required cofilteredness follows. The result is stronger than a statement that the comma category has contractible nerve.

The other factor is controlled by extension by the initial object. If $j\notin I$, the category $(w\downarrow j)$ is empty, because a cosieve admits no arrow from its objects to an object outside it. If $j\in I$, that category has terminal object $(j,1_j)$. Consequently
\[
 (w_!X)_j=\begin{cases}
 X_j,&j\in I,\\
 \varnothing_{\mathcal M},&j\in I'.
 \end{cases}
\]
An objectwise weak equivalence is sent to the same weak equivalence on $I$ and to an identity on $I'$. Both $w_!$ and $w^*$ are therefore homotopical, and their unit and counit descend. No inference of the kind disproved in Remark~3.6 is needed.

There are two small inheritance checks behind the application of the earlier cofibration theorem. A chain in $C(u)$ can cross from one piece to the other at most once, so artinianity passes from the two pieces to their collage. An increasing chain of open subsets stabilizes once its intersections with the two pieces stabilize, which proves the stated noetherianity assertion. Corollary~3.5 therefore applies to $v$, while the preceding extension argument applies to $w$. Compatibility of the already constructed left Kan extensions with composition identifies $u_!^W$ with $v_!^Ww_!^W$. Composing the two actual adjunctions proves Theorem~3.7.

For a small worked instance, take $I$ to be two isolated objects and $I'=e$. The cylinder has an initial vertex $*$ and two arrows $*\to a$, $*\to b$. Extending the pair $(A,B)$ along $w$ gives $(\varnothing\to A,\varnothing\to B)$; taking its colimit along $v$ gives $A\amalg B$. Its cocones to $Y$ are exactly the pairs of arrows $A\to Y$, $B\to Y$. With cofibrant replacements $QA\to A$ and $QB\to B$, the same construction gives the derived coproduct $QA\amalg QB$. The proof above, rather than the mere existence of these replacements, supplies its adjunction in the homotopy categories.
\end{dcomment}


\textbf{Théorème 3.7.} \textit{Soit $(\M,W,C)$ une catégorie à cofibrations stable par $\varinjlim$ finies, et $u : I \longrightarrow I'$ une application croissante entre ensembles ordonnés artiniens. On suppose qu'on est sous l'une des hypothèses suivantes.}
\begin{enumerate}
\item[a)] $\Esp(I)$, $\Esp(I')$ \textit{sont noethériens, p.ex. $I$, $I'$ finis.}

\printedpage{30}
% head: A. Grothendieck : Dérivateurs, Ch. XVII, Édition des Universités de Montpellier II et Paris VII
% foot: Version 9 février 2019
\item[b)] \textit{$C$ et $TC \overset{\text{déf}}{=} C\cap W$ sont stables par composition transfinie, et $\M$ stable par petites $\varinjlim$ filtrantes.}
\end{enumerate}
\textit{Alors le foncteur} $u_!^W : \Ho_W(I^\circ) \longrightarrow \Ho_W(I'^\circ)$ \textit{construit dans} (3.3.1) \textit{est adjoint à gauche du foncteur $u_W^*$ en sens opposé.}

\textsc{Démonstration.} On utilise le diagramme bien connu
\[
\begin{tikzcd}[row sep=large, column sep=large]
 & I \arrow[d, "\partial_1"'] \arrow[dd, bend left=60, "w"] \\
I \arrow[r, shift left, "\partial_0\ (\text{ouv.})"] \arrow[d, "u"'] \arrow[dr, phantom, "\text{cocart.}"] & I\times\boldsymbol{\Delta}^1 \arrow[l, dotted, shift left, "\pr_1"] \arrow[d] \\
I' \arrow[r, shift left, "\text{imm. ouv.}"] & C(u) \arrow[l, dotted, shift left, "v"]
\end{tikzcd}
\]
pour factoriser $u$ en
\[
u \;=\; vw\ ,
\]
où $w$ est une \textit{immersion fermée,} et $v$ est déduit de $\pr_1 : I\times\boldsymbol{\Delta}^1 \longrightarrow I$ par push-out. C'est une rétraction du mapping-cone C(u) sur l'ouvert $I'$, dont le complémentaire fermé s'identifie à $I$ via $w$.

\msp{page 61}
Le foncteur de recollement entre $I'$, $I$
\[
H \;:\; I'^\circ\times I \;\longrightarrow\; \mathrm{Ens}
\]
est donné par
\[
\underbrace{H(i',i)}_{=\,\Hom_{C(u)}(i',i)} \;=\; \Hom_{I'}(i',u(i))\ . \tag*{$(*)$}
\]
La rétraction $v$ de $C(u)$ dans $I'$ applique $i$ dans $u(i)$, et une flèche $i'\longrightarrow i$ correspondant à $i'\longrightarrow u(i)$ dans $I'$ via $(*)$, est transformé par $v$ en cette même flèche $i'\longrightarrow u(i)$. Calculons, % [sic] transformé -> transformée, accord avec « une flèche »
pour $i'_0$ donné, la catégorie $i'_0\backslash\Bigl(C(u),\ v : C(u)\longrightarrow I'\Bigr)$, des couple d'un objet $\xi$ de $C(u)$ % [sic] des couple -> des couples, accord du pluriel
et d'une flèche $i'_0\longrightarrow v(\xi)$ dans $I'$. On voit que $i'_0\backslash C(u)$ admet l'ouvert $i'_0\backslash I'$, avec le complémentaire fermé $i'_0\backslash I$,
\[
\begin{tikzcd}[column sep=large, row sep=huge]
I' \arrow[r, hook, "\text{ouvert}"'] & C(u) \arrow[d] & I \arrow[l, hook', "\begin{array}{l}\text{fermé}\\ \text{complémentaire}\\ \text{de } I'\end{array}"'] \\
 & I'\ ,
\end{tikzcd}
\]
ces deux catégories étant recollées par le bifoncteur
\[
\begin{aligned}
(i'_0\backslash I')^\circ \times (i'_0\backslash I) &\quad\longrightarrow\quad \mathrm{Ens}\\
\Bigl((\alpha : i'_0\longrightarrow i'), (\beta : i'_0\longrightarrow u(i), i)\Bigr) &\quad\longmapsto\quad \Hom_{i'_0\backslash I'}(i',u(i))
\end{aligned}
\]
\printedpage{31}
% head: A. Grothendieck : Dérivateurs, Ch. XVII, Édition des Universités de Montpellier II et Paris VII
% foot: Version 9 février 2019
induit par $(*)$\,; i.e. les flèches de $(\alpha : i'_0\longrightarrow i') \in \Ob\,(i'_0\backslash I')$ dans $(\beta : i'_0\longrightarrow u(i)) \in i'_0\backslash I$ sont les flèches $i' \xrightarrow{\gamma} u(i)$ telles que le triangle
\[
\begin{tikzcd}
 & i'_0 \arrow[dl, "\alpha"'] \arrow[dr, "\beta"] & \\
i' \arrow[rr, "\gamma"] & & u(i)
\end{tikzcd}
\]
commute.

\msp{page 62}
Si $i' = i'_0$, $\alpha = \id_{i'_0}$, alors cette commutativité revient à $\gamma = \beta$, ce qui signifie qu'il y a une flèche et une seule de $\alpha$ dans $(i,\beta)$. Cela montre donc que $(i'_0,\id_{i'_0})$ est un objet initial de $i'_0\backslash C(u)$. On a donc prouvé le lemme :

\textbf{Lemme 3.8.} \textit{Soit $u : I\longrightarrow I'$ une flèche dans} (Cat)\textit{, factorisons-le en}
\[
I \xrightarrow{\;w\;} C(u) \xrightarrow{\;v\;} I'\ ,
\]
\textit{où $C(u)$ est le mapping-cone défini page} 60\textit{, $w$ une immersion fermée de complémentaire $I'$, $v$ la rétraction canonique de $C(u)$ sur $I'$. Ceci posé, pour tout $i'_0\in\Ob I'$, la catégorie $i'_0\backslash(C(u),v)$ admet un objet initial – donc $v : C(u)\longrightarrow I'$ est (tout ce qu'il y a de) cofiltrante.}

Nous pouvons maintenant prouver 3.7. Tout d'abord, les hypothèses faites sur $I$, $I'$ se transportent à $C(u)$ (hypothèse artinienne, et éventuellement hypothèse noethérienne sur $\Esp(C(u))$). Donc
\msp{page 63}
on peut appliquer les résultats précédents, notamment 3.5, aux flèches $v$, $w$. Le lemme précédent montre bien que $v$ satisfait aux conditions de 3.5, donc le foncteur $v_!^W$ construit dans 3.3.1 est adjoint à gauche de $v_W^*$. Considérons d'autre part le foncteur $w_!^W$ construit dans 3.3.1, je dis que la même conclusion s'applique. Cela provient du fait que $w$ est une immersion fermée. En effet, dans ce cas il est immédiat que
\[
w_! \;:\; \M(I^\circ) \;\longrightarrow\; \M(I'^\circ), \qquad\quad \text{où } I = C(u), % [sic] I = C(u) -> I' = C(u), w a pour but le mapping-cone C(u)
\]
est compatible avec les localiseurs, puisque l'on a
\[
\left\{
\begin{array}{lcll}
w_!(X)(j) & \simeq & X(j) & \text{si } j\in I,\\[1ex]
w_!(X)(j) & \simeq & \emptyset_{\M} & \text{si } j\in J\backslash I,
\end{array}
\right.
\]
par le calcul des fibres
\[
w_!(X)(j) \;\simeq\; \varinjlim_{i\in I/j} X(i)\ .
\]
Donc la définition de $w_!^W$ donnée dans 3.3.1 revient en fait à passer directement aux catégories localisées pour $w_!$, sans avoir à se restreindre à $\M(I^\circ)_{\mathrm{c}}$, et le lemme des foncteurs adjoints nous dit que
\printedpage{32}
% head: A. Grothendieck : Dérivateurs, Ch. XVII, Édition des Universités de Montpellier II et Paris VII
% foot: Version 9 février 2019
\msp{page 64}
ce foncteur est adjoint à gauche de $w_W^*$.

Comme on a évidemment
\[
u_!^W \;\simeq\; v_!^W\circ w_!^W\ , \qquad\quad u_W^* \;\simeq\; w_W^*\circ v_W^*\ ,
\]
on en conclut que $u_!^W$ est bien adjoint à gauche de $u_W^*$, q.e.d.

\textbf{Corollaire 3.9.} \textit{Soit $(\M,W,C)$ une catégorie à cofibrations. Alors le prédérivateur $\mathbf{D}_W$, restreint à la catégorie} $\mathrm{Diag}_0$ \textit{formée des catégories $I$} ordonnées \textit{artiniennes et telles que} $\Esp(I)$ \textit{soit noethérien, satisfait aux axiomes} Dér 1 (décomposition), Dér 2 (localisation), Dér 3 (objet final), Dér 4 bis (existence de $u_!$), Dér 5 bis (mode de calcul des $u_!$), Dér 6 (exactitude des carrés $W$-cocartésiens). \textit{Si $C$, $TC = C\cap W$ sont stables par composition transfinie, on a le même résultat sur la catégorie} Diag \textit{formée des catégories $I$ ordonnées artiniennes sans plus.}

\dgrblock{dgrb933}{The invariant theory and the auxiliary structures}
\begin{dcomment}{context}{p64:on a le même résultat sur la catégorie}{The invariant theory and the auxiliary structures}
The point of Corollary~3.9 is that a one-sided structure on the coefficient category already constructs the corresponding half of the theory of diagrams. Cofibrations organize the calculation of left homotopy Kan extensions. A dual structure of fibrations organizes right homotopy Kan extensions. Grothendieck does not require the two chosen classes to determine one another by lifting properties before recognizing the resulting operations as part of the same homotopy theory.

This viewpoint is stated particularly clearly in the letter to Thomason of 2 April 1991: weak equivalences specify the essential data, while fibrations and cofibrations serve to construct operations and establish their existence \citep[point~1, printed pp.~3--4]{Grothendieck1991Letter}. Brown's categories of fibrant objects provided a one-sided precedent \citep[Section~1]{Brown1973}; the text here works dually and allows noncofibrant objects before passing to the cofibrant subcategory. A Quillen structure imposes the further agreement of the pairs $(C,F\cap W)$ and $(C\cap W,F)$ through lifting and factorization. That agreement is a powerful construction device, but it is a stronger condition than having two separate ways to calculate the operations attached to $W$.

Two conventions need attention when reading the corollary. First, its axiom numbers are Grothendieck's. Dér~2 is the localization or evaluation-conservativity assertion. Dér~3 concerns the full faithfulness of the constant-diagram functor when the indexing category has a final object; it does not assert that $\mathbf D(e)$ has a final object. Dér~4 bis and Dér~5 bis are the existence and pointwise-computation assertions for left Kan extensions. Dér~6 includes the corresponding exactness of base change on the specified homotopy-cocartesian squares. This older list should not be substituted term by term for a modern list numbered Der1--Der4.

Second, the preceding theorem is about covariant diagrams $\mathcal M(I^\circ)$, whereas the usual presheaf notation is $\mathcal M(I)=\operatorname{Fun}(I^{\mathrm{op}},\mathcal M)$. Passing to the presheaf prederivator reverses the admissible indexing orders. The artinian condition of Theorem~3.7 becomes a noetherian condition on $I$; the extra topological condition becomes noetherianity of $\operatorname{Esp}(I^\circ)$. Corollary~3.9 and the corresponding sentence of~3.11 omit this reversal in their wording. Grothendieck gives the corrected presheaf formulation explicitly in Chapter~XIX, Proposition~3.2 and the discussion in~3.4 \citep{Grothendieck1990XIX}. The finite examples are unaffected, but the infinite indexing convention should be fixed before applying the result.

This also explains why the comma symbols change. For a covariant diagram $X:I\to\mathcal M$ the value of $u_!X$ at $j$ is a colimit over $(u\downarrow j)=I/j$. For a presheaf $X:I^{\mathrm{op}}\to\mathcal M$ it is a colimit over $(j\downarrow u)^{\mathrm{op}}$. The printed $i'\backslash I$ in the square on manuscript p.~70 is therefore correct in the presheaf convention of~3.11. The square carries its canonical comma natural transformation; it is not meant as an equality of its two composites in the underlying ordinary category of categories.
\end{dcomment}


\msp{page 65}
\textbf{3.10.} Ainsi, pour avoir un dérivateur au plein sens du terme, soit sur $\mathrm{Diag}_0$, soit sur Diag, il ne manque que Dér 4 (existence des $u_*$), Dér 5 (calcul des $u_*$) et Dér 6 (exactitude des carrés $W$-cartésiens). On va les obtenir en principe à l'aide d'une structure de \textit{catégorie à fibrations} $(\M,F,W)$ associé à $(\M,W)$, i.e. l'espèce de structure duale de celle introduite % [sic] associé -> associée, accord avec « structure »
page 39. Ces deux structures n'ont pas à être reliées entre elles d'une façon autre que par le fait qu'elles correspondent au même $W$. Elles sont considérablement plus faciles à construire que des structures de Quillen $(W,C,F)$, où $C$, $F$ doivent être reliés de façon extrêmement stricte, avec les propriétés $C \longleftrightarrow \underbrace{TF}_{=\,F\cap W}\ ,\ \underbrace{TC}_{=\,C\cap W} \longleftrightarrow F$ (dont il n'est pas question dans l'axiomatique de K. S. \textsc{Brown}), et factorisation pour $(C,TF)$ et $(TC,F)$.

\msp{page 66}
Mais revenant à la construction des $u_!^W$, il reste le problème de les construire pour une flèche $u : I'\longrightarrow I$ \textit{quelconque} dans Cat. % [sic] u : I' -> I -> u : I -> I', u_!^W va de Ho_W(I°) dans Ho_W(I'°)

À vrai dire, la construction du foncteur
\[
u_!^W \;:\; \Ho_W(I^\circ) \;\longrightarrow\; \Ho_W(I'^\circ) \tag*{(3.3.1)}
\]
est valable sans supposer que $I'$ soit une catégorie ordonnée – ça peut être une catégorie quelconque. Cela provient du fait que le foncteur induit par $u_!$,
\[
u_{!\,\mathrm{cof}} \;:\; \M(I^\circ)_{\mathrm{c}} \;\longrightarrow\; \M(I'^\circ)
\]
est encore compatible avec les localiseurs, comme on voit par le calcul des fibres
\[
u_!(X)(i) \;\simeq\; \varinjlim_{i\in I/i'} X(i)\ . % [sic] u_!(X)(i) -> u_!(X)(i'), la fibre est prise en i'
\]
Ainsi, on est ramené aux cas où $I' = e$, quitte à remplacer $I$ par $I/i'$, et en notant que % [sic] aux cas -> au cas, singulier
$(I/i'\longrightarrow I$ étant étale) l'image inverse de $X$ sur $I/i'$ est cofibrant (cf. 3.2.a). % [sic] cofibrant -> cofibrante, accord avec « l'image inverse »

\dgrblock{dgrb934}{A finite source can require an infinite supply of coefficients}
\begin{dcomment}{development}{p66:Ainsi, on est ramené aux cas}{A finite source can require an infinite supply of coefficients}
Grothendieck observes that the target of a Kan extension need not itself be ordered: the calculation can be made over comma categories. This is correct in the transfinite setting, but it does not extend the finite-colimit alternative without a further hypothesis. The new comma categories may be much larger than the source. We make the obstruction precise in a case where weak equivalences introduce no additional difficulty.

Take $\mathcal M=\mathrm{FinSet}$, let $W$ be the isomorphisms, and let every map be a cofibration. All the finite cofibration axioms hold: finite colimits exist, pushouts preserve isomorphisms, and $f=1\circ f$ is a cofibration--weak-equivalence factorization. Every object is cofibrant. For the functor $e\to BG$, covariant restriction is the forgetful functor
\[
 U:\mathrm{FinSet}^{G}\longrightarrow\mathrm{FinSet}.
\]
The comma category $(e\downarrow *)$ is the discrete set $G$, so the proposed pointwise formula on a singleton requires $\coprod_G1$. If $G$ is infinite this colimit is absent from $\mathrm{FinSet}$. There is also a sharper test for whether any left adjoint exists.

\textbf{Finite-action criterion.}\ours\par
For a group $G$, set
\[
 N=\bigcap_{[G:H]<\infty}H.
\]
The forgetful functor $U$ has a left adjoint if and only if $[G:N]$ is finite. When it exists, the left adjoint sends a finite set $S$ to $(G/N)\times S$, with $G$ acting on the first factor.

\emph{Proof.}
Suppose $F\dashv U$, put $P=F(1)$, and let $p\in P$ be the point specified by the unit. Equivariance and the adjunction give a natural bijection
\[
 \operatorname{Hom}_G(P,A)\longrightarrow U(A),\qquad
 f\longmapsto f(p).
\]
In particular $P$ is nonempty. For the trivial two-element $G$-set, an equivariant map is precisely a function on the orbit set of $P$. If $P$ has $k$ orbits, the left side has $2^k$ elements and the right side has two. Thus $k=1$ and $P$ is transitive. Write $H=\operatorname{Stab}_G(p)$, so $P\simeq G/H$ and $H$ has finite index.

For every finite-index subgroup $K\subset G$, apply the bijection to $A=G/K$ and the distinguished point $K$. The required map taking $p$ to $K$ forces $H\subset K$. Hence $H\subset N$. The reverse inclusion follows because $H$ itself occurs in the intersection, so $H=N$ and $[G:N]<\infty$.

Conversely, the intersection defining $N$ is normal, since conjugation permutes the finite-index subgroups. For every finite $G$-set $A$ and every $a\in A$, its stabilizer has finite index, so $N$ fixes~$a$. All finite actions therefore factor through $G/N$. If $G/N$ is finite, $(G/N)\times S$ is a finite $G$-set, and an equivariant map from it to $A$ is uniquely specified by the arbitrary function $s\mapsto f(N,s)$. This is the required adjunction, naturally in both variables.\hfill$\square$

For $G=\mathbb Z$, the intersection of the subgroups $n\mathbb Z$ is zero and has infinite index. Thus $U$ has no left adjoint, even though the source of $e\to B\mathbb Z$ has one object and no nonidentity arrows. Equivalently, a hypothetical universal point in a finite $\mathbb Z$-set would have a finite orbit of length $r$, but could not map to the distinguished point in the cycle of length $r+1$. This disproves the unrestricted finite-colimit reading of the sentence on manuscript p.~66.

The criterion also separates existence of a left adjoint from the pointwise colimit formula. An infinite group with no nontrivial finite quotients has $N=G$, so all its finite actions are trivial and $U$ is an equivalence. Its left adjoint exists although the infinite discrete coproduct appearing in the proposed formula does not. In the general theory, Dér~5 bis is therefore substantive even when Dér~4 bis has been obtained.

We did not find this explicit finite-action criterion in the derived-functor and ABC sources checked here, nor in a targeted search for adjoints of the finite-action forgetful functor. The elementary argument is supplied as a diagnostic continuation of the text, with no claim of priority. The repair needed for Grothendieck's construction is to retain the transfinite colimit hypotheses, or to require that each comma category and its colimit remain within the range of the finite theorem.
\end{dcomment}


\msp{page 67}
Mais il se présente à nouveau la perplexité si le foncteur (3.3.1) ainsi construit est bien adjoint à gauche de $u_W^*$. On peut songer à utiliser encore la même factorisation en $u = vw$, en
\printedpage{33}
% head: A. Grothendieck : Dérivateurs, Ch. XVII, Édition des Universités de Montpellier II et Paris VII
% foot: Version 9 février 2019
utilisant le mapping-cone. La difficulté, cependant, c'est qu'à présent $J = C(u)$ n'est plus une catégorie ordonnée, donc on n'a pas une définition toute prête de $C_{J^\circ}\subset\Fl(\M(J^\circ))$, donnant lieu à une sous-catégorie $\M(J^\circ)_{\mathrm{c}}$ etc. En fait, il semblerait qu'on soit pratiquement paralysé, tant qu'on n'arrive à construire sur $\M(J^\circ)$ également une structure de catégorie à cofibrations, se comportant bien par rapport au foncteurs $v_!$. La % [sic] au foncteurs -> aux foncteurs, accord du pluriel
“démonstration” d'\textsc{Anderson} de l'existence de $u_!$ pour \textit{toute} flèche dans Cat, foire totalement (même pour $I\longrightarrow e$, $I$ ensemble ordonné fini – mais là j'ai pu rétablir la situation), car il se ramène au cas où $I$ est ordonné localement fini, et il
\msp{page 68}
admet sans plus que \textit{dans ce cas} $u_!^W$ est adjoint à gauche de $u_W^*$, et \textit{de plus,} que si $u$ est $\mathbf{W}_0$-connexe, donc $u^* : \M(I'^\circ)\longrightarrow\M(I^\circ)$ pleinement fidèle, qu'il en soit de même de $u_W^*$ déduit par passage aux localisés. Or même si on admet que $u_!^W$ est adjoint à gauche de
$u_*^W$, faute de disposer d'un morphisme d'adjonction tout près “en bas” (dans $\Ho_W(I'^\circ)$), % [sic] u_*^W -> u_W^*, notation employée partout ailleurs
déduit du morphisme d'adjonction dans $\M(I^\circ)$, il n'est absolument pas clair que $u_W^*$ soit pleinement fidèle, i.e. $u_!^W$ un foncteur de localisation. Le contrexemple 3.6 montre bien que, même si on a adjonction (i.e. si $e$ est un objet initial de $\mathcal{C} = \mathcal{C}'$), le foncteur $u^* = \id_{\mathcal{C}}$ a beau être pleinement fidèle (et même une équivalence de catégories), le foncteur localisé n'a aucune raison de l'être : cela pouverait que tout objet de Cat ayant un objet initial % [sic] pouverait -> prouverait, faute de frappe
est équivalent à $e$\,!

\msp{page 69}
\textbf{3.11.} Disons qu'on a une \textit{catégorie à cofibrations stricte,} si dans l'axiome $\mathrm{Cof}_3$ de factorisation (cf. p. 39), la factorisation peut être prise fonctoriellement en $f$. Il revient $\pm$ au même de dire que l'axiome de factorisation reste valable, quand on pense à une catégorie $\M(I)$ pour $I$ dans Cat, en y définissant $W(I)$ et $C(I)$ (à ne pas confondre avec $C_I$) fibre par fibre. Cela revient au même tout au moins si $\M$ est équivalente à une petite catégorie. Sinon, il faut supposer fonctorialité pour toute \textit{petite} sous-catégorie pleine de $\underline{\mathrm{Fl}}(\M)$. Il vaut mieux définir la catégorie à cofibrations stricte par cette propriété de stabilité.

On voit alors que le prédérivateur $\mathbf{D} = \mathbf{D}_{\M}$ sur Cat est un dérivateur \textit{relativement à} $\mathcal{D} = \mathrm{Diag}_0$ défini plus haut (catégories ordonnées $I$ artiniennes, telles que $\Esp(I)$ soit noethérien), ou même relativement à $\mathcal{D} = \mathrm{Diag}$ (catégories ordonnées $I$ artiniennes sans plus) lorsque $C$ et $TC$ sont stables par composition transfinie. J'entends par là que $u_!$ existe quand $u$ est une flèche de $\mathcal{D}$,
\msp{page 70}
et que ça reste valable quand on passe à un dérivateur induit $\mathbf{D}_Z$ ($Z$ dans Cat quelconque)\,; en d'autre termes, que $v_!$ existe quand $v$ est de la forme $\id_Z\times u$, où $u : I\longrightarrow I'$ est une % [sic] d'autre termes -> d'autres termes, accord du pluriel
flèche de $\mathcal{D}$, et que le calcul se fait \textit{au dessus de} $I'$ selon la règle Dér 5 bis. Cela signifie donc que la flèche de changement de base pour les $f_!$ est un isomorphismes dans le cas % [sic] un isomorphismes -> un isomorphisme, accord du singulier
suivant :
\[
\begin{tikzcd}[row sep=huge, column sep=huge]
I\times Z \arrow[d, "u\times\id_Z"] & (i'\backslash I)\times Z \arrow[l] \arrow[d, "\pr_2"] \\
I'\times Z & \{i'\}\times Z \arrow[l]
\end{tikzcd}
\ ,
\]
où $i'$ est un objet quelconque de $I'$.

% ---- from der-ch17_p036-040.tex
\printedpage{34}
% head: A. Grothendieck : Dérivateurs, Ch. XVII, Édition des Universités de Montpellier II et Paris VII
% foot: Version 9 février 2019
L'hypothèse \textit{stricte} pour une catégorie à cofibrations me paraît assez anodine, dans la plupart des cas à ma connaissance, elle est satisfaite. Mais même moyennant cette hypothèse, je ne vois aucun moyen de prouver l'existence de $u_!$ pour \textit{toute} flèche de Cat, et c'est une situation bien sûr très gênante. On voudrait absolument dégager des conditions qui assurent la validité de Dér 4 bis et de Dér 5 bis (existence de $u_!$ pour
\msp{page 71}
\textit{toute} flèche de Cat, et mode de calcul des fibres \dots).

\dgrblock{dgrb935}{Strict factorization and arbitrary diagram shapes}
\begin{dcomment}{later}{p70}{Strict factorization and arbitrary diagram shapes}
Grothendieck introduces strictness to make the theorem uniform in an extra parameter. A natural transformation of $Z$-diagrams is itself a functor $Z\to\operatorname{Arr}(\mathcal M)$. A functorial factorization on the arrow category factors all its components together and supplies the transition maps between the intermediate objects. Thus the coefficient category $\mathcal M^Z$, with weak equivalences and cofibrations tested componentwise, again has the required factorization axiom. Colimits and the stated closure properties are also tested componentwise. Applying the earlier theorem to these new coefficients constructs the extension along $u\times1_Z$.

There is a converse behind the qualification about size. If $\mathcal M$ is small, factor the universal arrow diagram indexed by $\operatorname{Arr}(\mathcal M)$; its two endpoint functors and their universal natural transformation recover a functorial factorization on all arrows. For an essentially small category, first choose a small equivalent one. In a large category the all-small-diagrams assertion only makes this choice on every small full subcategory of the arrow category, which is the formulation Grothendieck finally adopts. No global choice on a proper class is needed for the application to a particular small~$Z$.

The arbitrary-shape question at the end of~3.11 was subsequently answered under hypotheses appropriate to all small homotopy colimits. In \citet[Theorems~9.6.3 and~9.6.5]{RadulescuBanu2009}, the existence of the left adjoint and its pointwise comma-category formula are proved separately for ABC cofibration categories; Theorem~10.3.2 identifies the resulting strong left Heller derivator. Here we use the numbering of arXiv:math/0610009v4. Its axioms CF5 and CF6 provide small coproduct and sequential-composition control. The theorem therefore does not assert that the finite axioms alone suffice. The construction proceeds through direct simplicial resolutions of arbitrary categories, rather than requiring a projective Quillen structure on every original diagram category.

\citet[Theorem~6.11 and Corollary~6.21]{Cisinski2010} gives the corresponding presheaf and derivator statements in a fibration-oriented convention, with the required homotopical completeness assumptions; the dual statements concern cofibrations. These references establish the later affirmative results without turning the source's uncertainty into a present-day open problem. They also explain why strictness and specialness should remain distinct: strictness is a coherent factorization condition, while the next section adds a set of generators and transfinite closure to obtain a particular transfer construction.
\end{dcomment}


\section*{4\quad Catégories à cofibrations spéciales}

\textbf{4.1.} Je suppose à présent que
\begin{itemize}[leftmargin=3em]
\item[\circled{a}] $\mathcal{M}$ est stable par $\varinjlim$ filtrantes, et
\item[\circled{b}] que ses objets sont accessibles,
\end{itemize}
de sorte qu'on est sous les conditions d'application du théorème de factorisation (XV, \S1). On suppose :
\begin{itemize}[leftmargin=3em]
\item[\circled{c}] donné $W, C \subset \Fl(\mathcal{M})$, satisfaisant les conditions $\Cof_1$, $\Cof_2$ de page 39. Pour $\Cof_3$, cf. plus bas.
\item[\circled{d}] $C$, $TC \defeq C \cap W$ \textit{stables par composition transfinie} (c'est le cas si la structure provient d'une structure de Quillen close $(W, C, F)$).
\end{itemize}
\[
TF \defeq C_* \;\subset\; \Fl(\mathcal{M})\ , \tag*{(4.1)}
\]
il résulte du théorème de factorisations que toute flèche $f$ de $\mathcal{M}$ se factorise en
\[
f \;=\; pi\ , \qquad p \in TF,\ i \in C\ ,
\]
\textit{pourvu} qu'on fasse l'hypothèse
\begin{itemize}[leftmargin=3em]
\item[\circled{e}] \mbox{}
\[
\exists\, C_0 \subset C,\ \textit{petite} \text{ partie de } C, \text{ telle que } C \subset \tilde{C}_0, \text{ i.e. } TF = (C_0)_*. \tag*{(4.3)}
\]
\end{itemize}
Sous ces conditions, si on suppose
\begin{itemize}[leftmargin=3em]
\item[\circled{f}] \mbox{}
\[
TF \subset W\ , \tag*{(4.4)}
\]
\end{itemize}
\msp{page 72}
alors l'axiome de factorisations $\Cof_3$ est satisfait. Inversement, si $\Cof_3$ est satisfait, et si $W$ est stable par facteurs directs, alors on a 4.4. Donc si $W$ est stable par facteurs directs, $\Cof_3$ équivaut à (4.4).

\dgrblock{dgrb936}{What the small set of cofibrations contributes}
\begin{dcomment}{exposition}{p71}{What the small set of cofibrations contributes}
Grothendieck now changes the method. The induction over a well-founded indexing order is replaced by an induction that attaches solutions to lifting problems. A small set of maps can test all the required right lifting properties at once. This is why the set $C_0$ enters the argument precisely when the indexing category is allowed to have arbitrary arrows and endomorphisms.

For a class $S$ of maps, write
\[
 S^{\perp}=\{p:s\mathrel{\perp}p\text{ for every }s\in S\},
 \qquad
 \widetilde S={}^{\perp}(S^{\perp}).
\]
The tilde denotes closure for the lifting relation, not closure under weak equivalence. If $C_0\subset C\subset\widetilde{C_0}$, then $C^{\perp}=C_0^{\perp}$ by the two inclusions in the lifting relation. Under the smallness hypotheses, $\widetilde{C_0}$ is also the class of retracts of relative $C_0$-cell maps. These are the two descriptions connected by the small-object argument; they should not be identified without its hypotheses.

Here is the factorization mechanism in the form needed in~4.1. Given $f:X\to Y$, begin with $X_0=X$. At a successor stage, form the coproduct of all maps $a:A\to B$ in $C_0$, indexed further by all commutative squares from $a$ to $X_\alpha\to Y$. Push this coproduct out along its maps to $X_\alpha$, obtaining $X_\alpha\to X_{\alpha+1}$ and an induced map $X_{\alpha+1}\to Y$. At a limit stage take the colimit of the preceding objects. The set of generators and local smallness ensure that each attaching family is a set. Choose a regular cardinal large enough for the domains of all the generators, and carry the construction to a limit ordinal of at least that cofinality. Every map from a generating domain to the final object then factors through some stage, so its associated lifting square is solved at the next stage. The resulting map $p:X_\lambda\to Y$ lies in $C_0^{\perp}=C^{\perp}$.

The left factor $i:X\to X_\lambda$ lies in $C$. Indeed, a coproduct of cofibrations can be attached one summand at a time: well-order the indexing set, replace one source summand by its target at each step, and take colimits at limits. Each step is a pushout of a map in $C$, and transfinite closure finishes the argument. The same closure applies to the successive attaching maps and then to their composite. This gives $f=pi$ with $i\in C$ and $p\in C_*$; the extra assumption $C_*\subset W$ supplies the required weak equivalence. For the standard modern formulation, see \citet[Theorem~12.2.2 and Corollary~12.2.4]{Riehl2014}.

The converse on manuscript p.~72 is a short retract argument. Factor $p\in C_*$ as $p=q i$ with $i\in C$ and $q\in W$. In the square having top map $1_{\operatorname{dom}p}$, left map $i$, right map $p$, and bottom map $q$, a lift $h$ satisfies $hi=1$ and $ph=q$. Therefore $p$ is a retract of $q$ in the arrow category. Retract closure of $W$ gives $p\in W$. This is exactly where that additional hypothesis is used.

The accessibility assumption should also be read at its stated strength. It provides a smallness cardinal for each relevant object, and the set $C_0$ permits a common bound for the domains occurring in the construction. It does not, by itself, state that the whole category has a set of presentable generators. Thus a special cofibration category is not automatically being identified here with a combinatorial model category. A second lifting/factorization pair and its compatibility with $W$ are still to be constructed.
\end{dcomment}


\textbf{4.2.} Il me semble que ces conditions, en pratique, ne seront guère jamais remplies, que si $W$ (qui est donné en tout premier lieu, satisfaisant $\Cof_0$ \textsuperscript{(}\footnote[18]{i.e. saturation faible.}\textsuperscript{)}), n'est pas stable par $\varinjlim$
\printedpage{35}
% head: A. Grothendieck : Dérivateurs, Ch. XVII, Édition des Universités de Montpellier II et Paris VII
% foot: Version 9 février 2019
filtrantes. À ce moment, pour construire la structure à cofibrations, on part de $C_0$ petite partie de $\Fl(\mathcal{M})$, prise telle que l'on ait
\[
(C_0)_* \;\subset\; W\ , \tag*{(4.5)}
\]
on prend pour $C$ au choix le saturé $C'_0$ de $C_0$ par cochangement de base et composition transfinie, ou $C = \tilde{C}_0 =$ ensemble des facteurs directs de $C'_0$. Les conditions $\Cof_1$ (stabilités de $C$), $\Cof_3$ (factorisation) sont satisfaites automatiquement, et il reste à s'occuper de $\Cof_2$, qui revient à la stabilité de $C{\cap}W$ par cochangement de base (la stabilité par composition et que ça contienne les isomorphismes, étant immédiates). C'est une condition délicate, et la seule façon standard que je vois pour l'assurer, c'est en supposant
\[
C_0 \;\subset\; \Cof_W\ ; \tag*{(4.6)}
\]
\msp{page 73}
la stabilité de $W$ par $\varinjlim$ filtrantes (donc aussi par facteurs directs) assure alors qu'on a également
\[
C \;\subset\; \Cof_W\ , \tag*{(4.7)}
\]
et par suite
\[
\left\{
\begin{aligned}
& C \cap W \subset \Cof_W \cap W = W^{\mathrm{univ}}, \\
& \text{d'où }\ TC \defeq C \cap W = C \cap W^{\mathrm{univ}}\ ,
\end{aligned}
\right. \tag*{(4.8)}
\]
ce qui assure la stabilité de $C \cap W$ par cochangement de base. La stabilité de $TC$ par composition transfinie résulte de celle de $C$, et de la stabilité de $W$ par composition transfinie.

\textbf{4.3.} Je reviens aux conditions générales\circled{a} à\circled{f} de 4.1, sans supposer nécessairement $W$ stable par $\varinjlim$ filtrantes. Je veux à présent étendre la structure à cofibrations aux $\mathcal{M}(I)$, $I$ quelconque dans Cat. Nous n'avons pas pour ceci (je crois) à utiliser les développements des paragraphes 2, 3, mais nous ferons le lien avec ceux-ci par la suite.

Désignons par $I_0$ la catégorie discrète ayant mêmes objets que $I$, d'où un foncteur canonique
\[
v = v_I \ :\ I_0 \longrightarrow I\ . \tag*{(4.9)}
\]
L'idée est de définir la structure de $\mathcal{M}(I)$ via la structure évidente de $\mathcal{M}(I_0) \simeq \mathcal{M}^{I_0}$, en jouant sur le couple de foncteurs adjoints
\[
\begin{tikzcd}
\mathcal{M}(I_0) \arrow[r, shift left, "\varphi = v_!"] & \mathcal{M}(I) \arrow[l, shift left, "\psi = v^*"]
\end{tikzcd}
\ .
\]
\msp{page 74}
Nous allons présenter les choses de façon un peu plus générale.

\dgrblock{dgrb937}{Why intrinsic cofibrations control the weak equivalences}
\begin{dcomment}{exposition}{p73}{Why intrinsic cofibrations control the weak equivalences}
The difficulty in~4.2 is not producing a large class of cofibrations. Pushouts and transfinite compositions already do that. The problem is keeping those maps that are also weak equivalences stable under arbitrary cobase change. Grothendieck's intrinsic class $\operatorname{Cof}_W$ addresses precisely this issue: a map belongs to it when pushout along it preserves weak equivalences between objects under its source. This class is determined by $W$, while $C_0$ is a chosen supply of maps with which to construct factorizations.

Suppose $W$ is stable under filtered colimits of diagrams of arrows. Pushout stability of $\operatorname{Cof}_W$ follows from the associativity of pushouts. For a transfinite composite of intrinsic cofibrations, push a weak equivalence through the successive stages. At a successor stage it remains a weak equivalence by the defining property; at a limit stage it remains one by filtered-colimit stability. This proves the transfinite closure needed to pass from $C_0$ to its cellular closure.

The same hypothesis on $W$ also gives retract closure, by a useful elementary argument. If an arrow $f$ is a retract of an arrow $g\in W$, let $e$ be the corresponding idempotent endomorphism of $g$ in the arrow category. The sequence
\[
 g\xrightarrow{e}g\xrightarrow{e}g\xrightarrow{e}\cdots
\]
has colimit $f$. To see this, compose the retraction and section maps with the defining cocones; the identities for a split idempotent establish both universal properties. All objects of this sequence lie in $W$, so its colimit does as well. Retracts of intrinsic cofibrations are then intrinsic cofibrations: any pushout comparison for the retract is a retract of the corresponding comparison for the original arrow. This proves the passage from the cellular closure to $\widetilde{C_0}$.

The equality $\operatorname{Cof}_W\cap W=W^{\mathrm{univ}}$ used in the text relies here on the factorization already obtained, and admits a direct verification. Let $a:A\to B$ be both intrinsic and weakly invertible, and take any map $A\to A'$. Factor the latter as $A\xrightarrow c Z\xrightarrow q A'$ with $c\in C\subset\operatorname{Cof}_W$ and $q\in W$. Pushing $a$ along $c$ gives a weak equivalence $Z\to B\amalg_A Z$. Since the pushout $Z\to B\amalg_A Z$ of $a$ is itself intrinsic, pushing the weak equivalence $q$ along it gives a weak equivalence
\[
 B\amalg_A Z\longrightarrow B\amalg_A A'.
\]
The commutative square with these two maps and $q$ shows, by two-out-of-three, that $A'\to B\amalg_A A'$ is a weak equivalence. Conversely, if all pushouts of $a$ are weak equivalences, the two horizontal pushouts of $a$ in any cobase-change comparison are in $W$; two-out-of-three shows that the comparison is in $W$ whenever its input is. Thus $a$ is intrinsic.

Consequently $TC=C\cap W^{\mathrm{univ}}$ is pushout-stable. Its transfinite stability follows from that of $C$ and the limit-stage hypothesis on $W$. The negative wording ``n'est pas stable'' at the end of manuscript p.~72 is printed in the edition but conflicts with the positive stability hypothesis used throughout this argument. We retain the wording in the source and record the intended sign in the correction register.
\end{dcomment}


\textbf{Définition 4.4.} J'appelle structure de cofibrations \textit{spéciale} une structure $(\mathcal{M}, W, C)$ satisfaisant les conditions a) à f) de 4.1, que je reformule :
\begin{itemize}[leftmargin=3em]
\item[\circled{a}] $\mathcal{M}$ stable par petites limites inductives, et les objets de $\mathcal{M}$ sont accessibles.
\printedpage{36}
% head: A. Grothendieck : Dérivateurs, Ch. XVII, Édition des Universités de Montpellier II et Paris VII
% foot: Version 9 février 2019
\item[\circled{b}] $W$ est faiblement saturé (i.e. $\Cof_0$ de p. 39), $C$ stable par cochangement de base et par composition, $TC \defeq C{\cap}W$ stable par cochangement de base (et par composition automatiquement).
\item[\circled{b$'$}] $C$ et $TC$ sont stables même par composition \textit{transfinie.}
\item[\circled{c}] $\exists$ petite partie $C_0 \subset C$ telle que $C \subset \tilde{C}_0$ ($\Longrightarrow \tilde{C} = \tilde{C}_0$).
\item[\circled{d}] $TF \defeq C_* \subset W$ (ce qui implique, joint à\circled{a} à\circled{c}, l'axiome de factorisation $\Cof_3$ de page 39).
\end{itemize}

\textbf{4.5.} Nous allons donner une recette générale d'extension d'une telle structure, quand on dispose d'un couple de foncteurs adjoints
\[
\begin{tikzcd}
\mathcal{M} \arrow[r, shift left, "v_!"] & \mathcal{M}' \arrow[l, shift left, "v^*"]
\end{tikzcd}
\ ,
\]
$\mathcal{M}$ étant muni d'une structure de cofibrations spéciale $(W, C)$, et $\mathcal{M}'$ satisfaisant la condition liminaire\circled{a}. On définira alors les parties $W'$, $C'$ de $\Fl(\mathcal{M}')$
\msp{page 75}
par
\[
\left\{
\begin{aligned}
W' \;&=\; (v^*)^{-1}(W) \\
C' \;&=\; \text{clôture stable par cochangement de base} \\
&\phantom{\;=\;}\ \ \text{et composition transfinie de } v_!(C),
\end{aligned}
\right. \tag*{(4.5.1)}
\]
donc
\[
v_!(C) \;\subset\; C' \;\subset\; (v_!(C))\postilde \;=\; \tilde{C}'\ . \tag*{(4.5.2)}
\]
Je dis que $(W', C')$, et au choix $(W', \tilde{C}')$, est une structure de cofibrations spéciale sur $\mathcal{M}'$. On va vérifier les différentes conditions de 4.4. La condition a) est déjà acquise.

\dgrblock{dgrb938}{The transfer proof and its exact hypotheses}
\begin{dcomment}{exposition}{p75:La condition a) est déjà acquise.}{The transfer proof and its exact hypotheses}
The adjunction in~4.5 is meant to reconstruct a diagram from its values at objects. Grothendieck first proves the transfer abstractly so that the construction does not depend on a particular indexing category. Write $L=v_!$ and $U=v^*$. The proposed weak equivalences are detected by $U$, while the proposed cofibrations are built from the images under $L$ of the old cofibrations. These are different parts of the construction, and the proof must show that their intersection behaves correctly.

Let $C'=\operatorname{cell}(LC)$ denote closure under pushouts and transfinite composition, including isomorphisms, and let $W'=U^{-1}W$. The two substantive assumptions are
\[
 U\text{ preserves small colimits},\qquad UL(C)\subset C.
\]
Every operation forming $C'$ is then sent to the corresponding operation in the source category. Induction through a cellular construction gives $U(C')\subset C$. Therefore
\[
 U(C'\cap W')\subset C\cap W=TC.
\]
If $a'\in C'\cap W'$ is pushed out in the target, its image under $U$ is a pushout of a trivial cofibration. It remains in $TC$, so the target pushout lies in $W'$ as well as in $C'$. For a transfinite composite of maps in $C'\cap W'$, preservation of the relevant colimits gives a transfinite composite in $TC$; the same conclusion follows. This proves the two closure statements whose compatibility was initially in doubt.

The set of generators is transferred using the left adjoint, not the right one. If $C_0\subset C\subset\widetilde{C_0}$, then $L$ preserves pushouts and transfinite composites, and every functor preserves retract diagrams. Hence
\[
 L(C_0)\subset C'\subset\widetilde{L(C_0)}.
\]
The target smallness hypothesis permits the small-object argument for this new set. Notice that the assertion is about a generating set for the lifting closure; it does not claim that every map of $C'$ is a cell complex on $L(C_0)$ without taking retracts.

Finally, right lifting properties pass across an adjunction exactly. A commutative square from $L(a)$ to $p'$ corresponds to a square from $a$ to $U(p')$, and diagonal fillers correspond under the same bijection. Since left lifting properties survive pushouts, transfinite compositions and retracts,
\[
 (C')^{\perp}=(LC)^{\perp}
 =U^{-1}(C^{\perp})=U^{-1}(TF)\subset U^{-1}(W)=W'.
\]
The small-object factorization therefore has a weakly invertible right factor. These arguments prove Proposition~4.6 without assuming that $U$ preserves all weak equivalences by some separate theorem: that property is part of the definition of $W'$.

For the optional replacement of $C'$ by $\widetilde{C'}$, the proof needs $C$ to be retract-closed. Then $U(\widetilde{C'})\subset C$, and the pushout and transfinite arguments apply again to $\widetilde{C'}\cap W'$. This does not require silently adding retract closure of $W'$. Footnote~20 offers a different route, developed in Proposition~5.1, using additional stability assumptions on $W$. The announcement in~4.5 should be read with the hypotheses collected in the actual proposition, not as a theorem preceding them.

In the diagram application, $U$ is evaluation on the discrete object set and genuinely does preserve colimits. The composite $UL$ is a coproduct of the old objects indexed by arrows of the indexing category. The stability of $C$ under such coproducts, proved by attaching summands as above, is what makes the abstract assumption $UL(C)\subset C$ usable there. The next question concerns trivial cofibrations and is not answered merely by this cofibration calculation.
\end{dcomment}


b) \textsuperscript{(}\footnote[19]{\textbf{N.B.} Il vaut mieux commencer par prouver c).}\textsuperscript{)}. La saturation faible de $W'$ résulte de celle de $W$, via (4.5.1). La stabilité de $C'$ par cochangement de base et par composition (même transfinie) est acquise par définition de $C'$, et itou pour $\tilde{C}'$ (étant Q-close à gauche). Stabilité de $TC' \defeq C'{\cap}W'$ par cochangement de base ? Il suffit de voir que si on a un carré cocartésien
\[
\begin{tikzcd}
X'_0 \arrow[r, "i'"] \arrow[d] \arrow[dr, phantom, "\text{\scriptsize cocart.}"] & X'_1 \arrow[d] \\
X'_2 \arrow[r, "j'"] & X'_3
\end{tikzcd}
\]
dans $\mathcal{M}'$, avec $i'$ dans $TC' = C' \cap W'$, alors $j' \in W'$ (car on sait déjà $j' \in C'$), i.e.
\marginpar{$||$}% marginal sign || printed in the left margin of this line
$v^*(j') \in W'$. Nous allons supposer que $v^*$ \textit{commute aux sommes amalgamées,} et on est donc ramené à prouver que $v^*(i') \in TC$. Donc % [sic] $v^*(j') \in W'$ -> $v^*(j') \in W$, since $W' = (v^*)^{-1}(W)$
\printedpage{37}
% head: A. Grothendieck : Dérivateurs, Ch. XVII, Édition des Universités de Montpellier II et Paris VII
% foot: Version 9 février 2019
\msp{page 76}
on doit prouver le

\textbf{Lemme 4.5.2.} \textit{Sous les conditions} (4.5.3) \textit{et} (4.5.4) \textit{ci-dessous, on a}
\[
v^*(C') \;\subset\; C\ , \qquad v^*(TC') \;\subset\; TC\ ,
\]
\textit{et si $C = \tilde{C}$, on a $v^*(\tilde{C}') \subset C$, $v^*(T\tilde{C}') \subset TC$.}

Pour $v^*(C') \subset C$, il faut d'abord prouver
\[
v^*(v_!(C)) \;\subset\; C\ , \tag*{(4.5.3)}
\]
et je ne vois aucun moyen -- il faut le poser comme hypothèse. De plus, que $v^*$ commute aux $\varinjlim$ quelconques (cochangement de base déjà admis -- à présent on a besoin des limites inductives filtrantes)
\[
v^* \text{ commute aux petites } \varinjlim \text{ quelconques} \tag*{(4.5.4)}
\]
pour conclure de (4.5.3) que $v^*(C' =$ clôture de $v_!(C)$ pour cochangement de base et composition transfinie) est également dans $C$. Passant à $v^*(TC') \subset TC$, par ce qui précède, $v^*(C') \subset C$, il suffit de prouver $v^*(TC') \subset W$, mais $v^*(TC') \subset v^*(W') \subset W$ par définition de $W'$.

Supposons $C = \tilde{C}$, i.e. $C$ stable par facteurs directs, et prouvons $v^*(\tilde{C}') \subset C$, $v^*(T\tilde{C}') \subset TC$. Par c) plus bas, $\tilde{C}'$ est formé des facteurs directs
\msp{page 77}
d'éléments de $C'$, donc $v^*(C') \subset C$ implique $v^*(\tilde{C}') \subset C$ puisque $C$ stable par facteurs directs.

Prouvons enfin $v^*(T\tilde{C}') \subset TC \defeq C \cap W$. Comme $T\tilde{C}' = \tilde{C}' \cap W'$, cela résulte de $v^*(\tilde{C}') \subset C$, et de $v^*(W') \subset W$.

Prouvons b$'$). La stabilité de $C'$ et de $\tilde{C}'$ par composition transfinie est acquise, il reste à prouver la stabilité de $TC' = C' \cap W'$ et de $T\tilde{C}' = \tilde{C}' \cap W'$. Il suffit de voir que le composé transfinie $f$ en question est dans $W'$, donc que $v^*(f) \in W$, mais cela résulte du fait que $v^*$ commute à la composition transfinie par (4.5.4), et que $v^*$ transforme $TC'$ et $T\tilde{C}'$ dans $TC$.

\textit{Prouvons} c) pour $C'$ (donc aussi pour $\tilde{C}'$). Soit $C_0$ une petite partie de $C$ telle que $C \subset \tilde{C}_0$, i.e. $C$ est contenu dans la clôture de $C_0$ par les trois opérations de cochangement de base, de composition transfinie, des facteurs directs. On voit alors que $v_!(C)$ est dans la clôture de $v_!(C_0)$ pour les opérations, donc aussi $C'$, et on peut alors prendre $C'_0 = v_!(C_0)$.

\textit{Prouvons enfin} d), c'est la même assertion pour $C'$ et pour $\tilde{C}'$, puisque
\[
C'_* \;=\; \tilde{C}'_*\ (\ \defeq\ TF'\ )\ ,
\]
\msp{page 78}
il faut donc prouver
\[
TF' \defeq C'_* \;\subset\; W'\ .
\]
Cela va résulter du
\printedpage{38}
% head: A. Grothendieck : Dérivateurs, Ch. XVII, Édition des Universités de Montpellier II et Paris VII
% foot: Version 9 février 2019

\textbf{Lemme 4.5.3.} \textit{On a $TF' = (v^*)^{-1}(TF)$.}

C'est formel par définition de $TF'$ comme
\[
TF' \;=\; (v_!(C))_*\ ,
\]
de $TF$ comme
\[
TF \;=\; C_*\ ,
\]
et de l'adjonction de $v_!$, $v^*$. Comme $TF \subset W$, on en conclut
\[
TF' \;=\; (v^*)^{-1}(TF) \;\subset\; (v^*)^{-1}(W) = W'\ ,
\]
q.e.d.

Donc on a prouvé ceci :

\textbf{Proposition 4.6.} \textit{Considérons un couple de foncteurs adjoints}
\[
\begin{tikzcd}
\mathcal{M} \arrow[r, shift left, "v_!"] & \mathcal{M}' \arrow[l, shift left, "v^*"]
\end{tikzcd}
\ ,
\]
\textit{où $\mathcal{M}$ est muni d'une structure de cofibrations} spéciale $(W, C)$\textit{. On suppose de plus les conditions suivantes satisfaites.}
\[
\begin{aligned}
&1^\circ)\quad \mathcal{M}' \textit{ stable par petites } \varinjlim\ \textit{, ses éléments sont accessibles.} \\
&2^\circ)\quad v^* \textit{ commute aux petites } \varinjlim\ . \\
&3^\circ)\quad v^*(v_!(C)) \subset C.
\end{aligned} \tag*{(4.6.1)}
\]
\msp{page 79}
\textit{Considérons les parties $W'$, $C'$ de} $\Fl(\mathcal{M}')$ \textit{définies par}
\[
\left\{
\begin{aligned}
W' \;&=\; (v^*)^{-1}(W), \\
C' \;&=\; \textit{clôture de } v_!(C) \textit{ par cochangement de base et par composition} \\
&\phantom{\;=\;}\ \textit{transfinie.}
\end{aligned}
\right. \tag*{(4.6.2)}
\]
\textit{Alors $(\mathcal{M}', W', C')$ est une structure à cofibrations spéciale. Si de plus $C = \tilde{C}$, i.e. $C$ est stable par facteurs directs, alors $(\mathcal{M}', W', \tilde{C}')$ est également une structure à cofibrations spéciale} \textsuperscript{(}\footnote[20]{\textbf{N.B.} C'est là un fait général pour les catégories à cofibrations spéciales, et devrait être explicité plus bas (cf. 5.1). L'hypothèse sur $C$ est inutile ici (ou plutôt, il suffit au lieu de cela que $W$ soit stable par facteurs directs \textit{et} par composition transfinie, ce qui n'est pas une hypothèse plus réjouissante \dots).}\textsuperscript{)}.

\textbf{4.7.} J'aimerais, en vue d'énoncés du type 3.7 (établissant l'existence de foncteurs du type $u_!^W$), avoir une prise sur $TC' = C'{\cap}W'$, et sur $T\tilde{C}' = \tilde{C}'{\cap}W'$, de façon précise, j'aimerais avoir la relation
\[
v_!(TC)\ \subset TC' \subset v_!(TC)\postilde\ , \tag*{(4.7.1)}
\]
voire la relation légèrement plus forte
\[
\underbrace{v_!(TC)\ \subset TC'}_{} \ \overset{\text{tautol.}}{\subset}\ \underbrace{T\tilde{C}' \;\subset\; (v_!(TC))\postilde}_{}\ , \tag*{(4.7.2)}
\]

% ---- from der-ch17_p041-045.tex
\printedpage{39}
% head: A. Grothendieck : Dérivateurs, Ch. XVII, Édition des Universités de Montpellier II et Paris VII
qui se décompose en deux
\[
\left\{
\begin{array}{rcl}
v_{!}(TC) & \subset & TC' \\
T\tilde{C}' & \subset & (v_{!}(TC))\cl.
\end{array}
\right.
\tag*{(4.7.3)}
\]
La moins délicate des deux est la \msp{page 80} première, $v_{!}(TC) \subset TC'$, qui revient, en vertu des définitions, impliquant
\[
v_{!}(TC) \subset v_{!}(C) \subset C' \, ,
\]
à
\[
v_{!}(TC) \;\subset\; W' \; ,
\]
i.e. par définition de $W'$, à
\[
v^{*}v_{!}(TC) \;\subset\; TC \; ,
\tag*{(4.7.4)}
\]
ce qui est une hypothèse de type (4.6.1) $3^{\circ}$), à joindre aux hypothèses (4.6.1).

\dgrblock{dgrb939}{The remaining comparison is a comparison of lifting classes}
\begin{dcomment}{exposition}{p79:qui se décompose en deux}{The remaining comparison is a comparison of lifting classes}
Grothendieck has enough structure to factor arrows in the new category, but he needs more to derive arbitrary left Kan extensions in a controlled way. The further issue is whether the transferred trivial cofibrations are generated, for lifting purposes, by the images of the original trivial cofibrations. It is this question, rather than the already established existence of a special cofibration structure, that occupies~4.7.

The forward inclusion is transparent once the hypotheses are separated. Since $L(TC)\subset L(C)\subset C'$, one has
\[
 L(TC)\subset TC'
 \quad\Longleftrightarrow\quad UL(TC)\subset W
 \quad\Longleftrightarrow\quad UL(TC)\subset TC.
\]
The last equivalence uses $UL(C)\subset C$. Thus equation~(4.7.4) is an additional preservation condition, not a consequence of Proposition~4.6 alone.

Assume it from now on. The inclusion $LC\supset L(TC)$ does not control the reverse comparison $TC'\subset\widetilde{L(TC)}$: a cofibration built from $LC$ may happen to be weakly invertible without presenting itself as an attachment of maps from $L(TC)$. Passing to right lifting classes isolates the exact missing assertion. With $F=(TC)^{\perp}$ and $F'=(TC')^{\perp}$, adjunction gives
\[
 (L(TC))^{\perp}=U^{-1}(F),\qquad
 F'\subset U^{-1}(F).
\]
The reverse inclusion is equivalent to $TC'\subset{}^{\perp}((L(TC))^{\perp})$, which is precisely the desired Quillen-closure inclusion. The equivalence is a formal Galois calculation and does not assume an unproved identification of a lifting closure with a cellular closure.

The same argument applies after retract completion, replacing $TC'$ by $\widetilde{C'}\cap W'$. Proposition~5.1 will compare these two trivial classes; Proposition~5.8 will provide the stronger transfer theorem with the required factorization hypotheses. The question on manuscript p.~80 should therefore be read as a problem that directs the next section. The later propositions answer it under further hypotheses; they do not assert that every special cofibration structure already satisfies this stronger comparison.
\end{dcomment}


La relation hypothétique
\[
\underbrace{T\tilde{C}'}_{=\,\tilde{C}'\cap W'} \;\subset\; (v_{!}(TC))\cl, \quad \text{ou seulement} \quad TC' \subset (v_{!}(TC))\cl,
\tag*{$(*)$}
\]
m'apparaît de nature nettement plus délicate. Comme $T\tilde{C}' \supset TC' \supset v_{!}(TC)$, elle équivaut à l'\textit{égalité} de $(T\tilde{C}')_{*}$ resp. $(TC')_{*}$ et $(v_{!}(TC))_{*} = (v^{*})^{-1}(TC)_{*}$. Posons
\[
\left\{
\begin{array}{lcl}
(TC)_{*} & \defeq & F \\
(TC')_{*} & \defeq & F' \;\supset\; (T\tilde{C}')_{*} \;=\; F'_{\sim} \; .
\end{array}
\right.
\tag*{(4.7.5)}
\]
On a donc, en vertu de l'inclusion
\[
v_{!}(TC) \;\subset\; TC' \; ( \;\subset\; T\tilde{C}' \; ) \; ,
\]
la relation
\[
F'_{\sim} \;\subset\; F' \;\subset\; (v^{*})^{-1}(F) \; ,
\tag*{(4.7.6)}
\]
et la question est donc si on a égalité $F' = (v^{*})^{-1}(F)$, voire $F'_{\sim} = (v^{*})^{-1}(F)$ (d'où $F'_{\sim} = F'$).

\msp{page 81}

En fait, on montre assez facilement qu'on a $F'_{\sim} = F'$, i.e. $T\tilde{C}' \defeq \tilde{C}' \cap W' \subset (C' \cap W')\cl$ -- c'est une histoire générale de catégories à cofibrations spéciales \textsuperscript{(}\footnote[21]{cf. prop. 5.1 ci-dessous.}\textsuperscript{)}. Mais je ne vois guère comment prouver (sous des hypothèses convenables) que $F' = (v^{*})^{-1}(F)$. Il me faudra revenir là-dessus au besoin \textsuperscript{(}\footnote[22]{voir cependant prop. 5.8 ci-dessous.}\textsuperscript{)}. Pour l'instant, je me bornerai à récolter le fruit de 4.6 dans un cas particulier :

\dgrblock{dgrb940}{Free diagrams and the missing coproduct lemma}
\begin{dcomment}{exposition}{p81:Pour l'instant, je me bornerai}{Free diagrams and the missing coproduct lemma}
Grothendieck now has a way to give a diagram its cofibrations without choosing an order in which to construct its objects. The price of this freedom is to begin with a diagram having no relations, then attach relations by pushouts. This is the purpose of the adjunction attached to the inclusion of the objects of $I$. Write $L=v_!$ and $R=v^*$. Since $\mathcal M(I)=\operatorname{Fun}(I^{\mathrm{op}},\mathcal M)$, the free diagram on a family $(X_j)$ is
\[
 (LX)(i)=\coprod_{j\in\operatorname{Ob}I}
             \operatorname{Hom}_I(i,j)\cdot X_j.
\]
Here $S\cdot X=\coprod_{s\in S}X$ is a copower. The exponent in the source denotes this coproduct, not a product or an internal Hom. An arrow $i'\to i$ acts by precomposition on the indexing sets $\operatorname{Hom}_I(i,j)$. Thus an arbitrary collection of maps $X_j\to Y(j)$ extends uniquely to a natural transformation $LX\to Y$: on the summand indexed by $i\to j$, use the corresponding structure map of the contravariant diagram $Y$. This also verifies the variance of the formula directly.

The assertion that the product category has a special cofibration structure contains a smallness argument worth making explicit. If $C_0$ is a generating set for the right lifting tests in $\mathcal M$, put a map of $C_0$ in one coordinate and the identity of the initial object in every other coordinate. The union of these sets over $\operatorname{Ob}I$ is still a set. A family of maps has the right lifting property against it exactly when every component lies in $C_*$. This gives the required generating set for the product. Colimits and the stated closure properties are computed componentwise. Likewise, a diagram of accessible objects is accessible: choose a regular cardinal larger than the cardinality of the arrows of $I$ and larger than the smallness bounds for all its components. The naturality equations express its Hom functor as an equalizer of two products of component Hom functors. At this cardinal, the relevant filtered colimits of sets commute with those products and the equalizer. The same argument works with sufficiently long ordinal chains in the small-object formulation of accessibility. No common small generating family for the whole category $\mathcal M$ is required by this argument.

The coproduct assertion singled out in footnote~25 follows from precisely the two closure operations already assumed. Suppose $\mathcal A$ contains the isomorphisms and is stable under pushout and continuous transfinite composition. Given a set of maps $a_s:A_s\to B_s$ in $\mathcal A$, well-order its index set as an ordinal $\lambda$ and form
\[
 X_\alpha=
 \left(\coprod_{s<\alpha}B_s\right)
 \amalg
 \left(\coprod_{\alpha\leq s<\lambda}A_s\right)
 \qquad(\alpha\leq\lambda).
\]
The arrow $X_\alpha\to X_{\alpha+1}$ is the pushout of $a_\alpha$ along its coproduct inclusion. At a limit ordinal the displayed object is the colimit of the previous ones, as follows by specifying a map out of each summand. Therefore $X_0\to X_\lambda$, which is $\coprod_s a_s$, lies in $\mathcal A$. The empty family gives the identity of the initial object. Apply this argument separately to $C$ and to $TC$. It proves that $RL$ preserves both classes and supplies the lemma requested in the footnote. This is the familiar cellular coproduct argument; compare \citet[Lemma~2.1.13]{Hovey1999}.

The outcome of Proposition~4.8 is consequently quite concrete. A freely attached diagram cell is a family of copies of a coefficient cell, one for each incoming arrow at an object. Closing under pushouts and transfinite attachment produces diagram cofibrations; the right lifting tests are pointwise. The first formula (4.8.1) has $TC(I)$ where both occurrences must be $TC(I_0)$, since $RL$ is an endofunctor of the category of families. The proof itself uses the correctly typed assertion.
\end{dcomment}


\textbf{Proposition 4.8.} \textit{Soit} $(\mathcal{M}, W, C)$ \textit{une catégorie à cofibrations spéciale} (déf. 4.4, page 94)\textit{. Soit $I$ objet de} Cat\textit{, $I_0$ la catégorie discrète associée, considérons $v : I_0 \longrightarrow I$ et} % [sic] page 94 -> page 74, apparently: the same déf. 4.4 is cited as "page 74" on printed p. 42
% foot: Version 9 février 2019
\printedpage{40}
% head: A. Grothendieck : Dérivateurs, Ch. XVII, Édition des Universités de Montpellier II et Paris VII
\[
\begin{tikzcd}
\mathcal{M}(I_0) \arrow[r, shift left, "v_{!}"] & \mathcal{M}(I) \arrow[l, shift left, "v^{*}"]
\end{tikzcd} \; ,
\]
\textit{enfin munissons $\mathcal{M}(I_0) = \mathcal{M}^{I_0}$ de la structure de catégorie à cofibrations,} produit \textit{de celles des facteurs, i.e. définissons $W(I_0)$, $C(I_0)$ “fibre à fibre”. (Cela coïncide avec la définition générale donnée au paragraphe} 2\textit{, pour $I_0$ muni de la structure d'ordre discret, et fait de $\mathcal{M}(I_0)$ une catégorie à cofibrations spéciale.) Alors le couple $(v_{!}, v^{*})$ satisfait} \msp{page 82} \textit{aux conditions de la proposition} 4.6\textit{, ainsi qu'à la condition supplémentaire à} (4.6.1)\textit{,}
\[
4^{\circ}) \qquad v^{*}(v_{!}(TC(I))) \;\subset\; TC(I) \; .
\tag*{(4.8.1)}
\]
\textit{Donc on trouve une structure de catégorie à cofibrations spéciale sur $\mathcal{M}(I)$, en posant}
\[
W(I) \;=\; (v^{*})^{-1}(W(I_0))
\tag*{(4.8.2)}
\]
\textit{(c'est en fait équivalent à la définition habituelle de $W(I)$) et}
\[
C(I) \quad = \quad
\begin{array}[t]{l}
\textit{clôture de } v_{!}(C(I_0)) \textit{ pour les opérations de cochangement de}\\
\textit{base et composition transfinie.}
\end{array}
\tag*{(4.8.3)}
\]
\textit{On a de plus sous ces conditions :}
\[
\underbrace{TF(I)}_{\defeq\,(C(I))_{*}} \;=\; (v^{*})^{-1}( \; \underbrace{TF(I_0)}_{\substack{\defeq\,(C_{I_0})_{*},\\ \textit{se définit aussi}\\ \textit{fibre par fibre}}} \; ) \; ,
\tag*{(4.8.8)}
\]
\[
\left\{
\begin{array}{rcl}
v_{!}( \, \overbrace{(TC)(I_0)}^{C(I_0)\cap W(I_0)} \, ) & \subset & TC(I) \qquad\qquad\qquad \text{\fnm{23}} \\
F(I) \defeq (TC(I))_{*} & \subset & (v^{*})^{-1}(F(I_0)).
\end{array}
\right.
\tag*{(4.8.9)}
\]
\footnotetext[23]{Ces inclusions devraient être des égalités, mais je ne sais pas le prouver sans faire des hypothèses plus fortes, du type \textit{catégorie de Quillen spéciale.}}

\textsc{Démonstration}. On sait déjà que $\mathcal{M}(I)$ est stable par petites $\varinjlim$ , et que tout élément est accessible, comme il résulte aussitôt des propriétés idoines de $\mathcal{M}$. Il est clair aussi que $v^{*}$ commute aux petites $\varinjlim$ . Il reste à prouver (4.6.1. $3^{\circ}$), et $4^{\circ}$ dans (4.8.1), i.e. \msp{page 83} que $v^{*}v_{!} : \mathcal{M}(I_0) \longrightarrow \mathcal{M}(I_0)$ applique $C$ dans $C$, $TC$ dans $TC$. Or on a
\[
v_{!}(F_0)(i) \;=\; \varinjlim_{(\alpha,j)\in i\backslash I_0} F_0(j) \qquad\quad \text{\fnm{24}}
\]
\footnotetext[24]{où $j \in \mathrm{Ob}\, I_0 = \mathrm{Ob}\, I$, et $\alpha : i \longrightarrow v(j) = j$.}
et $i\backslash I_0$ est la catégorie discrète au dessus de $I_0$, dont la fibre en $j \in \mathrm{Ob}\, I_0 \;=\; \mathrm{Ob}\, I$ est $\mathrm{Hom}_I(i,j)$ :
\[
i\backslash I_0 \;\;=\;\; \coprod_{j\in \mathrm{Ob}\, I} \text{(catégorie discrète associée à $\mathrm{Hom}_I(i,j)$)} \; ,
\]
% foot: Version 9 février 2019
\printedpage{41}
% head: A. Grothendieck : Dérivateurs, Ch. XVII, Édition des Universités de Montpellier II et Paris VII
donc
\[
v_{!}(F_0)(i) \;=\; \coprod_{j\in \mathrm{Ob}\, I = \mathrm{Ob}\, I_0} \underbrace{F_0(j)^{\mathrm{Hom}_I(i,j)}}_{\substack{\text{somme directe}\\ \text{de } \mathrm{Hom}_I(i,j)\\ \text{copies de } F_0(j)}} \, ,
\]
donc on trouve
\[
v^{*}v_{!}F_0( \, \underbrace{\; i \;}_{\in \mathrm{Ob}\, I_0 \,=\, \mathrm{Ob}\, I} \, ) \;=\; v_{!}(F_0)(i) \;=\; \text{\tikz[baseline=0pt]\draw[->] (0,0) .. controls (0.4,0) .. (0.5,0.45);} \quad .
\]
% the curved arrow points up to the underbraced expression F_0(j)^{Hom_I(i,j)} of the preceding display
Comme $C$ et $TC$ sont stables par sommes directes quelconques (l'étant par sommes directes finies, à cause des conditions de stabilité finies sur $C$ et $TC$, et aussi par somme infinies à cause de l'axiome de stabilité transfinie \dots\textsuperscript{(}\footnote[25]{mérite un lemme général \dots}\textsuperscript{)}), il s'ensuit bien que $v^{*}v_{!}$ applique $C(I_0)$ et $TC(I_0)$ en eux-mêmes. % [sic] somme infinies -> sommes infinies, agreement with the plural adjective

Ainsi, la première assertion de 4.8 résulte de 4.6. De plus, (4.8.8) et (4.8.9) ont été vus dans 4.7, dans les conditions générales de 4.6 et en supposant de plus le $4^{\circ}$ de (4.7.4).

\dgrblock{dgrb941}{A free cell does not make every cofibration free}
\begin{dcomment}{examples}{p83:Ainsi, la première assertion}{A free cell does not make every cofibration free}
The two hoped-for equalities in footnote~23 have different meanings. The equality concerning fibrations is established under the stronger assumptions of Proposition~5.8. The literal equality between $v_!(TC(I_0))$ and $TC(I)$ cannot hold in that generality: the class on the left consists of free maps, whereas the class on the right is allowed to attach them to diagrams that already have relations. A small equivariant example makes the distinction visible without any homotopical calculation.

Give $\mathbf{Set}$ the model structure in which every map is a weak equivalence, cofibrations are injections, and fibrations are surjections. We work with the usual axiom of choice. Injections lift against surjections by choosing a preimage for each new element. Conversely, lifting against $\varnothing\to\{*\}$ characterizes surjections, and a map lifting against all surjections is injective, by the retract argument after factoring it through the injection
\[
 A\longrightarrow A\amalg B\longrightarrow B.
\]
Here the second map is $f$ on $A$ and the identity on $B$. Thus both weak factorization systems are the injection--surjection system. It is generated by the single map $\varnothing\to\{*\}$, so this is a special closed triple in Grothendieck's sense.

Let $I=BG$ for the group $G=\mathbb Z/2$, viewed as a one-object category. Contravariant diagrams are right $G$-sets, $R$ forgets the action, and $L(S)=S\times G$ is free. Consider the equivariant inclusion
\[
 k:\{*\}\longrightarrow\{*\}\amalg G,
\]
where $*$ is fixed and $G$ acts regularly on the second summand. This is the pushout of $L(\varnothing\to\{*\})$, so it belongs to $TC(I)$; all maps are weak equivalences in the transferred structure. It also has the requisite lifting property directly. Given a lifting square against an equivariant map whose underlying function is surjective, the image of $*$ is already specified. Choose a preimage of the image of the identity element of the free orbit, and extend that choice equivariantly to $G$. This defines the diagonal and verifies both triangle identities.

On the other hand, $k$ is not isomorphic in the arrow category to $L(a)$ for any injection $a:S\to T$. The domain of $L(a)$ is the free $G$-set $S\times G$. A nonempty free $G$-set has no fixed points, whereas the domain of $k$ has one. Even the identity of the fixed point would already disprove the literal equality. The nonidentity map $k$ shows why the failure belongs to the geometry of attachment rather than to a convention about including identities.

The correct conclusion under the hypotheses of Proposition~5.8 is
\[
 TC(I)=\widetilde{\,v_!(TC(I_0))\,},
 \qquad F(I)=(v^*)^{-1}(F(I_0)).
\]
The tilde in the first equality is indispensable. Theorem~6.1(b) gives exactly this statement. Modern descriptions by representable cells have the same distinction; see \citet[Theorem~12.3.2]{Riehl2014}. We therefore read the first hope in footnote~23 as needing this closure, and the second as the precise question that the next section resolves.
\end{dcomment}


\msp{page 84}

\section*{5\quad Relation des catégories à cofibrations spéciales avec les catégories de Quillen spéciales}

Soit $(\mathcal{M}, W, C)$ une catégorie à cofibrations spéciale. J'introduis
\[
TF \;=\; C_{*}, \qquad\quad TC \;=\; C \cap W, \qquad\quad F \;=\; (TC)_{*} \; ,
\tag*{(5.1)}
\]
donc on a
\[
\begin{tikzcd}[row sep=large, column sep=large]
TC \arrow[r, hook] \arrow[d, leftrightarrow] & C \arrow[d, leftrightarrow, "\shortmid"{marking, sloped}] \\
F & TF \arrow[l, hook']
\end{tikzcd} \; ,
\tag*{(5.2)}
\]
$(C, TF)$ étant un précouple de Quillen clos à droite. On sait alors (XV 4.2 cor. 1, page 31) que $\tilde{C}$ est formé des facteurs directs d'éléments de $C$. Posons
\[
T\tilde{C} \;=\; \tilde{C} \cap W \; ,
\tag*{(5.3)}
\]
d'où
\[
TC \;\subset\; T\tilde{C} \; .
\tag*{(5.4)}
\]

\textbf{Proposition 5.1} \textit{Soit $(\mathcal{M}, W, C)$ une catégorie à cofibrations spéciale. On a $T\tilde{C} \subset (TC)\cl$ (et même $T\tilde{C} \subset (TC)^{\natural}$ \textsuperscript{(}\footnote[26]{\textbf{N.B.} Pour une partie $A$ de $\mathrm{Fl}(\mathcal{M})$, $A^{\natural}$ désigne l'ensemble des facteurs directs d'éléments de $A$.}\textsuperscript{)}), ou ce qui revient au même en vertu de (5.4), $TC$ et $T\tilde{C}$ ont même $Q$-clôture à gauche :}
\[
T\tilde{C} = (T\tilde{C})\cl = (TC)\cl \qquad\quad \textit{(et même } T\tilde{C} \subset (TC)^{\natural}\textit{),}
\tag*{(5.5)}
\]
% foot: Version 9 février 2019
\printedpage{42}
% head: A. Grothendieck : Dérivateurs, Ch. XVII, Édition des Universités de Montpellier II et Paris VII
\textit{d'où}
\[
(\tilde{C})_{*} = C_{*} \, (= TF), \qquad\quad (T\tilde{C})_{*} = (TC)_{*} \, (= F) \qquad\quad \text{\fnm{27}}.
\tag*{(5.6)}
\]
\footnotetext[27]{Ces relations ne dépendent pas de l'hypothèse $W$ stable par facteurs directs.}
\textit{De plus, si $W$ stable par facteurs directs (ou si seulement tout facteur direct d'une flèche de $TC$ est dans $W$), et si de plus $W$ est stable par composition transfinie, alors $(\mathcal{M}, W, \tilde{C})$ est une catégorie à cofibrations spéciale. De plus, $T\tilde{C}$ est $Q$-saturé à gauche, donc} (par 5.5)
\[
T\tilde{C} \;=\; (TC)^{\natural} \; .
\tag*{(5.7)}
\]

\msp{page 85}

\dgrblock{dgrb942}{What the retract argument actually identifies}
\begin{dcomment}{exposition}{p84}{What the retract argument actually identifies}
Grothendieck introduces fibrations here in order to measure the remaining obstruction to constructing a model structure. The cofibration factorization is already available. His task is to decide what happens when the cofibrations are enlarged by retracts and whether the trivial cofibrations admit a second factorization. These are separate questions, and the calculation of Proposition~5.1 answers the first without assuming an answer to the second.

For a class $A$ of arrows write $A^{\perp}$ for its right lifting class and ${}^{\perp}A$ for its left lifting class. In this notation Grothendieck's $Q$-closure is
\[
 \widetilde A={}^{\perp}(A^{\perp}).
\]
It is a closure defined by lifting. By contrast, closure under pushouts, transfinite compositions and retracts is the class now usually called weak saturation. The two coincide for a set of generating maps when the small object argument applies. They must not be identified merely because a class is stable under those three operations. A left lifting class has all three stability properties, but the converse requires an additional argument.

Let $f:A\to B$ lie in $\widetilde C\cap W$. Use the available factorization $f=pi$, with $i:A\to E$ in $C$ and $p:E\to B$ in $TF=C^{\perp}$. Since $TF\subset W$, the two-out-of-three property gives $i\in C\cap W=TC$. The definition $f\in\widetilde C={}^{\perp}TF$ supplies a map $s:B\to E$ with
\[
 ps=\operatorname{id}_B,\qquad sf=i.
\]
In the arrow category under $A$, the maps $(\operatorname{id}_A,s)$ and $(\operatorname{id}_A,p)$ exhibit $f$ as a retract of $i$. This proves the strong, correctly qualified conclusion
\[
 TC\subset\widetilde C\cap W
       \subset\operatorname{Retr}(TC)
       \subset\widetilde{TC}.
\]
Taking $Q$-closures gives
\[
 \widetilde{\,\widetilde C\cap W\,}=\widetilde{TC},
 \qquad (\widetilde C\cap W)^{\perp}=(TC)^{\perp}.
\]
This is all that is needed to identify $F'_{\sim}$ with $F'$ on manuscript p.~81. In the source's proof, the occurrence $f\in C$ and the corresponding diagram label should be read as $f\in\widetilde C$. The lifting argument remains valid with exactly that hypothesis. Its modern general form is the retract argument of \citet[Lemma~11.2.3]{Riehl2014}.

The first equality in (5.5), identifying $\widetilde C\cap W$ with its own $Q$-closure, goes beyond this argument. The preceding prose asserts equality of the two closures, and that is the formulation proved here. Likewise, the equality with $\operatorname{Retr}(TC)$ printed at the start of the proof must at that point be read as the inclusion just established. If retracts of trivial cofibrations belong to $W$, the reverse inclusion follows: a retract of a map in $TC$ belongs to $\widetilde C$ and, by this additional hypothesis, to $W$. Hence
\[
 \widetilde C\cap W=\operatorname{Retr}(TC).
\]
We retain these source displays in the transcription and use their justified forms in the audit. No counterexample to the extra first equality of (5.5) under the full stated assumptions is asserted here.

There is a reason for keeping the retract under its fixed source. Push out the retract diagram in $A\backslash\mathcal M$ along $A\to A'$. Functoriality of pushout preserves its two splitting identities, so the new map is a retract of a pushout of $i\in TC$. That pushout still lies in $TC$. Retract stability of $W$ therefore proves the required weak-equivalence assertion after pushout. Transfinite stability of the enlarged trivial class follows under the further hypothesis used by Grothendieck, namely transfinite stability of $W$ itself. The other conditions for a special cofibration structure pass to $\widetilde C$ because its right lifting class and its generating set have not changed. This proves his stated sufficient result without silently adding the missing second factorization.
\end{dcomment}


\textsc{Démonstration}. Prouvons
\[
T\tilde{C} \;\subset\; (TC)\cl \, ,
\]
et plus précisément, que tout élément de $T\tilde{C}$ est facteur direct d'un élément de $TC$ :
\[
T\tilde{C} = (TC)^{\natural} \; .
\]
Soit donc $f \in T\tilde{C} \defeq \tilde{C}\cap W$, considérons leur factorisation de $f$ en $pi$, avec $i \in C$, $p \in TF$. Comme $TF \subset W$, on a $p \in W$, donc $i \in W$, donc $i \in C \cap W \defeq TC$. Comme $f \in C$, $p \in TF$, il existe $q : Y \longrightarrow X'$ tel que $pq = \mathrm{id}_Y$, $qf = i$,
\[
\begin{tikzcd}[column sep=small, row sep=small]
X \arrow[rr, "f\in C\cap W"] \arrow[dr, "i\in C"'] & & Y \arrow[dl, dotted, shift right, "q"'] \\
& X' \arrow[ur, shift right, "p\in TF"'] &
\end{tikzcd}
\]
donc $f$ apparaît comme facteur direct de $i \in TC$, q.e.d.

Prouvons que $(\mathcal{M}, W, \tilde{C})$ est une catégorie à cofibrations spéciale (déf. 4.4, page 74). Condition a) sur $\mathcal{M}$ acquise, condition b), la partie sur $W$, aussi. La stabilité de $\tilde{C}$ par cochangement de base et par composition OK, aussi (pour b$'$) la stabilité par composition transfinie. Stabilité de $T\tilde{C} = \tilde{C} \cap W$ par cochangement de base : il reste à prouver que si on a un carré cocartésien
\[
\begin{tikzcd}[row sep=large, column sep=large]
X \arrow[r, "i"] \arrow[d] \arrow[dr, phantom, "\text{cocart.}"] & Y \arrow[d] \\
X' \arrow[r, "i'"] & Y'
\end{tikzcd} \; ,
\]
si $i \in \tilde{C}\cap W$, alors $i' \in W$ (car on sait déjà que $i' \in \tilde{C}$). Mais la démonstration précédente montre que c'est facteur direct \textit{dans} $X\backslash\mathcal{M}$ d'une flèche dans $TC$, donc $i'$ est facteur direct d'une flèche de $TC \subset W$. Donc si $W$ est stable par facteurs directs, $i' \in W$.
% foot: Version 9 février 2019

\printedpage{43}
% head: A. Grothendieck : Dérivateurs, Ch. XVII, Édition des Universités de Montpellier II et Paris VII
\msp{page 86}

Cela achève de prouver b). Pour b$'$), il reste à voir que $T\tilde{C}$ est stable par composition transfinie. Je ne vois pas d'autre façon de le voir, qu'en supposant que $W$ lui-même est stable par composition transfinie (ce qui paraît à peine moins exigeant que la stabilité de $W$ par $\varinjlim$ \texttt{[filtrantes]}). Cela achève de prouver b, b$'$). La condition c) sur $C$ implique la même condition sur $\tilde{C}$. Enfin, comme $\tilde{C}_{*} = C_{*}$, la condition d) est la même pour $C$ et $\tilde{C}$. Cela achève de prouver que $(\mathcal{M}, W, \tilde{C})$ est une catégorie à cofibrations spéciale. (\textbf{N.B.} on avait besoin de deux hypothèses supplémentaires sur $W$, stabilité par facteurs directs et par composition transfinie, pour établir que $T\tilde{C}$ est stable par cochangement de base et par composition transfinie). Le dernière assertion de 5.1 et la relation (5.7) qui en résulte, revient à dire que $T\tilde{C} = \tilde{C} \cap W$ est stable par facteurs directs, ce qui est évident par la même propriété de $\tilde{C}$, et l'hypothèse faite sur $W$ de stabilité par facteurs directs (ou seulement que les facteurs directs d'éléments de $T\tilde{C}$ sont dans $W$, donc dans $T\tilde{C}$). % [sic] Le dernière -> La dernière, gender agreement with "assertion"

\msp{page 87}

\textbf{Proposition 5.2.} \textit{Soit $(\mathcal{M}, W, C)$ une catégorie à cofibrations spéciale. Alors avec les notations de} (5.1)\textit{, on a}
\[
W \;=\; TF \circ TC,
\tag*{(5.8)}
\]
\[
TF \;=\; F \cap W.
\tag*{(5.9)}
\]

\textsc{Démonstration}. Soit $f \in W$, écrivons $f = pi$, avec $i \in C$, $p \in TF$, donc $p \in W$, donc $i \in W$, donc $i \in C \cap W = TC$, d'où (5.8).

Prouvons (5.9). On a $TF \subset F \cap W$, prouvons inversement qu'une $f \in F \cap W$ est dans $TF$. Écrivons $f = pi$, avec $p \in TF$, $i \in TC$.
\[
\begin{tikzcd}[column sep=small, row sep=small]
X \arrow[rr, "f\in F\cap W"] \arrow[dr, "i\in TC"'] & & Y \\
& X' \arrow[ur, "p\in TF"'] &
\end{tikzcd}
\]
Comme $f \in F$, $i \in TC$, et $TC \longleftrightarrow F$ par définition de $F$, $\exists\, r : X' \longrightarrow X$ tel que $ri = \mathrm{id}_X$, $fr = p$. Donc $f$ est facteur direct de $p \in TF$, et comme $TF = C_{*}$ est stable par facteurs directs, on conclut que $f \in TF$.

\textbf{Corollaire 5.3.} \textit{Soit $(\mathcal{M}, W, C)$ une catégorie à cofibrations spéciale. Pour que $(W, C, F)$, $(F \defeq (TC)_{*})$ soit un triple de Quillen dans $\mathcal{M}$, il faut et il suffit que le couple $(TC, F)$ satisfasse à l'axiome de factorisation (i.e. que ce soit un précouple de Quillen). Pour que cette structure de Quillen soit clos, il faut et il suffit que $C$ soit stable par facteurs directs.} % [sic] soit clos -> soit close, agreement with "cette structure"

\msp{page 88}

\textsc{Démonstration}. La nécessité de l'axiome de factorisation pour $(TC, F)$, pour avoir un triple de Quillen, est tautologique. Pour la suffisance, ayant déjà les deux précouples de Quillen $(C, TF)$, $(TC, F)$ avec
% foot: Version 9 février 2019

% ---- from der-ch17_p046-050.tex
\printedpage{44}
% head: A. Grothendieck : Dérivateurs, Ch. XVII, Édition des Universités de Montpellier II et Paris VII
\[
\begin{tikzcd}
TC \arrow[r, hook] \arrow[d, leftrightarrow] & C \arrow[d, leftrightarrow] \\
F & TF \arrow[l, hook'] \ ,
\end{tikzcd}
\tag*{(5.2)}
\]
et $TC = C \cap W$, le fait que $(W, C, F)$ est un triple de Quillen est le corollaire 4 p. 36 de XV : Les conditions $TC = C \cap W$, $TF = F \cap W$, $W = TF \circ TW$ sont équivalentes, et équivalentes au fait que $(W, C, F)$ soit un triple de Quillen.

Par la définition de Quillen des triples \textit{clos,} par les propriétés
\[
\underbrace{W = TF \circ TC}_{\substack{\text{vraie pour tous}\\ \text{le triples}\\ \text{de Quillen}}}, \qquad C = (TF)^*, \qquad \underbrace{F = (TC)_*}_{\substack{\text{vraie dans}\\ \text{le cas actuel}}}\ , % [sic] le triples -> les triples, agreement with "tous"
\]
il reste à exprimer que $C = (TF)^*$, ce qui équivaut à $C$ stable par facteurs directs (cf. loc. cit. cor. 1, p. 31), q.e.d.

\textbf{Remarque 5.4.} Les énoncés 5.2 et 5.3, et la relation $T\tilde{C} \subset (TC)^\natural$ dans proposition 5.1, ne font pas appel à la totalité de l'hypothèse que $(\mathcal{M}, W, C)$ soit une catégorie à cofibrations spéciale, mais seulement que l'on ait une catégorie à cofibrations, avec $C \mathrel{\leftarrow\!\!\!+\!\!\!\rightarrow} \underbrace{TF}_{\defeq C_*} \subset W$, % ?? uncertain: arrow symbol between C and TF -- double-headed arrow crossed by a small bar/plus sign
\msp{page 89}
i.e. une catégorie avec $W$, $C$ satisfaisant Cof$_0$, Cof$_1$, Cof$_2$ de la page 39, plus $TF \defeq C_* \subset W$, plus l'axiome de factorisation pour $(C, TF)$ (lequel implique Cof$_3$).

\dgrblock{dgrb943}{The second factorization is the remaining requirement}
\begin{dcomment}{exposition}{p88}{The second factorization is the remaining requirement}
Grothendieck's comparison with Quillen theory becomes short once the role of the two factorizations is isolated. The factorization associated with $C$ already determines which fibrations must be trivial, even before a general fibration factorization has been constructed. Proposition~5.2 is the calculation that makes this assertion exact.

Set $F=(TC)^{\perp}$. If $f\in W$, factor $f=pi$ with $i\in C$ and $p\in TF$. Since $p\in W$, two-out-of-three gives $i\in TC$; conversely, a composite of an arrow in $TC$ and an arrow in $TF$ belongs to $W$. Thus $W=TF\circ TC$. To identify $TF$, take $f:A\to B$ in $F\cap W$ and use this factorization again. The lifting property of $f$ against $i$ gives $r:E\to A$ satisfying
\[
 ri=\operatorname{id}_A,\qquad fr=p.
\]
The maps $(i,\operatorname{id}_B)$ and $(r,\operatorname{id}_B)$ exhibit $f$ as a retract of $p$. A right lifting class is retract-stable, so $f\in TF$. The reverse inclusion is immediate from $TC\subset C$ and $TF\subset W$. We have proved $TF=F\cap W$ without needing $W$ itself to be retract-stable.

It follows that the new requirement for a Quillen triple is precisely a factorization by $(TC,F)$. If $TC_0\subset TC$ is a set with $(TC_0)^{\perp}=(TC)^{\perp}$, the small object argument factors every map as a relative $TC_0$-cell map followed by an arrow of $F$. The left map belongs to $TC$ because this class is stable under pushouts and transfinite compositions; coproducts were accounted for above. This proves Corollary~5.5. The source's formula $W=TF\circ TW$ on p.~88 has $TC$ as its intended second factor.

Once the second factorization exists, $W$ is in fact retract-stable. Here is a proof that does not presuppose this conclusion or retract stability of $TC$. Let $f:A\to B$ be a retract of $w:C\to D\in W$, with splitting pairs $(a,b):f\to w$ and $(r,s):w\to f$, so $ra=\operatorname{id}_A$ and $sb=\operatorname{id}_B$. First suppose $f\in F$. Factor $w=vu$, with $u:C\to E\in TC$ and $v:E\to D\in TF$. Lifting $u$ against $f$ in the square with top $r$ and bottom $sv$ gives $\ell:E\to A$ such that $\ell u=r$ and $f\ell=sv$. The pairs $(ua,b)$ and $(\ell,s)$ now exhibit $f$ as a retract of $v$. Thus $f\in TF\subset W$, because $TF=C^{\perp}$ is retract-stable.

For arbitrary $f$, use the second factorization $f=hg$, with $g:A\to Z\in TC$ and $h:Z\to B\in F$. Push out $g$ along $a$, writing $E=C\amalg_A Z$ and $j:C\to E$, $t:Z\to E$ for the structure maps. Pushout stability gives $j\in TC$. The compatible maps $w$ and $bh$ induce $k:E\to D$ with $kj=w$ and $kt=bh$, so two-out-of-three gives $k\in W$. The compatible maps $gr$ and $\operatorname{id}_Z$ induce $q:E\to Z$ satisfying
\[
 qj=gr,\qquad qt=\operatorname{id}_Z,\qquad hq=sk.
\]
The last identity can be checked on the two pushout summands. Therefore $(t,b)$ and $(q,s)$ exhibit $h$ as a retract of $k$. The special case already proved gives $h\in W$, and hence $f=hg\in W$. This is the argument of \citet[Lemma~11.3.3]{Riehl2014}, observed here to require only the indicated factorization, lifting, and pushout properties.

If $C$ is retract-stable, so is $TC=C\cap W$. Factoring a map of ${}^{\perp}F$ by $(TC,F)$ and retracting its left factor therefore proves $TC={}^{\perp}F$. The first factorization similarly gives $C={}^{\perp}TF$. Thus the source's criterion for a closed triple supplies the full modern lifting characterizations as well. The second factorization is essential in this explanation; this argument cannot be inserted earlier into Proposition~5.1 to justify its extra equality.

The terminology ``special closed triple'' thus points toward cofibrant generation: there is one set of tests for trivial fibrations and another for fibrations. Its smallness assumption is imposed on all coefficient objects. Modern treatments generally impose only the smallness needed for the chosen generators; see \citet[Theorem~2.1.14 and Definition~2.1.17]{Hovey1999}. If the underlying category is locally presentable, the resulting model category is combinatorial in the now customary terminology. Grothendieck's condition that each object is accessible does not itself state local presentability, which also requires an appropriate set of objects generating the category by filtered colimits. Corollary~5.7 should therefore not be translated by replacing the word ``special'' with ``combinatorial'' without that extra hypothesis.
\end{dcomment}


\textbf{Corollaire 5.5.} \textit{Soit $(\mathcal{M}, W, C)$ une catégorie à cofibrations spéciale, telle qu'il existe une} petite \textit{partie $TC_0 \subset TC = C \cap W$, avec}
\[
TC \ \subset \ TC_0\cl\ ,
\tag*{(5.5.1)}
\]
\textit{ou ce qui revient au même, $(TC)\cl = (TC_0)\cl$, i.e.}
\[
\underbrace{F}_{\defeq (TC)_*} \ = \ (TC_0)_*\ .
\tag*{(5.5.2)}
\]
\textit{Alors $(W, C, F)$ est un triple de Quillen. Ce triple est clos si et seulement si $C$ est stable par facteurs directs.}

En effet, le théorème de factorisation dans XV §1 implique que moyennant l'hypothèse faite sur $TC_0 \subset TC$, le couple $(TC, F)$ satisfait à la condition de factorisation. En effet, par les conditions b, b$'$ de la définition 4.4, $TC$ a les propriétés de stabilité voulues (contient les isomorphismes, est stable par cochangement de base et par composition transfinie).

\msp{page 90}
\textbf{Définition 5.6.} Les triples de Quillen déduits comme dans 5.5 d'une catégorie à cofibrations spéciale sont dits \textit{spéciaux.} Ainsi un triple de Quillen $(W, C, F)$ dans $\mathcal{M}$ est spécial s'il satisfait les conditions a) à d) dans la définition (4.4), concernant $(W, C)$ exclusivement, plus les deux conditions suivantes.
% foot: Version 9 février 2019

\printedpage{45}
% head: A. Grothendieck : Dérivateurs, Ch. XVII, Édition des Universités de Montpellier II et Paris VII
\begin{itemize}[leftmargin=3em]
\item[\encircled{c$'$}] $\exists$ petite partie $TC_0$ de $TC$, telle que $TC \subset (TC_0)\cl$.
\item[\encircled{e}] $F = (TC)_*$.
\end{itemize}

\textbf{Corollaire 5.7.} \textit{Soit $(W, C, F)$ un triple de Quillen} clos \textit{dans $\mathcal{M}$. Pour que ce triple soit spécial, il faut et il suffit qu'il satisfasse les conditions suivantes.}
\begin{itemize}[leftmargin=3em]
\item[\encircled{a}] \textit{$\mathcal{M}$ stable par petites $\varinjlim$\,, et ses éléments sont accessibles.}
\item[\encoval{c, c$'$}] \textit{Il existe une petite partie $C_0$ de $C$, $TC_0$ de $TC$, telles que l'on ait}
\[
C \ = \ C_0\cl\ , \qquad TC \ = \ (TC_0)\cl\ .
\tag*{(5.7.1)}
\]
\end{itemize}
En effet, toutes les autres conditions a) à e), y inclus c$'$) de 5.6, sont automatiquement satisfaites.

\msp{page 91}
\textbf{Proposition 5.8.} \textit{Considérons un couple de foncteurs adjoints}
\[
\begin{tikzcd}
\mathcal{M} \arrow[r, shift left, "v_!"] & \mathcal{M}' \arrow[l, shift left, "v^*"] \ ,
\end{tikzcd}
\tag*{(5.8.1)}
\]
\textit{où $\mathcal{M}$ est munie d'un triple de Quillen spécial clos $(W, C, F = (TC)_*)$, avec $TC = C \cap W$} (cf. définition 5.6). \textit{On suppose satisfaites les conditions suivantes} (comparer prop. 4.6, relations (4.6.1), et (4.7.4)) :
\[
\left\{
\begin{array}{ll}
\text{a)} & \mathcal{M}' \ \textit{stable par petites}\ \varinjlim\ \textit{, ses éléments sont accessibles.}\\[1.5ex]
\text{b)} & v^* \ \textit{commute aux petites}\ \varinjlim\ .\\[1.5ex]
\text{c)} & v^*(v_!(C)) \subset C,\ v^*(v_!(TC)) \subset TC.
\end{array}
\right.
\tag*{(5.8.2)}
\]
\textit{Définissons $(W', C', F')$ dans $\mathcal{M}'$} (comparer prop. 4.6, notamment 4.6.2) \textit{par}
\[
\left\{
\begin{array}{lcl}
W' & = & (v^*)^{-1}(W)\\[1ex]
C' & = & v_!(C)\cl\\[1ex]
F' & = & (TC')_*\ , \qquad \textit{où}\ \ TC' \ \defeq \ C' \cap W'\ .
\end{array}
\right.
\tag*{(5.8.3)}
\]
\textit{Alors $(W', C', F')$ est un triple de Quillen clos spécial dans $\mathcal{M}'$, et on a de plus}
\[
F' \ = \ (v^*)^{-1}(F)\ , \qquad TF' \ = \ (v^*)^{-1}(TF)\ ,
\tag*{(5.8.4)}
\]
\[
TC' \ = \ (v_!(TC))\cl\ .
\tag*{(5.8.5)}
\]

\dgrblock{dgrb944}{The two factorizations meet in the transfer theorem}
\begin{dcomment}{exposition}{p91:est un triple de Quillen clos spécial}{The two factorizations meet in the transfer theorem}
Grothendieck now returns to the equality of fibrations that had resisted the argument of \S4. The extra structure supplied by the second generating set does exactly what is needed: it produces a candidate trivial cofibration factorization in the target, and the previously constructed cofibrations then force that candidate to be the correct one. The proof is a transfer argument, but its decisive step is again a retract under a fixed source.

Write $L=v_!$, $R=v^*$, and choose generating sets $I_0\subset C$ and $J_0\subset TC$ in the coefficient category. To avoid the ambiguity between the raw image of a class and its closure, put
\[
 C'=\operatorname{cof}(LI_0),\qquad
 K'=\operatorname{cof}(LJ_0),\qquad
 W'=R^{-1}(W),
\]
where $\operatorname{cof}$ denotes the class of retracts of relative cell maps. The smallness assumptions ensure that these are also the corresponding left lifting classes. Adjointness transports lifting squares and their diagonals, giving
\[
 (LI_0)^{\perp}=R^{-1}(TF),\qquad
 (LJ_0)^{\perp}=R^{-1}(F).
\]
The equality for trivial fibrations is therefore formal. The equality for fibrations will follow only after $K'$ has been identified with $C'\cap W'$.

First, $K'\subset C'$ because $J_0\subset C$. It is also contained in $W'$. Indeed, $RL(TC)\subset TC$ by hypothesis, while $R$ preserves pushouts and transfinite colimits. It follows by induction along an attachment that $R$ carries every relative $LJ_0$-cell map into $TC$. Functors preserve retract diagrams, and $TC$ is retract-stable in the closed coefficient model structure, so $R(K')\subset TC\subset W$. This proves $K'\subset C'\cap W'$. Notice that preservation of weak equivalences alone would not prove this step: it is the cellular construction and the stronger assertion about $RL(TC)$ that control the entire attachment.

For the converse, let $f:A\to B$ belong to $C'\cap W'$. Factor it by the small object argument for $LJ_0$ as
\[
 A\xrightarrow{j}E\xrightarrow{p}B,
 \qquad j\in K',\quad Rp\in F.
\]
The first part of the proof gives $j\in W'$. As $f\in W'$, two-out-of-three gives $p\in W'$, hence $Rp\in F\cap W=TF$. Thus $p\in R^{-1}(TF)=(C')^{\perp}$. The lifting property of $f\in C'$ supplies a section $s:B\to E$ with $ps=\operatorname{id}_B$ and $sf=j$. Consequently $f$ is a retract of $j$, and $K'$ is retract-stable. This proves
\[
 C'\cap W'=K',\qquad
 (C'\cap W')^{\perp}=R^{-1}(F).
\]
Both target factorizations, both lifting properties, and the equality of the target weak equivalences with the prescribed inverse image are now established. The generating sets are $LI_0$ and $LJ_0$.

This presentation also explains the status of Remark~5.9. If the source triple is not closed, use the pushout--transfinite closure $C'$ of $L(C)$ and the relative $LJ_0$-cell class for the left factor. Its image under $R$ still lies in $TC$. The same argument shows that every map in $C'\cap W'$ is a retract of one of these cell maps, which suffices to identify the right lifting class. The cell maps themselves lie in $C'\cap W'$, so their small-object factorization supplies the required factorization for that larger class. One need not first assert that either of these classes is closed under retracts.

In the modern transfer theorem, the condition that the right adjoint preserve all colimits can be replaced by a direct acyclicity condition on the transferred trivial cells, together with the smallness needed to form the two factorizations. This is the form of \citet[Theorem~11.3.2]{Hirschhorn2003}. Grothendieck's hypotheses give a transparent sufficient test for that acyclicity condition. Their main application here is exceptionally well suited to them: restriction from diagrams to their objects computes colimits pointwise, and the free diagram functor produces coproducts of coefficient cells.
\end{dcomment}


\msp{page 92}
\textsc{Démonstration}. Par 4.6, on sait déjà que $(\mathcal{M}', W', C')$ est une catégorie à cofibrations spéciale. (\textbf{N.B.} Ce qui est noté $C'$ ici, était noté $\tilde{C}'$ dans la proposition 4.5. On doit utiliser la dernière assertion de 4.6, où on suppose $C$ stable par facteurs directs, OK car $(W, C, F)$ supposé clos.) Pour prouver que $(W', C', F' \defeq (TC')_*)$ est un triple de Quillen
% foot: Version 9 février 2019
\printedpage{46}
% head: A. Grothendieck : Dérivateurs, Ch. XVII, Édition des Universités de Montpellier II et Paris VII
spécial, il reste, par 5.3 et définition 5.6, à prouver qu'il existe une petite partie $TC'_0$ de $TC'$, telle que
\[
(\ TC'_0 \ \subset\ )\ TC' \ \subset \ (TC'_0)\cl\ .
\tag*{(5.8.6)}
\]
(Ce triple de Quillen sera de plus clos, puisque $C'$ est Q-clos à gauche par définition.) Soit donc $TC_0$ une petite partie de $TC$ telle que
\[
TC \ \subset \ (TC_0)\cl\ ,
\tag*{(5.8.7)}
\]
qui existe par hypothèse sur $(W, C, F)$. On prendra
\[
TC'_0 \ \defeq \ v_!(TC_0)\ .
\tag*{(5.8.8)}
\]
On a vu dans 4.7 qu'on a
\[
v_!(TC) \ \subset \ TC'
\tag*{(5.8.9)}
\]
(résulte de l'hypothèse $v^*(v_!(TC)) \subset TC$ dans (5.8.2. c)), a fortiori
\[
TC'_0 \ = \ v_!(TC_0) \ \subset \ TC'\ ,
\]
\msp{page 93}
il reste à prouver
\[
TC' \ \overset{?}{\subset} \ (TC'_0)\cl\ ,
\tag*{(5.8.10)}
\]
ce qui équivaut à $(TC')\cl \subset (TC'_0)\cl$, ou encore à
\[
(TC'_0)_* \ \overset{?}{\subset} \ (TC')_* \ \defeq \ F'\ ,
\tag*{(5.8.10 bis)}
\]
alors qu'on a a priori l'inclusion inverse (déduite de $TC'_0 \subset TC'$)
\[
F' \ \subset \ (TC'_0)_*\ .
\tag*{(5.8.11)}
\]
Or par adjonction, on a
\[
\begin{aligned}
(\ \underbrace{TC'_0}_{v_!(TC_0)}\ )_* \ &= \ (v_!(TC_0))_*\\
&= \ (v^*)^{-1}(\ \underbrace{(TC_0)_*}_{= (TC)_* \defeq F}\ )\\
&= \ (v^*)^{-1}(F)\ ,
\end{aligned}
\]
donc (5.8.11) s'écrit
\[
F' \ \subset \ (v^*)^{-1}(F)\ , \qquad \text{i.e.}\ \ v^*(F') \ \subset \ F
\tag*{(5.8.12)}
\]
(comparer (4.7.6)), et il faut montrer que cette inclusion est une égalité (sur quoi j'avais buté p. 80, 81 \ldots), ce qui n'est autre que la première des deux relations (5.8.4). Mais c'est sous la forme initiale (5.8.10) (transposée de (5.8.4)) que nous allons prouver la chose, finalement. Notons que l'on a
% foot: Version 9 février 2019
\printedpage{47}
% head: A. Grothendieck : Dérivateurs, Ch. XVII, Édition des Universités de Montpellier II et Paris VII
\msp{page 94}
\[
(\underbrace{v_!(TC_0)}_{TC'_0})\cl \ = \ (v_!(TC))\cl\ ,
\]
puisque $TC$ est la clôture de $TC_0$ pour les opérations de cochangement de base, de composition transfinie, et de facteurs directs, et que $v_!$ commute à ces opérations. Donc on doit montrer simplement
\[
TC' \ \subset \ (v_!(TC))\cl
\tag*{(5.8.12 bis)}
\]
(où il n'est plus question de $TC_0$). Mais le théorème de factorisation s'applique au couple
\[
\Big( (v_!(TC))\cl,\ (v_!(TC))_* = (v^*)^{-1}(F) \Big)\ ,
\]
puisque $(v_!(TC))\cl = (TC'_0)\cl$, $TC'_0$ petit. Si donc on a $f' \in TC' = C' \cap W'$, factorisons $f'$ en $p'i'$, avec
\[
i' \ \in \ (v_!(TC))\cl\ , \qquad p' \ \in \ (v^*)^{-1}(F)\ .
\tag*{(5.8.13)}
\]
\[
\begin{tikzcd}[column sep=small, row sep=small]
X' \arrow[rr, "f' \in C' \cap W'"] \arrow[dr, "i'"'] & & Y' \arrow[dl, dotted, shift right, "q'"'] \\
& X'_1 \arrow[ur, shift right, "p'"'] &
\end{tikzcd}
\tag*{$(*)$}
\]
Notons qu'on a
\msp{page 95}
\[
(v_!(TC))\cl \ \subset \ TC' \ (\ \defeq \ C' \cap W'\ )
\tag*{(5.8.14)}
\]
car $v_!(TC) \subset TC' \subset C'$, donc $(v_!(TC))\cl \subset \tilde{C}' = C'$, et il reste à prouver que $v_!(TC)\cl \subset W'$, i.e. $v^*((v_!(TC))\cl) \subset W$. Or je dis qu'on a même
\[
v^*((v_!(TC))\cl) \ \subset \ TC\ ,
\]
i.e.
\[
v_!(TC)\cl \ \subset \ (v^*)^{-1}(TC)\ .
\tag*{(5.8.15)}
\]
Cela résulte du fait que $TC$ est Q-clos, donc Q-saturé à gauche, ce qui implique, comme $v^*$ commute aux petites $\varinjlim$\,, que $(v^*)^{-1}(T)$ l'est également. D'autre part, il contient $v_!(TC)$, donc aussi le Q-saturé à gauche de $v_!(TC)$. Mais ce dernier est \textit{égal} à la Q-clôture $(v_!(TC))\cl$, par le théorème de factorisation, vu que $TC'_0 \subset v_!(TC) \subset (TC'_0)\cl$, avec $TC'_0$ petit. % [sic] (v^*)^{-1}(T) -> (v^*)^{-1}(TC), the set in question is TC

Ainsi, (5.8.13) et (5.8.14) impliquent $i' \in TC' = C' \cap W'$, donc $i' \in W'$, d'où $p' \in W'$, i.e. $v^*(p') \in W$, et comme par hypothèse
% foot: Version 9 février 2019
\printedpage{48}
% head: A. Grothendieck : Dérivateurs, Ch. XVII, Édition des Universités de Montpellier II et Paris VII
\msp{page 96}
(5.8.13) on a $v^*(p') \in F$, on a $v^*(p') \in TF$, i.e.
\[
p' \ \in \ (v^*)^{-1}(TF) \ = \ TF' \ (\ \defeq \ C'\ )\ , % [sic] C' -> C'_*, cf. (5.8.16) TF' ( = C'_* )
\]
l'égalité $(v^*)^{-1}(TF) = TF'$, i.e.
\[
(v^*)^{-1}(C_*) \ = \ (v_!(C))_*\ ,
\]
résultant de l'adjonction de $v_!$, $v^*$. Au total, on a donc
\[
p' \ \in \ TF' \ (\ = \ C'_*\ )\ ,
\tag*{(5.8.16)}
\]
et comme $f' \in C'$, le diagramme $(*)$, page 114, nous montre qu'il existe
\[
q' \ : \ Y' \longrightarrow X'_1\ ,
\]
avec $p'q' = \mathrm{id}$, $q'f' = i'$. Donc $f'$ est facteur direct de $i' \in v_!(TC)\cl$, donc on a $f' \in v_!(TC)\cl$, d'où la relation
\[
TC' \ \subset \ v_!(TC)\cl\ ,
\]
ce qui achève de prouver que $(W', C', F')$ est un triple de Quillen spécial et clos.

Les deux relations (5.8.4) ont été prouvées chemin faisant (la deuxième étant essentiellement triviale, et la première représentant la difficulté principale dans la démonstration). Reste à établir (5.8.5). Mais on sait déjà (5.8.9) que
\[
v_!(TC) \ \subset \ TC'\ ,
\]
\msp{page 97}
d'où
\[
v_!(TC)\cl \ \subset \ TC',
\]
puisque $TC'$ est Q-clos à gauche (ou parce que ça a été prouvé dans (5.8.14)), on a déjà prouvé l'inclusion en sens inverse (5.8.12 bis), assertion équivalente à (5.8.4) (et qui représentait la principale difficulté).

\textbf{Remarque 5.9.} Si on ne suppose pas que la structure sur $\mathcal{M}$ soit close, alors la démonstration donnée de 5.8 montre qu'on trouve pourtant sur $\mathcal{M}'$ une structure de Quillen spéciale, en définissant cette fois $C'$ comme dans 4.6. Mais je doute qu'on aura jamais à travailler avec des triples de Quillen spéciaux qui ne soient aussi clos. De toute façon, si $W$ est stable par petites $\varinjlim$ filtrantes\footnote[28]{ce qui implique la stabilité par facteurs directs, et implique aussi la même propriété pour $W' \defeq (v^*)^{-1}(W)$.} (et je doute qu'on trouve des triples de Quillen spéciaux en dehors de ce cas), alors si $(W, C, F)$ est un triple de Quillen spécial, il en est de même de $(W, \tilde{C} = C^\natural, F)$, et je présume qu'on a tout intérêt à remplacer le triple initial par celui-ci. % footnote mark printed in the text as superscript (28)
% foot: Version 9 février 2019

\dgrblock{dgrb945}{Special triples need not preserve ordinary filtered colimits}
\begin{dcomment}{development}{p97:Mais je doute qu'on aura jamais}{Special triples need not preserve ordinary filtered colimits}
Grothendieck closes the section by expressing doubt that special triples occur outside the case where weak equivalences are stable under filtered colimits. The smallness conditions do explain why sufficiently filtered colimits behave well. They do not force the bound to be finite. A countably presented coefficient gives a concrete separation, already in abelian groups.

\textbf{A special closed triple with a failing sequential colimit.}\ours\par
Let $A=\mathbb Z[1/2]$ and, in $\mathbf{Ab}$, define
\[
 \begin{split}
 W_A&=\{f\mid\operatorname{Hom}(A,f)\text{ is bijective}\},\\
 I_A&=\{0\longrightarrow A,\quad A\oplus A\xrightarrow{\nabla}A\},\\
 C_A&={}^{\perp}W_A,\qquad F_A=\operatorname{Mor}(\mathbf{Ab}).
 \end{split}
\]
Then $(W_A,C_A,F_A)$ is a special closed Quillen triple. Its weak equivalences are not closed under sequential colimits in the arrow category.

\emph{Proof.} We first identify the right lifting tests. Lifting against $0\to A$ says that every map $A\to Y$ lifts through $f:X\to Y$; equivalently, $\operatorname{Hom}(A,f)$ is surjective. A lifting square from the fold $\nabla$ to $f$ is a pair of homomorphisms $a,b:A\to X$ with $fa=fb$. A diagonal exists exactly when $a=b$: it must be the common map $A\to X$. The right lifting property against the fold is therefore injectivity of $\operatorname{Hom}(A,f)$. We obtain
\[
 I_A^{\perp}=W_A.
\]
The domains of $I_A$ are small. More explicitly, $A$ has the countable presentation with generators $x_n$, $n\geq0$, and relations $2x_{n+1}=x_n$. For a map from this presentation into an $\aleph_1$-filtered colimit, choose representatives of the countably many images of the generators and stages witnessing the countably many relations. Filteredness at $\aleph_1$ provides a common stage at which they define a homomorphism. The same argument applied to the countably many generator equalities proves eventual equality of two such homomorphisms. Thus $A$, and likewise $A\oplus A$, is $\aleph_1$-presentable. Every abelian group has a presentation of some set cardinality, so the analogous argument supplies the accessibility condition for every object of $\mathbf{Ab}$.

The small object argument now gives the weak factorization system $(C_A,W_A)$, with $C_A=\operatorname{cof}(I_A)$. The standard theorem being used is \citet[Theorem~2.1.14]{Hovey1999}; all its smallness conditions have just been checked. Since $W_A$ is the inverse image of the isomorphisms of abelian groups under the functor $\operatorname{Hom}(A,-)$, it contains the isomorphisms, satisfies two-out-of-three, and is closed under retracts.

If $f:X\to Y$ lies in both $C_A$ and $W_A$, apply its left lifting property to the square having $f$ on both vertical sides and identities horizontally. Its diagonal $h:Y\to X$ satisfies $hf=\operatorname{id}_X$ and $fh=\operatorname{id}_Y$. Therefore
\[
 C_A\cap W_A=\operatorname{Is}(\mathbf{Ab}).
\]
The second weak factorization system is consequently $(\operatorname{Is}(\mathbf{Ab}),\operatorname{Mor}(\mathbf{Ab}))$: factor a map as its identity followed by itself, and invert the left map to solve every lifting problem. This verifies all model axioms. Both left classes are defined by lifting and hence are stable under pushout, transfinite composition and retracts. The generating cofibrations are $I_A$; the generating trivial cofibrations form the empty set, whose right lifting class is all arrows and whose left lifting class is the isomorphisms. Together with cocompleteness and accessibility of the coefficient objects, this verifies every condition of a special closed triple.

To see the failure, consider the sequence in the arrow category whose $n$th arrow is $0\to\mathbb Z$ and whose transition square is $(\operatorname{id}_0,\times2)$. Every homomorphism $A\to\mathbb Z$ is zero, because the image of each $x_n$ must be divisible by arbitrarily high powers of $2$. Thus all these arrows lie in $W_A$. Their colimit is
\[
 0\longrightarrow
 \operatorname*{colim}(\mathbb Z\xrightarrow{2}\mathbb Z
              \xrightarrow{2}\cdots)
 \cong\mathbb Z[1/2]=A.
\]
The map $\operatorname{Hom}(A,0)\to\operatorname{Hom}(A,A)$ is not surjective, since its target contains the nonzero identity. The colimit arrow is not in $W_A$, as required.\hfill$\square$

The example does preserve $\aleph_1$-filtered colimits of weak equivalences, by the presentability calculation in the proof. The general distinction is formalized in \citet[Proposition~4.1]{RaptisRosicky2015}: generating cofibrations between $\kappa$-presentable objects give stability of weak equivalences under $\kappa$-filtered colimits in a combinatorial model category. We checked that result, the standard small-object and transfer constructions, and searches for this particular $\operatorname{Hom}(\mathbb Z[1/2],-)$ example; we did not find this precise illustration in the literature. The construction is a specialization of established methods, and the mark \ours{} records the argument assembled here, with no claim of priority.
\end{dcomment}


% ---- from der-ch17_p051-055.tex
\printedpage{49}
% head: A. Grothendieck : Dérivateurs, Ch. XVII, Édition des Universités de Montpellier II et Paris VII
% foot: Version 9 février 2019

\msp{page 98}

\section*{6\quad Application à la construction des dérivateurs (associés aux triples de Quillen spéciaux)}

\textbf{Théorème 6.1.}\footnote[29]{C'est le résultat principal dans cette partie XVII des notes.} \textit{Soit $(W,C,F)$ un triple de Quillen clos spécial sur la catégorie $\mathcal{M}$} (cf. déf. 5.6, p. 90, et cor. 5.7). \textit{Pour tout $I$ dans} Cat\textit{, considérons la catégorie}
\[
\mathcal{M}(I) \;=\; \underline{\mathrm{Hom}}(I^{\circ},\mathcal{M}) \;, \tag*{(6.1.1)}
\]
\textit{et considérons}
\[
W(I),\ F(I) \;\subset\; \Fl(\mathcal{M}(I)) \;,
\]
\textit{définis par}
\[
\begin{aligned}
&\text{\textit{$u : X \longrightarrow Y$ dans $\mathcal{M}(I)$ est dans $W(I)$ (resp. $F(I)$) si et seulement}}\\
&\text{\textit{si $\forall\, i \in \Ob I$, $u(i) : X(i) \longrightarrow Y(i)$ est dans $W$ (resp. $F$).}}
\end{aligned}
\tag*{(6.1.2)}
\]

\begin{itemize}[leftmargin=2.5em]
\item[a)] \textit{Il existe une structure de Quillen close (nécessairement unique) dans $\mathcal{M}(I)$, ayant $W(I)$ et $F(I)$ comme ensembles des quasi-isomorphismes et de fibrations respectivement. On a donc aussi}
\[
\left\{
\begin{aligned}
TF(I) \;&=\; F(I)\cap W(I) \;=\; \{u\in \Fl(\mathcal{M}(I)) \mid u(i)\in TF \ \ \forall\, i\in \Ob I\},\\
C(I) \;&=\; TF(I)^{*}.
\end{aligned}
\right.
\tag*{(6.1.3)}
\]
\textit{Cette structure de Quillen est} spéciale.

\msp{page 99}
\item[b)] \textit{Soit $I_0$ la catégorie discrète ayant même ensemble d'objets que $I$, et}
\[
v_I \ \ \text{\textit{ou}}\ \ v \;:\; I_0 \longrightarrow I
\tag*{(6.1.4)}
\]
\textit{le foncteur canonique. Désignons par}
\[
C(I) \;=\; (TF(I))^{*}, \qquad TC(I) \;=\; C(I)\cap W(I) \;=\; F(I)^{*}
\tag*{(6.1.5)}
\]
\textit{les ensembles de cofibrations resp. de cofibrations triviales dans $\mathcal{M}(I)$, et utilisons les mêmes notations $C(I_0)$, $TC(I_0)$ pour $I_0$. On a alors}
\[
\begin{aligned}
&\text{\textit{$u : X_0 \longrightarrow Y_0$ dans $\mathcal{M}(I_0)$ est dans $C_0$ (resp. $TC_0$) si et seulement si}}\\
&\text{\textit{pour tout $i$ dans $\Ob I$, $u(i) : X_0(i) \longrightarrow Y_0(i)$ est dans $C$ (resp. $TC$).}}
\end{aligned}
\tag*{(6.1.6)}
\]
\textit{Enfin on a}
\[
\left\{
\begin{aligned}
C(I) \;&=\; (v_!(C(I_0)))\tl\\
TC(I) \;&=\; (v_!(TC(I_0)))\tl\,.
\end{aligned}
\right.
\tag*{(6.1.7)} % ?? uncertain: shape of the small superscript mark after the closing parenthesis -- tilde (\tilde{}) or breve (\breve{})
\]

\item[c)] \textit{Soit}
\[
u \;:\; I \longrightarrow I'
\]
\textit{une flèche dans} Cat\textit{, considérons le couple des foncteurs adjoints}
\[
\begin{tikzcd}
\mathcal{M}(I) \arrow[r, shift left, "u_!"] & \mathcal{M}(I') \arrow[l, shift left, "u^{*}"]
\end{tikzcd}
\;.
\tag*{(6.1.8)}
\]

\printedpage{50}
% head: A. Grothendieck : Dérivateurs, Ch. XVII, Édition des Universités de Montpellier II et Paris VII
% foot: Version 9 février 2019

\msp{page 100}
\textit{Ce couple satisfait aux conditions de Quillen pour les couples de foncteurs adjoints entre catégories de Quillen} (Hom. Alg. Chp. I, §4, th. 3, p. 4.5), \textit{de sorte qu'ils définissent par passage aux localisés un couple des foncteurs adjoints}
\[
\begin{tikzcd}
\mathrm{Ho}_W(I) \arrow[r, shift left, "u_!^{W}"] & \mathrm{Ho}_W(I') \arrow[l, shift left, "u^{*}_{W}"]
\end{tikzcd}
\;.
\tag*{(6.1.9)}
\]
\textit{On a, plus précisément, les relations}
\[
\left\{
\begin{aligned}
u_!(C(I)) \;&\subset\; C(I')\\
u_!(TC(I)) \;&\subset\; TC(I')\;,
\end{aligned}
\right.
\tag*{(6.1.10)}
\]
\textit{d'où}
\[
(u_!(\mathcal{M}(I)_{\mathrm{c}})) \;\subset\; \mathcal{M}(I')_{\mathrm{c}}\;,
\tag*{(6.1.11)}
\]
\textit{et de plus}
\[
\begin{aligned}
&\text{\textit{pour toute flèche $f$ dans $\mathcal{M}(I)_{\mathrm{c}}$, telle que $f \in W(I)$, on a}}\\
&\text{\textit{$u_!(f) \in W(I')$, i.e. $(u_!)_{\mathrm{c}} : \mathcal{M}(I)_{\mathrm{c}} \longrightarrow \mathcal{M}(I')_{\mathrm{c}}$ est compatible}}\\
&\text{\textit{avec les localiseurs.}}
\end{aligned}
\tag*{(6.1.12)}
\]
\textit{D'autre part on a}
\[
\left\{
\begin{aligned}
u^{*}(F(I')) \;&\subset\; F(I)\\
u^{*}(TF(I')) \;&\subset\; TF(I)\\
u^{*}(W(I')) \;&\subset\; W(I)\;,
\end{aligned}
\right.
\tag*{(6.1.13)}
\]
\textit{donc}
\[
\left\{
\begin{aligned}
&u^{*}(\mathcal{M}(I')_{\mathrm{f}}) \;\subset\; \mathcal{M}(I)_{\mathrm{f}}\\
&u^{*}_{\mathrm{f}} : \mathcal{M}(I')_{\mathrm{f}} \longrightarrow \mathcal{M}(I)_{\mathrm{f}} \text{ \textit{est compatible avec les localiseurs.}}
\end{aligned}
\right.
\tag*{(6.1.14)}
\]

\dgrblock{dgrb946}{The projective structure and its free diagrams}
\begin{dcomment}{exposition}{p99}{The projective structure and its free diagrams}
Grothendieck has now supplied the structure that was missing from the mapping-cone argument. A diagram can be replaced without discarding its transition maps, and the replacement is good for every left Kan extension. The weak equivalences and fibrations are still tested at individual objects; it is the cofibrations that encode the compatibility needed to make these tests into a theory of diagrams.

We fix the variance before using the theorem. Here
\[
 \mathcal M(I)=\operatorname{Fun}(I^{\mathrm{op}},\mathcal M),
 \qquad
 (u^*X)(i)=X(u(i)).
\]
For $j\in I$, evaluation has a left adjoint $F_j$, given by
\[
 (F_jA)(i)=\coprod_{I(i,j)}A.
\]
A square from $F_jf$ to a diagram map $p$ is exactly a square from $f$ to $p(j)$; the same correspondence identifies their fillers. Thus, if $S$ and $T$ generate the two lifting systems in $\mathcal M$, the sets
\[
 \{F_js:j\in I,\ s\in S\},\qquad
 \{F_jt:j\in I,\ t\in T\}
\]
generate the projective cofibrations and projective trivial cofibrations. The smallness conditions in the preceding sections justify the required factorizations. Formula (6.1.7) is the version using all objectwise cofibrations at once instead of selected generators. This is the modern projective model structure; the construction and its acyclicity check appear in \citet[Theorem~12.3.2]{Riehl2014}.

The Quillen property of $u_!\dashv u^*$ can also be seen without repeating the cellular construction. If $f\in C(I)$ and $p\in TF(I')$, transpose a lifting square from $u_!f$ to $p$ across the adjunction. Its right side is $u^*p$, which is still an objectwise trivial fibration. A filler exists by the definition of $C(I)$ and transposes back. Replacing $(C,TF)$ by $(TC,F)$ gives the second preservation statement. In particular, the conclusion does not depend on $u$ preserving the order properties used earlier to construct cofibrant diagrams.

There is a useful economy in (6.1.9). Restriction preserves every weak equivalence, so its functor on homotopy categories is already defined. Only the left adjoint needs to be derived: for a projectively cofibrant replacement $QX\to X$,
\[
 \mathbf L u_!(X)=u_!(QX).
\]
The adjunction asserts a natural bijection of morphisms in the localized categories, not merely the existence of this formula on objects. This is the precise additional conclusion that the earlier counterexample made necessary. The elementary lifting criterion for a Quillen adjunction and the preservation of weak equivalences between cofibrant objects are recorded in \citet[Lemmas~11.3.11 and~11.3.14]{Riehl2014}.
\end{dcomment}


\msp{page 101}
\item[d)] \textit{Supposons que $I$ soit une catégorie ordonnée, et considérons l'ensemble $C_I \subset \Fl(\mathcal{M}(I))$ défini dans §2}\footnote[30]{Au paragraphe 2, c'était noté $C_{I^{\circ}}$ au lieu de $C_I$.} (cf. page 25, cf. aussi th. 2.1, p. 26, et ses corollaires). \textit{On a alors}
\[
\left\{
\begin{aligned}
C_I \;&\supset\; C(I)\\
TC_I \;&\supset\; TC(I)\;,
\end{aligned}
\right.
\tag*{(6.1.15)}
\]
\textit{et on a égalité quand $I$ est un ensemble ordonné} artinien.
\end{itemize}

\dgrblock{dgrb947}{The order hypothesis follows the variance}
\begin{dcomment}{development}{p101:et on a égalité quand $I$ est un ensemble ordonné}{The order hypothesis follows the variance}
The comparison with the earlier order-theoretic cofibrations is meant to show that the new construction recovers the old one. At this point a change of notation matters: the earlier argument used covariant diagrams on an artinian poset, whereas (6.1.1) now uses contravariant diagrams on $I$. Accordingly, the sufficient condition is that $I^{\mathrm{op}}$ be artinian, or that $I$ be noetherian. Corollary~6.3 states exactly this condition. The word \emph{artinien} in 6.1(d), repeated in the sentence preceding 6.3, is printed in the edition and is retained in the source text.

\textbf{An inverse-sequence counterexample}\ours.
The condition printed in 6.1(d) is false in general. We work in sets with the axiom of choice and take
\[
 W=\operatorname{Mor}(\mathrm{Set}),\qquad
 C=\{\text{injections}\},\qquad F=\{\text{surjections}\}.
\]
These classes give a closed special Quillen structure. Indeed, injections and surjections are mutually determined by lifting, every map $A\to B$ factors as
\[
 A\longrightarrow A\amalg B\longrightarrow B,
\]
and the single map $\varnothing\to\{*\}$ generates the lifting system. The other specialness conditions hold because sets have all small colimits, each set is small for a sufficiently large regular cardinal, and injections are stable under pushout and transfinite composition. Since every map is in $W$, the two required factorizations are the same one.

Let $I=\mathbb N$ with its usual order. It is artinian but not noetherian. In $\mathcal M(I)$ consider $f:\varnothing\to\mathbf1$, where $\mathbf1$ is the constant singleton inverse sequence. Every nonempty upward closed subset of $I$ is a tail $[n,\infty)$. The colimit of $\mathbf1$ over its opposite is a singleton, while the colimit over the empty subset is empty. Therefore every relative-colimit map associated with closed subsets $U\subseteq V$ is one of
\[
 \varnothing\longrightarrow\varnothing,\qquad
 \varnothing\longrightarrow\{*\},\qquad
 \{*\}\xrightarrow{1}\{*\}.
\]
All are injections. In particular $f\in C_I$, whether one uses the defining principal closed-set tests or the equivalent tests for arbitrary closed pairs.

Nevertheless, $f$ is not projectively cofibrant. Define a presheaf $E$ by
\[
 E(n)=\{k\in\mathbb N:k\geq n\},\qquad
 E(n+1)\longrightarrow E(n)
 \quad\text{the inclusion}.
\]
The map $E\to\mathbf1$ is an objectwise surjection, hence a projective trivial fibration. A filler for the square from $\varnothing\to\mathbf1$ to $E\to\mathbf1$ would choose compatible elements $k_n\in E(n)$. Compatibility forces all $k_n$ to be the same integer, which would have to belong to every tail. There is no such integer. Thus $C(I)\ne C_I$.

This example isolates the missing hypothesis: the partial-lift argument must start at objects with no predecessors in the direction of the diagram. For inverse sequences indexed by $\mathbb N$, that direction has no starting object. We checked the standard projective and Reedy constructions in \citet[\S\S12.3 and~14.2]{Riehl2014} and \citet[Observation~4.16]{RiehlVerity2014Reedy}; the example is our explicit diagnostic for the printed variance error, with no claim of priority.
\end{dcomment}


\textsc{Démonstration}. L'assertion a) est un cas particulier de 5.8, appliqué au couple de foncteurs adjoints
\[
\begin{tikzcd}
\mathcal{M}(I_0) \arrow[r, shift left, "v_!"] & \mathcal{M}(I) \arrow[l, shift left, "v^{*}"]
\end{tikzcd}
\;.
\]
Les conditions (5.8.2) sont satisfaites, cela a été vu dans (4.8). La proposition 5.8 nous dit alors l'existence d'une structure de triple de Quillen close sur $\mathcal{M}(I)$, dont l'ensemble des quasi-isomorphismes $W(I)$ et l'ensemble des fibrations $F(I)$, sont donnés par la première formule (5.8.3) et par la première formule (5.8.4), ce qui équivaut aux définitions (6.1.2). (\textbf{N.B.} Le fait que $\mathcal{M}(I_0)$, avec la structure $W(I_0)$, $C(I_0)$, $F(I_0)$ évidente (produit), est
\printedpage{51}
% head: A. Grothendieck : Dérivateurs, Ch. XVII, Édition des Universités de Montpellier II et Paris VII
% foot: Version 9 février 2019
\msp{page 102}
bien une catégorie de Quillen spéciale close, est $\pm$ tautologique.)

\textit{Prouvons} b). La relation (6.1.6) est tautologique. Les relations (6.1.7) sont contenues dans 5.8 (cf. deuxième relation dans (5.8.3) pour $C' = v_!(C)\tl$, i.e. $C(I) = (v_!(C(I_0)))\tl$, et (5.8.5) pour $TC(I) = (v_!(TC(I_0)))\tl$).

\textit{Prouvons} c). La première assertion est, par \textsc{Quillen}, un cas particulier des relations (6.1.10) à (6.1.14), c'est elles qu'il s'agit d'établir. Les relations (6.1.13) et par suite (6.1.14) sont tautologiques en vertu des définitions, donc il s'agit d'établir \textit{(6.1.10)} et (6.1.12).

Pour les formules (6.1.10), j'utilise le diagramme commutatif
\[
\begin{tikzcd}[row sep=large, column sep=large]
I \arrow[d, "u"'] & I_0 \arrow[l, "v"'] \arrow[d, "u_0"] \\
I' & I'_0\;, \arrow[l, "v'"']
\end{tikzcd}
\]
et je note que les foncteurs $u_!^{\mathcal{M}}$, $u_{0\,!}^{\mathcal{M}}$, commutant aux petites limites inductives, se
\msp{page 103}
comportent bien pour l'opération
\[
\Phi \longmapsto \Phi\tl\,,
\]
quand $\Phi$ est une partie de $\Fl(\mathcal{M}(I))$ resp. de $\Fl(\mathcal{M}(I_0))$ qui est Q-accessible, i.e. telle qu'il existe une \textit{petite} partie $\Phi_0 \subset \Phi$ avec $\Phi \subset \tilde{\Phi}_0$, plus précisément, on a
\[
u_!(\Phi\tl) \;\subset\; (u_!(\Phi))\tl\,.
\]
Ceci noté, on a par b)
\[
C(I) \;=\; (v_!(C(I_0)))\tl\,,
\]
d'où
\[
u_!(C(I)) \;\subset\; (\underbrace{u_!(v_!(C(I_0)))}_{(uv)_!(C(I_0))})\tl\,,
\]
et le deuxième membre s'écrit aussi
\[
(v'_!u_{0\,!}(C(I_0)))\tl \;\overset{?}{\subset}\; (v'_!(I'_0))\tl \;=\; C(I')\;. % [sic] (v'_!(I'_0)) -> (v'_!(C(I'_0))), the C( ) around I'_0 apparently omitted
\]
Donc tout revient à montrer que
\[
u_{0\,!}(C(I_0)) \;\subset\; C(I'_0)\;,
\tag*{(6.1.16)}
\]
ce qui résulte aussitôt du calcul de $u_{0\,!}$ à coups de sommes directes, et du fait que $C$ est stable par sommes directes quelconques. Cela prouve donc $u_!(C(I)) \subset C(I')$, et on prouve de même $u_!(TC(I)) \subset TC(I')$, qui se ramène de même au cas particulier
\[
u_{0\,!}(TC(I_0)) \;\subset\; TC(I'_0)\;.
\tag*{(6.1.17)}
\]

\printedpage{52}
% head: A. Grothendieck : Dérivateurs, Ch. XVII, Édition des Universités de Montpellier II et Paris VII
% foot: Version 9 février 2019

\msp{page 104}

Il reste à prouver (6.1.12), et on procède comme pour 3.3 (p. 54), en utilisant le fait qu'une flèche $f$ dans $\mathcal{M}(I)_{\mathrm{c}}$ qui est dans $W(I)$, s'écrit dans $\mathcal{M}(I)_{\mathrm{c}}$ sous la forme $f = pi$, où $i \in TC(I)$, et où $p$ est inverse à gauche (rétraction) d'une flèche $j \in TC(I)$. On aura donc $u_!(f) = u_!(p)u_!(i)$, et $u_!(i) \in TC(I')$ par (6.1.10), et $u_!(p)$ est inverse à gauche de $u_!(j) \in TC(I')$. On a alors $u_!(i), u_!(j), u_!(p) \in W(I')$, d'où $u_!(f) \in W(I')$, q.e.d.

Prouvons enfin d). Compte tenu de la définition de $C_I$, pour établir $C(I) \subset C_I$, il suffit de voir ceci :

\dgrblock{dgrb948}{Relative colimits measure the new attachment}
\begin{dcomment}{exposition}{p104:pour établir $C(I) \subset C_I$}{Relative colimits measure the new attachment}
Grothendieck wants a cofibration of diagrams to remain a cofibration when part of the indexing category has already been assembled. The map to inspect is therefore not just the map between two colimits. It is the remaining attachment after the source diagram and the chosen part of the target have been put together:
\[
 X_V\amalg_{X_U}Y_U\longrightarrow Y_V.
\]
For a cosieve $U\subseteq V$ in $I$, the actual diagram direction is $V^{\mathrm{op}}$, and $U^{\mathrm{op}}$ is a sieve. These are exactly the subsets on which a partial covariant diagram can be constructed before its complement. In these formulas $X_U$ means $\operatorname{colim}_{U^{\mathrm{op}}}X$, and similarly for $V$; the repeated $U$ in the edition's definition of $X_V$ is a slip.

The reduction to an arrow is useful because it separates the categorical bookkeeping from the lifting argument. After restricting to $V$, take the characteristic map $V\to[1]$ whose fibre at $1$ is $U$. Left Kan extension turns the pair of colimits into the arrow $X_U\to X_V$, and turns a map of diagrams into its square of colimits. The relative map of that square is then the displayed attachment. Only the preservation of projective cofibrations by left Kan extension and the arrow-category calculation remain to be checked.

The restriction step is explicitly deferred in the source. It is justified by 6.4(c) and 6.5: a cosieve inclusion is a discrete opfibration, and its right Kan extension is given objectwise by either evaluation or an empty product. Thus restriction preserves the two projective left classes by adjunction. The temporary assumption in the proof of 6.2 is discharged within these pages.
\end{dcomment}


\textbf{Corollaire 6.2.}\footnote[31]{\textbf{N.B.} Ce résultat se généralise, sous les conditions générales dans 6.2.5.} \textit{Soit $I$ quelconque dans} Cat\textit{, soient $U \subset V$ deux fermés dans $I$ (i.e. des ouverts dans $I^{\circ}$), et soit $f : X_{\bullet} \longrightarrow Y_{\bullet}$ une flèche dans $C(I)$ (resp. dans $TC(I)$). Posons encore}
\msp{page 105}
\[
X_U \;=\; \varinjlim_{i\in U} X(i), \qquad X_V \;=\; \varinjlim_{i\in U} X(i)\;, % [sic] second index i\in U -> i\in V, X_V should be the limit over V
\]
\textit{et de même pour $Y$, et considérons le diagramme}
\[
\begin{tikzcd}[row sep=large, column sep=large]
X_U \arrow[r, "\text{trans.}"] \arrow[d, "u_U"'] \arrow[dr, phantom, "\text{cocart.}"] & X_V \arrow[d] \arrow[dd, dotted, bend left=60, "u_V"] \\
Y_U \arrow[r] \arrow[dr, dotted, "\text{trans.}"'] & Y_{U,V} \arrow[d, "u_{U,V}"'] \\
 & Y_V\;.
\end{tikzcd}
\tag*{(6.2.1)}
\]
\textit{Alors la flèche $u_{U,V}$ est dans $C$ (resp. $TC$).}

\textsc{Démonstration} de 6.2. Admettons que la restriction d'une flèche de $C(I)$ resp. $TC(I)$ à un fermé $V$ de $I$ est dans $C(V)$ resp. $TC(V)$ (cf. \dots{} ci-dessous). Cela nous permet de supposer $V = I$. La donnée de $U$, qui est un ouvert de $I^{\circ}$, équivaut à la donnée d'une flèche $v : I^{\circ} \longrightarrow \Dl^1$ (en prenant la fibre en 0), ou encore à la donnée d'une flèche $u : I \longrightarrow (\Dl^1)^{\circ} \simeq \Dl^1$ (en prenant le fibre en 1). Or % [sic] le fibre -> la fibre, gender (cf. "la fibre en 0" just before)
\[
u \;:\; I \longrightarrow \Dl^1\;, \qquad U \;=\; I_1 \overset{\text{déf}}{=} u^{-1}(\{1\})\;.
\tag*{(6.2.2)}
\]

\msp{page 106}
On a alors
\[
X_U \;=\; u_!(X)(1), \quad X_I \;=\; u_!(X)(0)\;,
\]
et le morphisme de transition $X_U \longrightarrow X_I$ est déduit de la flèche $0 \xrightarrow{\ \partial\ } 1$ dans $\Dl^1$, et de même pour $Y$. Soit
\[
\xi \;=\; u_!(X), \qquad \eta \;=\; u_!(Y)\;,
\]
\printedpage{53}
% head: A. Grothendieck : Dérivateurs, Ch. XVII, Édition des Universités de Montpellier II et Paris VII
% foot: Version 9 février 2019
donc le diagramme (6.2.1) s'identifie au diagramme déduit de
\[
\varphi \;:\; \xi \longrightarrow \eta\;, \qquad \varphi \;=\; u_!(f)\;,
\]
comme
\[
\begin{tikzcd}[row sep=large, column sep=large]
\xi(1) \arrow[r, "\xi(\partial)"] \arrow[d, "\varphi(1)"'] \arrow[dr, phantom, "\text{cocart.}"] & \xi(0) \arrow[d] \arrow[dd, dotted, bend left=60, "\varphi(0)"] \\
\eta(1) \arrow[r] \arrow[dr, dotted, "\eta(0)"'] & \tilde{\eta} \arrow[d, "\tilde{\varphi}"'] \\
 & \eta(0) \qquad .
\end{tikzcd}
\tag*{(6.2.3)}
\]
Or par 6.1.10 appliqué à $u : I \longrightarrow \Dl^1$, la relation $f \in C(I)$ resp. $f \in TC(I)$ implique $\varphi \in C(\Dl^1)$ resp. $\varphi \in TC(\Dl^1)$. Donc pour établir 6.1, on est ramené au cas particulier où $I = \Dl^1$, $U = \{1\}$, $V = I$, donc à prouver le

\textbf{Lemme 6.2.4.} \textit{Soit $\varphi : \xi \longrightarrow \eta$ une flèche dans $\mathcal{M}(\Dl^1)$. Si $\varphi \in C(\Dl^1)$ (resp.}
\msp{page 107}
\textit{$\varphi \in TC(\Dl^1)$), alors la flèche $\tilde{\varphi}$ dans le diagramme} (6.2.3) \textit{est dans $C$ (resp. $TC$).}

Plus généralement, et plus précisément :

\textbf{Lemme 6.2.5.} \textit{Soit $\mathcal{M}$ une catégorie stable par sommes amalgamées, et $\Phi$ une partie Q-close à gauche} (on prendra $\Phi = C$, et $\Phi = TC$, dans le cas du lemme 6.2.4). \textit{Soit $\Psi = \Phi_{*}$, et considérons, pour toute $I$ dans} Cat\textit{, $\Phi(I)$, $\Psi\langle I\rangle$ dans $\mathcal{M}(I)$, définis par :}
\[
\begin{aligned}
&\text{\textit{$u : X \longrightarrow Y$ dans $\mathcal{M}(I)$ est dans $\Psi\langle I\rangle$ si et seulement si $\forall\, i \in \Ob I$,}}\\
&\text{\textit{$u(i) : X(i) \longrightarrow Y(i)$ est dans $\Psi$,}}
\end{aligned}
\tag*{(6.2.5.1)}
\]
\[
\Phi(I) \;=\; \Psi\langle I\rangle^{*}.
\tag*{(6.2.5.2)}
\]

Prenons le cas $I = \Dl^1 = \{0 \xrightarrow{\ \partial\ } 1\}$, et soit $u : X \longrightarrow Y$ une flèche dans $\mathcal{M}(\Dl^1)$, correspondant à un carré commutatif
\[
\begin{tikzcd}[row sep=large, column sep=large]
X(1) \arrow[r, "X(\partial)"] \arrow[d, "u(1)"'] & X(0) \arrow[d, "u(0)"] \\
Y(1) \arrow[r, "Y(\partial)"] & Y(0)\;.
\end{tikzcd}
\tag*{(6.2.5.3)}
\]

Considérons le diagramme correspondant

% ---- from der-ch17_p056-060.tex
\noindent\printedpage{54}\msp{page 108}
% head: A. Grothendieck : Dérivateurs, Ch. XVII, Édition des Universités de Montpellier II et Paris VII
\[
\begin{tikzcd}
X(1) \arrow[r, "X(\partial)"] \arrow[d, "u(1)"'] \arrow[dr, phantom, "\text{cocart.}"] & X(0) \arrow[d] \arrow[dd, dotted, bend left=40, "u(0)"] \\
Y(1) \arrow[r] \arrow[dr, dotted, "Y(\partial)"'] & \tilde{Y} \arrow[d, "\tilde{u}"'] \\
 & Y(0)
\end{tikzcd}
\qquad .
\tag*{(6.2.5.4)}
\]

Pour qu'on ait $u \in \Phi(\boldsymbol{\Delta}^1)$, il faut et il suffit que l'on ait
\[
u(0),\ u(1),\ \tilde{u} \ \in\ \Phi\ .
\tag*{(6.2.5.5)}
\]

J'ai déjà vérifié la chose (sous forme duale), mais ne retrouve plus l'endroit (peut-être n'ai-je pas pris des notes). Je n'ai pas envie de rechercher la démonstration ici.

\dgrblock{dgrb949}{The omitted arrow-category calculation}
\begin{dcomment}{exposition}{p108:Je n'ai pas envie de rechercher la démonstration ici.}{The omitted arrow-category calculation}
The short lemma that Grothendieck leaves unwritten is the local lifting step used throughout the remainder of the proof. We give it with precisely the pushout hypothesis he states. This also explains why the three maps in (6.2.5.5) are slightly more than one needs to test.

Let $\Phi={}^\perp\Psi$ and $\Psi=\Phi^\perp$. Write a map of arrows as
\[
\begin{tikzcd}
 A_1\arrow[r]\arrow[d,"f_1"']&A_0\arrow[d,"f_0"]\\
 B_1\arrow[r]&B_0,
\end{tikzcd}
\qquad
 P=A_0\amalg_{A_1}B_1,\qquad t:P\longrightarrow B_0.
\]
Then $f\in{}^\perp(\Psi\langle[1]\rangle)$ if and only if $f_1,t\in\Phi$. When these conditions hold, $f_0\in\Phi$ as well.

For sufficiency, take a square from $f$ to a map of arrows $p:(E_1\to E_0)\to(F_1\to F_0)$ whose two components belong to $\Psi$. First choose a filler $h_1:B_1\to E_1$ against $f_1$ and $p_1$. The given map $A_0\to E_0$ and the composite $B_1\xrightarrow{h_1}E_1\to E_0$ agree on $A_1$. They define $P\to E_0$. Together with $B_0\to F_0$, this is a square from $t$ to $p_0$. A filler $h_0:B_0\to E_0$ solves it. Its restriction to $B_1$ is the prescribed composite, so $(h_1,h_0)$ is a natural filler of the original square.

For necessity of the relative condition, take any square from $t$ to $q:E\to F$ in $\Psi$, with upper map $a:P\to E$ and lower map $b:B_0\to F$. Consider the map of arrows
\[
 (B_1\xrightarrow{a|_{B_1}}E)
 \longrightarrow
 (B_1\xrightarrow{b|_{B_1}}F)
\]
with components $(1_{B_1},q)$. The maps $(f_1,a|_{A_0})$ and $(1_{B_1},b)$ define a square from $f$ to this map. A filler has first component $1_{B_1}$; its second component is exactly a filler for the prescribed square from $t$ to $q$. Thus $t\in\Phi$.

To test $f_1$, one must make an arbitrary square at the first vertex into a square of arrows. Suppose that square has maps $a:A_1\to E$, $b:B_1\to F$ and right side $q:E\to F$ in $\Psi$. Form
\[
 K=F\amalg_{B_1}B_0.
\]
There is a map of arrows $(E\to K)\to(F\to K)$ with components $(q,1_K)$. The given maps at vertex $1$, together with $A_0\to B_0\to K$ and $B_0\to K$ at vertex $0$, form a square from $f$. A filler supplies the desired map $B_1\to E$. Hence $f_1\in\Phi$. This construction uses a pushout in place of any terminal-object argument.

Finally, $A_0\to P$ is the pushout of $f_1$, and $f_0$ is its composite with $t$. A left lifting class is closed under pushout and composition, so $f_0\in\Phi$. This proves the printed criterion, including its redundant component condition. The same two-stage construction is the direct-category case of the relative-latching lifting argument of \citet[Observation~4.16]{RiehlVerity2014Reedy}; no model structure or factorization axiom is used here.
\end{dcomment}


Pour achever de prouver 6.1. d), i.e. l'égalité $C(I) = C_I$ pour $I$ artinien, j'ai envie de prouver encore quelque chose de plus générale que voici : % [sic] plus générale -> plus général, agreement with « quelque chose »

\textbf{Corollaire 6.3.} \textit{Soient $(\mathcal{M}, \Phi, \Psi)$ et $I \longmapsto \Phi(I) \subset \mathrm{Fl}(\mathcal{M}(I))$ comme dans \textup{6.2.5.} Soit $I$ un ensemble ordonné, et considérons aussi $\Phi_I$ défini (dans la situation duale) page \textup{25 (2.2)} (mais avec des notations duales \dots). Je dis que l'on a}
\[
\Phi_I \ \supset\ \Phi(I)\ ,
\tag*{(6.3.1)}
\]
\msp{page 109}\textit{et on a égalité si $I$ est noethérien, i.e. $I^{\circ}$ est artinien.}

\textsc{Démonstration}. La démonstration de $\Phi(I) \subset \Phi_I$ est celle qu'on vient de donner. Il reste à prouver que si $I$ est noethérien, on a $\Phi_I \subset \Phi(I)$. Étant donnée $i : A \longrightarrow B$ dans $\Phi_I$, il faut prouver que $u \in \Phi(I)$, % [sic] u -> i, the arrow in question is i : A -> B
i.e. que pour un diagramme commutatif dans $\mathcal{M}(I)$,
\[
\begin{tikzcd}
A \arrow[r, "\alpha"] \arrow[d, "i"'] & X \arrow[d, "p"] \\
B \arrow[r, "\beta"'] \arrow[ur, dotted, "h"] & Y
\end{tikzcd}
\tag*{(6.3.2)}
\]
avec $p \in \Psi\langle I\rangle$, i.e. $p(i) : X(i) \longrightarrow Y(i)$ dans $\Psi$ pour tout $i$, $\exists\ h : B \longrightarrow X$ laissant le diagramme commutatif. On va revenir aux notations duales, en travaillant dans $\mathcal{M}(I^{\circ})$, avec $I$ artinien. On considère l'ensemble $\mathcal{H}$ des couple % [sic] des couple -> des couples, plural
$(U, h_U)$ d'un ouvert $U$ de $I$ et d'un $h_U : B|U \longrightarrow X|U$, rendant commutatif la restriction à $U$ de (6.3.2). % [sic] commutatif -> commutative, agreement with « la restriction »
On ordonne $\mathcal{H}$ par prolongement, on note que toute partie totalement ordonnée $(U_\alpha, h_{U_\alpha})$ dans $\mathcal{H}$ a une borne supérieure, savoir $U = \bigcup_\alpha U_\alpha$, $h_U$ défini par les $h_{U_\alpha}$ sur les $U_\alpha$.
% foot: Version 9 février 2019

\dgrblock{dgrb950}{The lift is built in the direction of the arrows}
\begin{dcomment}{exposition}{p109:On va revenir aux notations duales}{The lift is built in the direction of the arrows}
The long cube calculation has a simple purpose: when a new vertex is added, its lift must extend the lifts already chosen on all incoming arrows. The relative latching map incorporates that compatibility in a single lifting problem. We formulate the argument in a fixed covariant convention so that its order hypothesis is visible throughout.

Let $J$ be a small artinian poset, let $A,B,X,Y:J\to\mathcal M$, and suppose $p:X\to Y$ is objectwise in $\Psi$. Assume that the relevant colimits exist. For $j\in J$, put
\[
 L_jA=\operatorname{colim}_{k<j}A_k,\qquad
 L_jB=\operatorname{colim}_{k<j}B_k,\qquad
 \ell_j(f):A_j\amalg_{L_jA}L_jB\longrightarrow B_j.
\]
If every $\ell_j(f)$ belongs to $\Phi$, any square from $f:A\to B$ to $p$ has a natural filler.

Indeed, order partial fillers on downward closed subsets by extension. The empty subset supports a filler, and a chain has an upper bound obtained by taking the union of its domains and its compatible component maps. Zorn's lemma gives a maximal partial filler $h_U$. If $U\ne J$, choose a minimal element $j$ of $J\setminus U$. Every $k<j$ lies in $U$, so the partial filler gives a map $L_jB\to X_j$. Together with $A_j\to X_j$ it induces
\[
 A_j\amalg_{L_jA}L_jB\longrightarrow X_j.
\]
The original square and the naturality of the maps make this a square with lower map $B_j\to Y_j$ and right side $p_j$. The lifting relation $\Phi\perp\Psi$ supplies $h_j:B_j\to X_j$. Its restriction along $L_jB\to B_j$ is the map already prescribed by $h_U$, hence it is compatible with every incoming arrow. There are no arrows from $j$ to an old object of $U$, since $U$ is downward closed. The result is a filler on $U\cup\{j\}$, contradicting maximality. Therefore $U=J$.

This proves the difficult inclusion in 6.3 without shortening any of its mathematical content. Grothendieck's $\widetilde B$, $\widetilde i$ and $\widetilde h_0$ are the pushout, relative latching map and upper arrow just displayed. The identities on pp.~113--116 verify that the proposed upper arrow and lower arrow make the required square commute; they hold by the pushout universal property. The sentence at the end must use the given hypothesis $i\in\Phi_I$, not the printed $i\in\Phi(I)$, which is the conclusion.

The passage from $J$ back to the source's contravariant notation is $J=I^{\mathrm{op}}$. Thus $I$ must be noetherian. The colimit assumption is inherited from the use of the classes defined in \S2; pushouts alone, the hypothesis of the arrow lemma, do not ensure that arbitrary $L_jA$ exist. No small-object argument enters this proof. Its induction is legitimate for any small well-founded poset, including those whose ranks are transfinite.
\end{dcomment}


\printedpage{55}\msp{page 110}
% head: A. Grothendieck : Dérivateurs, Ch. XVII, Édition des Universités de Montpellier II et Paris VII
(\textbf{N.B.} L'ensemble $\mathcal{H}$ n'est pas vide, car $(\emptyset, h_{\emptyset})$ est dans $\mathcal{H}$ -- mais c'est contenu dans l'énoncé précédent que toute partie totalement ordonnée de $\mathcal{H}$ y a une borne supérieure \dots). Par Zorn, il existe un élément maximal $(U, h_U)$. À prouver que $U = I$, et on aura terminé. Ou encore que si $(U, h_U)$ tel que $U \neq I$, ce n'est \textit{pas} maximal. Soit $F = I\backslash U \neq \emptyset$ le fermé complémentaire. Par l'hypothèse artinienne sur $I$, $F$ a un élément minimal $i_0$. Soit
\[
U' \ =\ U \cup \overbrace{I/i_0}^{V} \ =\ U \cup \{i_0\}\ ,
\]
on va prolonger $h_U$ en $h_{U'}$ sur $U'$, on aura terminé. Pour ceci, soit $V_0 = V \cap U = V\backslash\{i_0\}$, on est ramené à prolonger la restriction $h_{V_0}$ de $h_U$ à $V_0$, de $V_0$ à $V$. Au total, on peut supposer $I = I/i_0$ muni d'un élément final $i_0$, $U = I\backslash\{i_0\}$, et montrer qu'une $(U, h_U) \in \mathcal{H}$ se prolonge à $I$. Soient $A'$, $B'$, $X'$, $Y'$ les restrictions de $A$, $B$, $X$, $Y$ à $U$ (qu'on pourra aussi noter $I'$). La donné de $A$ % [sic] La donné -> La donnée, missing feminine ending
équivaut à celle de $A'$, celle de $A_0 = A(i_0) \in \mathrm{Ob}\,\mathcal{M}$, et d'une flèche de
\[
A'_{I'} \ \overset{\text{déf}}{=}\ \varinjlim\limits_{i' \in I'} A'(i') \ \xrightarrow{\ \lambda_A\ }\ A_0\ (\ \simeq A_I\ )\ ,
\]
\msp{page 111}et de même pour $B$, $X$, $Y$. Plus précisément, la catégorie $\mathcal{M}(I^{\circ})$ s'interprète comme celle des triples $X = (X', X_0, \lambda_X)$, où $X' \in \mathcal{M}(I'^{\circ})$, $X_0 \in \mathcal{M}$, et $\lambda_X : X'_{I'} \longrightarrow X_0$. À ce compte, la donnée du diagramme (6.3.2) équivaut à celle d'un diagramme carré similaire
\[
\begin{tikzcd}
A' \arrow[r, "\alpha'"] \arrow[d, "i'"'] & X' \arrow[d, "p'"] \\
B' \arrow[r, "\beta'"'] \arrow[ur, dotted, "h'"] & Y'
\end{tikzcd}
\tag*{(6.3.3)}
\]
\textit{dans} $\mathcal{M}(I'^{\circ})$, d'où un diagramme dans $\mathcal{M}$
\[
\begin{tikzcd}
A'_{I'} \arrow[r, "\alpha'_{I'}"] \arrow[d, "i'_{I'}"'] & X'_{I'} \arrow[d, "p'_{I'}"] \\
B'_{I'} \arrow[r, "\beta'_{I'}"'] \arrow[ur, dotted, "h'_{I'}"] & Y'_{I'}\ ,
\end{tikzcd}
\tag*{(6.3.4)}
\]
d'un autre diagramme carré \textit{dans} $\mathcal{M}$
\[
\begin{tikzcd}
A_0 \arrow[r, "\alpha_0"] \arrow[d, "i_0"'] & X_0 \arrow[d, "p_0"] \\
B_0 \arrow[r, "\beta_0"'] & Y_0\ ,
\end{tikzcd}
\tag*{(6.3.5)}
\]
et enfin d'un homomorphisme $\lambda = (\lambda_A, \lambda_B, \lambda_X, \lambda_Y)$ du carré (6.3.4) dans le carré (6.3.5) :
% foot: Version 9 février 2019
\printedpage{56}\msp{page 112}
% head: A. Grothendieck : Dérivateurs, Ch. XVII, Édition des Universités de Montpellier II et Paris VII
\[
\begin{tikzcd}
A'_{I'} \arrow[rr, "\alpha'_{I'}"] \arrow[dr, "i'_{I'}"'] \arrow[dd, "\lambda_A"'] & & X'_{I'} \arrow[dr, "p'_{I'}"] \arrow[dd, "\lambda_X" near start] & \\
 & B'_{I'} \arrow[ur, dotted, "h'_{I'}"] \arrow[rr, crossing over, "\beta'_{I'}" near end] & & Y'_{I'} \arrow[dd, "\lambda_Y"] \\
A_0 \arrow[rr, "\alpha_0" near start] \arrow[dr, "i_0"'] & & X_0 \arrow[dr, "p_0"] & \\
 & B_0 \arrow[from=2-2, crossing over, "\lambda_B"' near start] \arrow[ur, dotted, "h_0"'] \arrow[rr, "\beta_0"] & & Y_0
\end{tikzcd}
\qquad\qquad \text{\footnotemark[32]}
\tag*{(6.3.6)}
\]
\footnotetext[32]{\textbf{N.B.} La condition $p \in \Psi\langle I\rangle$ s'exprime par $p' \in \Psi\langle I'\rangle$ et \fbox{$p_0 \in \Psi.$}}

D'autre part, la question de prolonger $(I', h') \in \mathcal{H}$ en $(I, h)$, revient à celle de trouver $h_0 : B_0 \longrightarrow X_0$ dans $\mathcal{M}$, de façon qu'avec $h'_{I'}$ cela fasse encore un diagramme commutatif de (6.3.6), ce qui revient à dire que d'une part les deux triangles formés dans le carré de base sont commutatifs :
\[
p_0h_0 \ =\ \beta_0, \qquad\quad h_0i_0 \ =\ \alpha_0\ ,
\tag*{(6.3.7)}
\]
et d'autre part le carré diagonal défini par $h'_{I'}$, $h_0$ soit commutatif (ce qui exprime qu'on a une flèche $h$ de $B$ dans $X$) :
\[
\lambda_X h'_{X'} \ =\ h_0\lambda_B\ . % [sic] h'_{X'} -> h'_{I'}, the arrow is h'_{I'} everywhere else (cf. (6.3.16))
\tag*{(6.3.8)}
\]

\msp{page 113}Mais le diagramme (6.3.6), où la flèche en traits pointillés denses $h'_{I'}$ est déjà acquise, et $h_0$ en pointillés espacés reste à trouver, donne naissance au diagramme, dans lequel est mis en évidence la flèche $\tilde{i} \in \Phi$ qui exprime l'hypothèse $i \in C(i)$, % [sic] C(i) -> C(I), cf. i' \in C(I')
compte tenu de $i' \in C(I')$ :
\[
\begin{tikzcd}
A'_{I'} \arrow[rr, "\alpha'_{I'}"] \arrow[dr, "i'_{I'}"'] \arrow[dd, "\lambda_A"'] & & X'_{I'} \arrow[dr, "p'_{I'}"] \arrow[dd, "\lambda_X" near start] & \\
 & B'_{I'} \arrow[ur, dotted, "h'_{I'}"] \arrow[rr, crossing over, "\beta'_{I'}" near end] & & Y'_{I'} \arrow[dd, "\lambda_Y"] \\
A_0 \arrow[ur, phantom, "\text{cocart.}"] \arrow[rr, "\alpha_0" near start] \arrow[dr, "r"'] \arrow[ddr, dotted, "i_0"'] & & X_0 \arrow[dr, "p_0\in\Psi"] & \\
 & \tilde{B} \arrow[from=2-2, crossing over, "s"' near start] \arrow[ur, dotted, "\tilde{h}_0"] \arrow[rr, dotted, "\tilde{\beta}_0"] \arrow[d, "\tilde{i}\in\Phi"'] & & Y_0 \\
 & B_0\ , \arrow[uur, dotted, "h_0"'] \arrow[urr, "\beta_0"'] & &
\end{tikzcd}
\tag*{(6.3.9)}
\]
% foot: Version 9 février 2019
\printedpage{57}
% head: A. Grothendieck : Dérivateurs, Ch. XVII, Édition des Universités de Montpellier II et Paris VII
où la flèche $\tilde{h}_0$ en traits pointillés denses de $\tilde{B} \overset{\text{déf}}{=} A_0 \amalg_{A'_{I'}} B'_{I'}$ dans $X_0$, est définie par les conditions de commutativité
\[
\tilde{h}_0r \ =\ \alpha_0, \qquad\quad \tilde{h}_0s \ =\ \lambda_X h'_{I'}\ ,
\tag*{(6.3.10)}
\]
ce qui est licite, vu que
\[
\begin{aligned}
\alpha_0\lambda_A \ &=\ (\lambda_X h'_{I'})i'_{I'} \\
&=\ \lambda_X(h'_{I'}i'_{I'}) \\
&=\ \lambda_X\alpha'_{I'}.
\end{aligned}
\]

\msp{page 114}Quant à la flèche $\tilde{\beta}_0 : \tilde{B} \longrightarrow Y_0$, elle est donnée par la condition de commutativité
\[
\tilde{\beta}_0r \ =\ p_0\alpha_0, \qquad\quad \tilde{\beta}_0s \ =\ \lambda_Y\beta'_{I'}\ ,
\tag*{(6.3.11)}
\]
ce qui est encore licite, puisque
\[
(p_0\alpha_0)\lambda_A \ =\ (\lambda_Y\beta'_{I'})i'_{I'}
\]
en vertu de la commutativité du cube (6.3.6). Pour établir la commutativité totale du diagramme (6.3.9) (traits pleins, et traits pointillés denses pour $h'_{I'}$, $\tilde{h}_0$ et $\tilde{\beta}_0$ -- à l'exclusion de $h_0$ pas encore construit), il faut encore vérifier d'une part les relation % [sic] les relation -> les relations, plural
\[
\left\{
\begin{aligned}
\tilde{h}_0r \ &=\ \alpha_0, \quad p_0\tilde{h}_0 \ =\ \tilde{\beta}_0, \\
\lambda_X h'_{I'} \ &=\ \tilde{h}_0s
\end{aligned}
\right.
\qquad\qquad \text{\footnotemark[33]},
\tag*{(6.3.12)}
\]
\footnotetext[33]{\textbf{N.B.} la première et la troisième de ces relations sont déjà connues, cf. (6.3.10).}
établissant la commutativité du cube supérieur avec le carré diagonal (diagramme du type (6.3.6), la flèche pointillée $h_0$ y incluse), d'autre part les relations
\[
\left\{
\begin{aligned}
\tilde{i}r \ &=\ i_0 \qquad \text{(et $\tilde{i}s = \lambda_B$ pour mémoire)} \\
\beta_0\tilde{i} \ &=\ \tilde{\beta}_0\ ,
\end{aligned}
\right.
\tag*{(6.3.13)}
\]
où la première ligne est déjà connue par définition de $\tilde{B}$ et de $\tilde{i}$. La deuxième
\msp{page 115}se ramène aux deux relations
\[
\begin{aligned}
(\beta_0\tilde{i})r \ &=\ \beta_0(\tilde{i}r) \\
&=\ \beta_0i_0 \\
&=\ p_0\alpha_0 \\
&=\ \tilde{\beta}_0r\ ,
\end{aligned}
\]
OK par commutativité du carré de base de (6.3.5),
\[
\begin{aligned}
(\beta_0\tilde{i})s \ &=\ \beta_0(\tilde{i}s) \\
&=\ \beta_0\lambda_B \\
&=\ \lambda_Y\beta'_{I'} \\
&=\ \tilde{\beta}_0s\ ,
\end{aligned}
\]
% foot: Version 9 février 2019
\printedpage{58}
% head: A. Grothendieck : Dérivateurs, Ch. XVII, Édition des Universités de Montpellier II et Paris VII
OK par commutativité du carré antérieur dans (6.3.6).

\dgrblock{dgrb951}{A diamond and the compatibility at its last vertex}
\begin{dcomment}{examples}{p113}{A diamond and the compatibility at its last vertex}
Grothendieck's cube records information that an objectwise cofibration test would miss. A finite example makes the attachment visible. Work in simplicial sets, whose cofibrations are the monomorphisms \citep[Example~11.3.5]{Riehl2014}, and let $J$ be the poset with objects $a,b,c,d$ and relations $a<b<d$ and $a<c<d$. We use covariant diagrams, or equivalently the source's notation $\mathcal M(J^{\mathrm{op}})$.

Define a diagram $B$ by
\[
 B_a=\Delta^0,\qquad
 B_b=B_c=\Delta^1,\qquad
 B_d=\Delta^1\amalg_{\Delta^0}\Delta^1,
\]
where the maps out of $B_a$ choose the initial endpoint of each interval, and the maps into $B_d$ are the two inclusions. The latching objects are
\[
 L_aB=\varnothing,\qquad
 L_bB=L_cB=\Delta^0,\qquad
 L_dB=\Delta^1\amalg_{\Delta^0}\Delta^1.
\]
Every latching map is a simplicial monomorphism: at $d$ it is the identity. Consequently $B$ is projectively cofibrant by the lifting argument just given. Its colimit is $B_d$, since $d$ is terminal in $J$; its realization is two intervals joined at an endpoint. A compatible map out of $B_b$ and $B_c$ is exactly a pair of paths with the same value at their initial endpoints. The pushout in the latching object expresses precisely this requirement.

Now retain the first three values but replace $B_d$ by $\Delta^0$, with both maps from the intervals collapsing them. All four values remain cofibrant simplicial sets. However, the last latching map becomes
\[
 \Delta^1\amalg_{\Delta^0}\Delta^1\longrightarrow\Delta^0,
\]
which is not a monomorphism, since it identifies distinct vertices. The diagram is therefore not projectively cofibrant. A choice of lifts at its four objects need not be compatible along the two paths through the diamond. This is why the proof asks for $\widetilde i\in\Phi$, rather than merely the component map $i_0\in\Phi$.
\end{dcomment}


Donc il reste à vérifier la deuxième relation (6.3.12),
\[
p_0\tilde{h}_0 \ =\ \tilde{\beta}_0\ ,
\]
ce qui se ramène à
\[
\begin{aligned}
(p_0\tilde{h}_0)r \ &=\ p_0(\tilde{h}_0r) \\
&=\ p_0\alpha_0 \\
&=\ \tilde{\beta}_0r\ ,
\end{aligned}
\]
OK par (6.3.11),
\[
\begin{aligned}
(p_0\tilde{h}_0)s \ &=\ p_0(\tilde{h}_0s) \\
&=\ p_0(\lambda_X h'_{I'}) \\
&=\ (p_0\lambda_X)h'_{I'} \\
&=\ (\lambda_Y p'_{I'})h'_{I'} \\
&=\ \lambda_Y(p'_{I'}h'_{I'}) \\
&=\ \lambda_Y\beta'_{I'} \\
&=\ \tilde{\beta}_0s\ ,
\end{aligned}
\]
OK.

Ceci posé, je dis que les contraintes (6.3.7), (6.3.8) sur $h_0$ reviennent à la commutativité des deux triangles dans les carré % [sic] les carré -> le carré, singular (one square (6.3.14))
\[
\begin{tikzcd}
\tilde{B} \arrow[r, "\tilde{h}_0"] \arrow[d, "\tilde{i}\in\Phi"'] & X_0 \arrow[d, "p_0\in\Psi"] \\
B_0 \arrow[r, "\beta_0"'] \arrow[ur, dotted, "h_0"] & Y_0\ ,
\end{tikzcd}
\tag*{(6.3.14)}
\]
\msp{page 116}i.e. aux deux relations
\[
h_0\tilde{i} \ =\ \tilde{h}_0\ , \qquad\quad p_0h_0 \ =\ \beta_0\ .
\tag*{(6.3.15)}
\]

La deuxième de ces relations n'est autre que la première relation (6.3.7), et il reste à voir que la première relation (6.3.15) équivaut à l'ensemble des deux relations restantes dans (6.3.7), (6.3.8) :
\[
h_0i_0 \ =\ \alpha_0\ , \qquad\quad \lambda_X h'_{I'} \ =\ h_0\underbrace{\lambda_B}_{\tilde{i}s}
\tag*{(6.3.16)}
\]
% foot: Version 9 février 2019

% ---- from der-ch17_p061-062.tex
\printedpage{59}
% head: A. Grothendieck : Dérivateurs, Ch. XVII, Édition des Universités de Montpellier II et Paris VII
% foot: Version 9 février 2019

Or la relation en question $h_0\tilde{i} = \tilde{h}_0$ équivaut à l'ensemble des deux relations
\[
\begin{aligned}
(h_0\tilde{i})r &= h_0(\tilde{i}r)\\
&= h_0i_0\\
&\overset{?}{=} \alpha_0\\
&= \tilde{h}_0r,
\end{aligned}
\]
équivalent à $h_0i_0 = \alpha_0$,
\[
\begin{aligned}
(h_0\tilde{i})s &= h_0(\tilde{i}s)\\
&= h_0\lambda_B\\
&\overset{?}{=} \lambda_Xh'_{I'}\\
&= \tilde{h}_0s,
\end{aligned}
\]
équivalent à $h_0\lambda_B = \lambda_Xh'_{I'}$, c'est OK.

Ceci posé, dans le carré (6.3.14), les hypothèses $i \in \Phi(I)$, $p \in \Psi\langle I\rangle$ impliquent $\tilde{i}_0 \in \Phi$, $p_o \in \Psi$. Donc la flèche cherchée $h_0$ existe, q.e.d. % [sic] p_o -> p_0, index printed as letter o, elsewhere index 0 (h_0, i_0)

Cela achève la démonstration du corollaire 6.3, et par là aussi, celle de la dernière partie, d), de 6.1. On est heureux\,!

\msp{page 117}

\dgrblock{dgrb952}{The direct part of Reedy theory}
\begin{dcomment}{later}{p116}{The direct part of Reedy theory}
The proof just completed has become the direct part of Reedy theory. For a category with degrees and unique direct--inverse factorizations, a new vertex has both incoming and outgoing constraints. Relative latching maps encode the former and relative matching maps the latter. In a direct category the inverse part has only identities, so matching objects are terminal and Reedy fibrations are objectwise fibrations. Whenever the projective structure exists, the two model structures then agree because their fibrations and weak equivalences agree.

\citet[Definition~4.14, Observation~4.16 and Theorem~4.18]{RiehlVerity2014Reedy} give the relative maps, the lifting construction and the resulting model structure; their Lemma~7.1 also identifies trivial Reedy cofibrations by trivial relative latching maps. The familiar natural-number degree convention in that treatment is sufficient for finite posets. Grothendieck's well-founded-poset argument allows transfinite rank directly, without requiring a uniform finite bound on the lengths of chains entering a vertex. The equality in 6.3 concerns this direct situation, after the variance has been fixed; it is not an equality of the projective and arbitrary Reedy structures on every diagram category.
\end{dcomment}


\textbf{Proposition 6.4.}\footnote[34]{Sous réserve d'existence des $\varinjlim$ et des $\varprojlim$ pertinentes dans $\mathcal{M}$. On peut supposer pour simplifier $\mathcal{M}$ stable par $\varinjlim$ et $\varprojlim$ .} % mark printed as superscript (34)
\textit{Soit $\mathcal{M}$ une catégorie, et $\Phi \longleftrightarrow \Psi$ un couple Q-clos en Q-dualité ($\Phi = \Psi^*$, $\Psi = \Phi_*$). Soit $u : I \longrightarrow I'$ une flèche dans} Cat\textit{, d'où une suite de foncteurs adjoints}
\[
\begin{tikzcd}[column sep=large]
\mathcal{M}(I) \arrow[r, shift left=3, "u_!"] \arrow[r, shift right=3, "u_*"'] & \mathcal{M}(I') \arrow[l, "u^*"']
\end{tikzcd}\ .
\]
\textit{On définit $\Phi(I)$, $\Psi\langle I\rangle$, $\Phi(I')$, $\Psi\langle I'\rangle$ comme dans} 6.2.5. \textit{On a alors :}

\cercle{a}
\[
\left\{
\begin{aligned}
u_!(\Phi(I)) &\ \subset\ \Phi(I')\ ,\\
u^*(\Psi\langle I'\rangle) &\ \subset\ \Psi\langle I\rangle\ .
\end{aligned}
\right.
\tag*{(6.4.1)}
\]

\cercle{b} \textit{On a l'équivalence}
\[
u^*(\Phi(I')) \subset \Phi(I) \qquad\Longleftrightarrow\qquad u_*(\Psi\langle I\rangle) \subset \Psi\langle I'\rangle\ .
\tag*{(6.4.2)}
\]

\cercle{c} \textit{Les relations équivalentes} (6.4.2) \textit{ont lieu dans chacun des trois cas suivants :}
\begin{enumerate}[label={}, leftmargin=2.5em]
\item 1°) \textit{$u$ est coétale.}
\item 2°) \textit{$I$ est discrète.}
\item 3°) \textit{Pour tout $i' \in \operatorname{Ob} I'$, les composantes connexes de $I/i'$ admettent un objet final.}
\end{enumerate}

\newpage
\printedpage{60}
% head: A. Grothendieck : Dérivateurs, Ch. XVII, Édition des Universités de Montpellier II et Paris VII
% foot: Version 9 février 2019

\textbf{Corollaire 6.5.} \textit{Soit $(W,C,F)$ un triple de Quillen clos spécial dans $\mathcal{M}$, et définissons, pour toute $I$ dans} Cat\textit{, $C(I)$, $TC(I)$ et $F(I)$, $TF(I)$ comme dans le théorème} 6.1. a)\textit{. On a alors}

\msp{page 118}

\textit{les équivalences}
\[
\Bigl(u^*(C(I'))\ \subset\ C(I)\Bigr) \quad\Longleftrightarrow\quad \Bigl(u_*(TF)\ \subset\ TF'\Bigr)\ ,
\tag*{(6.5.1)}
\]
\[
\Bigl(u^*(TC(I'))\ \subset\ TC(I)\Bigr) \quad\Longleftrightarrow\quad \Bigl(u_*(F)\ \subset\ F'\Bigr)\ .
\tag*{(6.5.2)}
\]
\textit{Les conditions équivalentes} (6.5.1) \textit{ainsi que} (6.5.2) \textit{sont satisfaites quand on est dans l'un des cas} 1°) 2°) 3°) \textit{de} 6.4. c).

\dgrblock{dgrb953}{Why terminal comma components preserve cofibrations}
\begin{dcomment}{exposition}{p118:Les conditions équivalentes}{Why terminal comma components preserve cofibrations}
The final proposition asks when restriction can serve on the left of a Quillen adjunction as well as on the right. The distinction is useful in applying the construction to a subdiagram. Restriction always preserves the objectwise right classes. To preserve the left classes as well, its right adjoint must preserve the corresponding objectwise right classes. The commentary on Chapter~XII, manuscript p.~114 (``A terminal fibre controls all the ordinary direct images''), discusses terminal fibres and ordinary direct images; the additional point here is the product stability of right lifting classes.

Here is the formal detail behind the equivalence in 6.4(b). With all required copowers available, one has
\[
 \Psi\langle I\rangle=\Phi(I)^\perp.
\]
One inclusion follows from the definition $\Phi(I)={}^\perp(\Psi\langle I\rangle)$. For the other, use the free evaluation diagrams $F_iA=I(-,i)\cdot A$. If $a\in\Phi$, then $F_ia\in\Phi(I)$ by adjunction. A map right-orthogonal to $\Phi(I)$ therefore has each component right-orthogonal to $\Phi$, hence in $\Psi$. Transposing lifting problems across $u^*\dashv u_*$ now proves both directions of (6.4.2). Factorizations are not needed.

The geometric hypothesis makes right Kan extension particularly simple. Its actual pointwise formula for presheaves is
\[
 (u_*X)(i')=\lim_{(I/i')^{\mathrm{op}}}X.
\]
If each connected component of $I/i'$ has a terminal object $e_\alpha$, that object is initial in the opposite indexing category. A cone on the restricted diagram is then determined by its map into $X(e_\alpha)$. Splitting the index into its components gives
\[
 (u_*X)(i')\cong\prod_\alpha X(e_\alpha).
\]
For a morphism of diagrams the same formula gives a product of its components. Right lifting classes are closed under products: solve a square against the product by projecting to each factor, choosing the factorwise fillers, and using the universal property of the product. The empty product gives an identity map of a terminal object. Thus the displayed formula proves the required preservation.

For completeness, if $u$ is a discrete opfibration and $(i,\gamma:u(i)\to i')$ is an object of $I/i'$, lift $\gamma$ uniquely to $i\to\gamma_!i$. Its target lies over $i'$, and $(i,\gamma)$ has a unique map to $(\gamma_!i,1_{i'})$. A morphism in the comma category leaves this target unchanged by uniqueness of lifting. Consequently the comma category is a disjoint union of components, each with one of these terminal objects. If $I$ is discrete, the comma category itself is discrete and the assertion is immediate. These facts prove the two implications that the source abbreviates as $1^\circ\Rightarrow3^\circ$ and $2^\circ\Rightarrow3^\circ$.
\end{dcomment}


\textsc{Démonstration} de 6.4. Les assertions \cercle{a} , \cercle{b} sont $\pm$ tautologiques, par adjonction, comparer XIII, prop. 2, et cor., page 46. Pour prouver c), il suffit de traiter le cas 3°), car 1° $\Longrightarrow$ 3°, 2° $\Longrightarrow$ 3°. Mais le calcul de
\[
u_*(X)(i')\ \simeq\ \varprojlim_{i\in I/i'} X(i)
\]
montre que si les composantes connexes $(I/i')_\alpha$ de $I/i'$ ont chacune un objet final $e_\alpha$, alors
\[
u_*(X)(i')\ \simeq\ \prod_\alpha X(e_\alpha)\ .
\]
Or $\Psi$, étant clos à droite, est stable par produits quelconques, donc $X \in \Psi\langle I\rangle$ implique que les $u_*(X)(i')$ sont dans $\Psi$, donc $u_*$ est dans $\Psi\langle I'\rangle$, q.e.d.


\dgrblock{dgrb954}{A converse when the source is a poset}
\begin{dcomment}{development}{p118}{A converse when the source is a poset}
Grothendieck's terminal-component condition is sufficient for every Q-dual pair. The elementary model of sets used above shows that, for a poset source, it is also a necessary condition if one wants such preservation even in this one test theory. This identifies exactly how much of the comma category the assertion detects.

\textbf{Proposition}\ours.
Let $u:I\to I'$ be a functor between small categories, with $I$ a poset. Give sets the model structure with all maps weak equivalences, injections cofibrations and surjections fibrations, and give presheaf categories their projective structures. The following conditions are equivalent:
\begin{enumerate}
\item Every connected component of every $I/i'$ has a terminal object.
\item $u^*$ preserves projective cofibrations.
\item $u^*$ takes every representable presheaf on $I'$ to a projective presheaf on $I$.
\end{enumerate}
Here projective means having the lifting property against objectwise surjections, and the axiom of choice is understood.

We first establish the elementary description of projectives needed for the converse. If $P$ is a presheaf on a small poset $I$, form the canonical objectwise surjection
\[
 q:\coprod_{(i,x),\ x\in P(i)} I(-,i)\longrightarrow P.
\]
It is surjective at $j$ because the summand indexed by $(j,z)$ sends $1_j$ to every given $z\in P(j)$. If $P$ is projective, choose a section $s$ of $q$. The induced map of categories of elements takes each connected component of $\int P$ into a single representable summand, say the summand indexed by $(i,x)$. For any element $(j,z)$ in that component, $qs=1$ says that the arrow selected by $s$ has the form $j\to i$ and sends $x$ to $z$. Hence $(j,z)$ has an arrow to $(i,x)$ in $\int P$, so $(i,x)$ belongs to the same component. Because $I$ is a poset there is at most one such arrow. Thus $(i,x)$ is terminal in that component.

Conversely, if every component of $\int P$ has a terminal element $(i_\alpha,x_\alpha)$, the maps represented by these elements give an isomorphism
\[
 \coprod_\alpha I(-,i_\alpha)\xrightarrow{\ \sim\ }P.
\]
At $j$, existence and uniqueness of the arrow to the appropriate terminal element give respectively surjectivity and injectivity. A representable is projective by evaluation, and a coproduct of representables is projective by choosing a lift in each summand. We have proved that projective presheaves on a poset are exactly coproducts of representables, equivalently presheaves whose categories of elements have terminal connected components. The general background is the description of indecomposable projectives as retracts of representables and its Cauchy-complete refinement in \citet[Proposition~0.7 and Theorem~0.8]{vanOosten2021Topos}; the poset case eliminates nonidentity idempotents.

Condition~1 implies condition~2 by 6.4 and 6.5. Condition~2 implies condition~3 because $\varnothing\to I'(-,i')$ is a projective cofibration: any lifting problem against an objectwise surjection is a lifting problem for an element at $i'$. Restriction preserves the initial presheaf. Finally,
\[
 \int\bigl(u^*I'(-,i')\bigr)=I/i',
\]
including both its objects and its morphisms. The description of projectives just proved therefore makes condition~3 imply condition~1, completing the proof.

The poset hypothesis has a precise role. In a general category a projective presheaf may be a proper retract of a representable, controlled by a nontrivial idempotent; a terminal object in its category of elements need not exist. We checked the projective-presheaf results just cited and the projective diagram construction in \citet[\S12.3]{Riehl2014}. We did not locate this exact formulation of the converse to 6.4(c) in those sources. The proposition is an elementary application of those established results, developed here to clarify the scope of Grothendieck's criterion, not a claim of priority.
\end{dcomment}


\dgrblock{dgrb955}{The pointwise formula and the derivator theorem}
\begin{dcomment}{later}{p118}{The pointwise formula and the derivator theorem}
Grothendieck has now constructed enough compatible diagram structures to return to the purpose of the chapter: computing a homotopy Kan extension at one object. The last restriction criterion is the bridge. We spell out this application because the printed chapter stops immediately after the criterion, before drawing the conclusion in these terms.

We keep the convention of \S6, so that $\mathcal M(I)=\operatorname{Fun}(I^{\mathrm{op}},\mathcal M)$. Given $u:I\to I'$ and $i'\in I'$, put $S=i'\backslash I$, with objects $(i,\alpha:i'\to u(i))$, and let $q:S\to I$ and $p:S\to e$ be the projections. An arrow of $S$ over $a:i\to j$ is forced by the equation $\beta=u(a)\alpha$. Thus $q$ is a discrete opfibration. More explicitly, for each $j\in I$, the category $S/j$ splits according to the composite $i'\to u(i)\to u(j)$; the part indexed by $\gamma:i'\to u(j)$ has terminal object $((j,\gamma),\mathrm{id}_j)$. This verifies exactly condition~3 of Proposition~6.4, with no claim that arbitrary restrictions preserve projective cofibrations.

There is also a direct check of the required preservation of cofibrations which uses only coproducts, and therefore does not import the right Kan extension existence proviso in Proposition~6.4. Write $F_i(A)(j)=\coprod_{I(j,i)}A$ for the free presheaf at $i$. Unique outgoing lifts along $q$ give
\[
 q^*F_i(A)\cong\coprod_{s\in q^{-1}(i)}F_s(A).
\]
Restriction preserves colimits and retracts. This formula and the generating descriptions in Theorem~6.1(b) show that it preserves projective cofibrations and trivial cofibrations. Choose a projectively cofibrant replacement $QX\to X$; then $q^*QX$ is projectively cofibrant on $S$. Ordinary left Kan extension gives
\[
 (u_!QX)(i')\cong p_!q^*QX
   =\mathop{\mathrm{colim}}_{S^{\mathrm{op}}}q^*QX.
\]
The first expression computes the value of $\mathbf L u_!X$, and the second computes $\mathbf Lp_!(q^*X)$. Evaluation and restriction preserve all objectwise weak equivalences, so the replacement maps identify this comparison independently of the chosen cofibrant representative. We obtain the canonical pointwise formula
\[
 i'^*\mathbf L u_!X\cong
 \mathop{\mathrm{hocolim}}_{(i'\backslash I)^{\mathrm{op}}}q^*X.
\]
Here the opposite on the comma category is essential. In the covariant convention of \S3 the corresponding indexing category is $I/i'$, without this opposite. The formula is the content of the left Kan extension axiom, not merely an objectwise recipe unrelated to the adjunction.

The other elementary properties can also be read directly from the diagram model structures. Diagrams on a disjoint union form a product of diagram categories; cofibrant and fibrant representatives show that their homotopy categories decompose accordingly. Evaluations jointly detect isomorphisms: represent a homotopy-category morphism between cofibrant--fibrant objects, and use saturation of weak equivalences in a closed model category in each component. If $I$ has a terminal object $t$, the constant-presheaf functor is left adjoint to evaluation at $t$, with identity unit; both functors preserve weak equivalences, so constant diagrams remain fully faithful after localization. These are the precise elementary assertions behind the older axiom names.

The full, two-sided existence theorem is now available without the chapter's special-generation hypotheses, for a model category admitting all small limits and colimits. \citet[Theorem~6.11 and its dual]{Cisinski2003} establishes the homotopy Kan extension and pointwise assertions; his ``left'' terminology there is attached to the right Kan extension convention. For combinatorial model categories, \citet[Theorem~1.28, Proposition~1.30, and Theorem~1.32]{Groth2013} presents the projective--injective argument and records the general theorem. The present chapter provides the projective, colimit side; existence of all ordinary right Kan extensions is not part of its cocompleteness assumption.
\end{dcomment}


\dgrblock{dgrb956}{What the coherent diagram remembers}
\begin{dcomment}{later}{p118}{What the coherent diagram remembers}
Grothendieck's decision to localize the category of whole diagrams has a precise counterpart in the language of higher categories. Suppose, for a direct comparison, that $\mathcal M$ is a combinatorial simplicial model category, and let $\mathcal C_\infty$ be the coherent nerve of its simplicial subcategory of cofibrant--fibrant objects, with universes fixed as usual. Rectification identifies
\[
 \operatorname{Ho}\bigl(\mathcal M^{I^{\mathrm{op}}}\bigr)
 \simeq h\operatorname{Fun}\bigl(N(I^{\mathrm{op}}),\mathcal C_\infty\bigr).
\]
One precise reference is \citet[Proposition~4.2.4.4, Remark~4.2.4.5, and Corollary~4.2.4.8]{Lurie2009HTT}; the statement is made there with size-controlled subcategories, which give the displayed comparison after choosing an ambient universe. Projective and injective presentations have the same weak equivalences. The homotopy Kan extensions above present the corresponding higher categorical Kan extensions.

The comparison is with the homotopy category of a coherent diagram category. It does not identify that category with $\operatorname{Fun}(I^{\mathrm{op}},h\mathcal C_\infty)$. Passing to the homotopy category before forming diagrams discards the compatibility of homotopies. The latching maps in the proof of Corollary~6.3 provide a concrete model-level account of precisely this compatibility: each new lift must extend the lifts already chosen, and its domain is a pushout recording those earlier choices.
\end{dcomment}
\begin{otherlanguage}{english}
\dgroverview{XVII-b}{Chapter XVII \textperiodcentered\ \textup{XVII-b} \textperiodcentered\ Overview}
\selectlanguage{english}
\noindent\textbf{From cofibrant diagrams to a calculus of homotopy Kan extensions}

The second half of Chapter XVII supplies a construction that the preceding discussions have repeatedly needed: a homotopy theory of diagrams in a category of coefficients, compatible with changing the category of indices. Grothendieck begins with the economical axioms of a category with cofibrations and ends with projective model structures on arbitrary small diagram categories. The intervening work explains why the latter are useful. They provide a uniform source of derived adjunctions, while the earlier constructions on ordered sets explain how the resulting cofibrations can be recognized locally. The text therefore moves between two kinds of argument: a formal calculus of lifting and adjunction, and an inductive construction in which the part of a diagram already built controls the next attachment.

This is one concrete realization of the programme described in \citet[\S69, opening discussion]{Grothendieck1983}. A single derived category does not retain enough information about compatible constructions on complexes or spaces. One must also retain the localized categories of diagrams, together with restriction and the operations adjoint to restriction. The letter to Thomason of 2 April 1991 explains the role of the auxiliary model structures in this programme: the weak equivalences specify the homotopy theory; cofibrations and fibrations make its operations constructible \citep[point~1, printed pp.~3--4]{Grothendieck1991Letter}. Chapter XVII shows both the reach of this principle and the difficulty of implementing it with a minimal axiom system.

A convention must remain visible throughout. In \S3 the diagrams under construction are covariant, written $\mathcal M(I^{\circ})=\operatorname{Fun}(I,\mathcal M)$. In \S\S4--6 the notation returns to $\mathcal M(I)=\operatorname{Fun}(I^{\mathrm{op}},\mathcal M)$. A left Kan extension therefore uses $I/i'$ in the first convention and $(i'\backslash I)^{\mathrm{op}}$ in the second. Likewise, well-founded induction concerns an artinian poset in the covariant convention and a noetherian poset in the presheaf convention. The corrections and commentary preserve the printed text where it slips between these formulations, but do not treat a legitimate change of convention as a transcription error.

\medskip\noindent\textbf{The adjunction must survive the construction}

The first obstruction is already present in ordinary localization. An adjunction $L\dashv R$ does not automatically yield an adjunction between a functor obtained by restricting $L$ to a convenient subcategory and the localization of $R$. An equivalence between the localized subcategory and the whole localization is insufficient. Grothendieck's counterexample in Remark~3.6 is decisive: localizing all arrows in a contractible category can identify it with the terminal category, while selecting a terminal object of the original category does not make that object initial. A left adjoint from the terminal category would have to select an initial object. The defect concerns the universal property of the constructed functor, not its existence as a functor.

Theorem~3.7 repairs the argument for the allowed ordered index categories. A functor is factored through its categorical mapping cylinder, or collage. One part is a closed inclusion; its left extension is extension by the initial object. The other part is a retraction with an initial object in every relevant comma category. Each part admits the localized adjunction already justified by the earlier cofibration arguments. Their composite therefore has the desired adjoint property. This is a proof, not an appeal to a general principle that localization preserves adjunctions without hypotheses.

The subsequent attempt to allow arbitrary targets exposes a separate size issue. Finite colimits do not provide the coproduct indexed by every hom-set of an arbitrary target. For example, finite sets admit the finite-branch cofibration axioms with every map a cofibration and only isomorphisms weak equivalences, but restriction from finite $\mathbb Z$-sets to finite sets has no free-action left adjoint. The commentary gives the argument and states precisely which additional colimits repair the construction. Later results for cocomplete cofibration categories require corresponding infinite coproduct and sequential-colimit axioms; they should not be read back into the finite axioms of these pages.

\medskip\noindent\textbf{Small generation and transfer}

Sections~4 and~5 organize a second route. A special cofibration structure has enough colimits, individually accessible objects, a small family generating its cofibrations in the lifting sense, and the inclusion $C^{\square}\subseteq W$. The accessibility condition is used to run the small object argument. It does not by itself assert that the category has a small generating family of presentable objects; identifying it without qualification with local presentability would strengthen the source's hypotheses.

Given $L\dashv R$, the candidate weak equivalences in the target are detected by $R$, while the candidate cofibrations are built from $L(C)$ by pushouts and transfinite compositions. The central calculation is the lifting identity
\[
 (L(C))^{\square}=R^{-1}(C^{\square}).
\]
To make this into a cofibration structure, the source requires $R$ to preserve the relevant colimits and $RL$ to preserve cofibrations. Preservation of trivial cofibrations is a distinct condition. This distinction explains the difficulty of identifying the transferred fibrations: they are initially defined by lifting against the transferred trivial cofibrations, whose description is the very point at issue.

The free-diagram adjunction is particularly favourable. If $I_0$ is the discrete category of objects of $I$, the left adjoint to restriction along $I_0\to I$ sends a family $(A_j)$ to
\[
 i\longmapsto\coprod_{j\in I} I(i,j)\cdot A_j.
\]
For a map of families, every component of its free-diagram image is a coproduct of coefficient maps. Pushout stability and transfinite closure give the needed coproduct stability. The commentary supplies both the cardinal argument for accessibility of an entire diagram and the elementary transfinite construction of these coproducts.

Section~5 then isolates the remaining model-category problem. Put $TC=C\cap W$, $TF=C^{\square}$ and $F=(TC)^{\square}$. The established factorization gives $W=TF\circ TC$; a retract argument gives $TF=F\cap W$. To obtain a Quillen triple one needs the second factorization, into a trivial cofibration followed by a fibration. A small generating family for $TC$ supplies it. With this second factorization, a pushout and retract argument also proves retract stability of $W$, justifying the later comparison with closed model categories. Under the corresponding preservation hypotheses on $RL$, Proposition~5.8 identifies transferred trivial cofibrations as the lifting closure of $L(TC)$ and transferred fibrations as $R^{-1}(F)$.

The distinction between a generating image, its cellular closure, its retract closure, and its double-orthogonal closure matters here. The equality suggested in footnote~23 cannot be true for the raw image of the free-diagram functor; a worked equivariant example shows why pushouts are indispensable. Proposition~5.1 also contains a displayed equality stronger than the argument preceding the additional hypotheses establishes. The audit records the precise inclusions proved and separates this unresolved strengthening from the later, correctly proved transfer theorem. We do not infer falsity merely from an incomplete justification.

\medskip\noindent\textbf{The local structure of a projective cofibration}

Theorem~6.1 constructs, for every small $I$, the model structure on $\mathcal M(I)$ whose weak equivalences and fibrations are objectwise. Thus every left Kan extension is a left Quillen functor and its derived functor is left adjoint to restriction. Its cofibrations are generally more restrictive than objectwise cofibrations. The rest of the chapter explains what their compatibility condition says.

For a map of arrows, the decisive map is the relative latching map
\[
 X_0\amalg_{X_1}Y_1\longrightarrow Y_0.
\]
Along with the cofibration at the source of the arrow, it characterizes the projective lifting condition. Lemma~6.2.5 states this criterion but leaves its proof unwritten. The commentary proves both directions under exactly the stated pushout hypothesis. No terminal object or model-category factorization is needed for that lemma.

For a well-founded diagram, the same pushout occurs every time one extends an already constructed lift to a new object. Grothendieck uses Zorn's lemma to organize the extensions and writes out the final compatibility cube in detail. The mathematical content of that cube is that one lift from a relative latching object enforces both the pointwise lifting equation and naturality with all previous components. Corollary~6.3 correctly places the noetherian hypothesis on $I$ in the presheaf convention. The contrary word in Theorem~6.1(d) is a genuine printed variance slip: the inverse sequence of nonempty tails of $\mathbb N$, with empty inverse limit, gives a counterexample to the literal artinian version. Chapter XIX, Proposition~3.2 and \S3.4, subsequently uses the corrected noetherian formulation \citep{Grothendieck1990XIX}.

The final proposition identifies restrictions that preserve projective cofibrations. When the connected components of each $I/i'$ have terminal objects, right Kan extension is a product of evaluations. Right lifting classes are closed under products, and adjunction gives the preservation assertion. The commentary also develops a converse test with set-valued presheaves on a poset, making precise how much the terminal-component condition can detect.

\medskip\noindent\textbf{What is settled, and what the commentary adds}

The arbitrary-shape derivator construction is no longer an open problem for model categories with small limits and colimits. Cisinski proves the required homotopy Kan extension theorem without the special-generation restriction \citep[Theorem~6.11 and its dual]{Cisinski2003}; Groth gives a projective--injective treatment in the combinatorial case \citep[Proposition~1.30 and Theorem~1.32]{Groth2013}. For a combinatorial simplicial presentation, rectification identifies the localized diagram categories with $h\operatorname{Fun}(N(I^{\mathrm{op}}),\mathcal C_\infty)$ \citep[Proposition~4.2.4.4]{Lurie2009HTT}. These results concern coherent diagrams, retaining the compatibility which the latching construction makes explicit.

The developments marked \ours\ are proved in the accompanying blocks and make no claim of priority. In particular, a special closed model structure on abelian groups detected by $\operatorname{Hom}(\mathbb Z[1/2],-)$ shows that weak equivalences need not be closed under all filtered colimits. The counterexample addresses the expectation expressed on manuscript p.~97 while respecting the chapter's generation and accessibility hypotheses. Together with the poset restriction criterion and the complete missing arrow-category proof, it turns the commentary into a mathematical continuation of the arguments at the points where a reader needs one.
\end{otherlanguage}
\dgrchunkend
